\documentclass[11pt]{article}
\usepackage[letterpaper,margin=1in]{geometry}
\usepackage[T1]{fontenc}
\usepackage{lmodern,microtype,mathrsfs}
\usepackage{amsmath,amssymb,amsthm,mathtools}
\usepackage{enumitem,needspace}
\usepackage{xcolor,fancyhdr}
\usepackage[backend=biber,style=numeric,sorting=none,giveninits=true,
 maxbibnames=8,isbn=false,url=false,doi=true,eprint=true]{biblatex}
\addbibresource{references.bib}
\renewcommand*{\bibfont}{\small}
\AtEveryBibitem{%
 \iffieldequalstr{entrykey}{Fine2011}{\clearfield{note}}{}%
 \iffieldequalstr{entrykey}{BahuaudGuentherIsenbergMazzeo2025Bach}{\clearfield{note}}{}%
 \iffieldequalstr{entrykey}{Baker2026ChiralMetricEvolution}{\def\currentnote{Working manuscript}%
   \restorefield{note}{\currentnote}}{}%
}
\usepackage[colorlinks=true,linkcolor=blue!40!black,citecolor=blue!40!black,
  urlcolor=blue!40!black]{hyperref}
\hypersetup{pdftitle={Twistor sections and chiral Weyl curvature on homotopy four-spheres},
 pdfauthor={Edward B. Baker III},bookmarksdepth=2}
\pagestyle{fancy}
\fancyhf{}
\fancyhead[L]{\small Twistor sections and chiral Weyl curvature}
\fancyhead[R]{\small Edward B. Baker III}
\fancyfoot[C]{\thepage}
\setlength{\headheight}{14pt}
\setlength{\emergencystretch}{2em}
\setcounter{tocdepth}{1}
\allowdisplaybreaks[1]
\numberwithin{equation}{section}
\newtheorem{theorem}{Theorem}[section]
\newtheorem{proposition}[theorem]{Proposition}
\newtheorem{lemma}[theorem]{Lemma}
\newtheorem{corollary}[theorem]{Corollary}
\theoremstyle{definition}
\newtheorem{definition}[theorem]{Definition}
\theoremstyle{remark}
\newtheorem{remark}[theorem]{Remark}
\DeclareMathOperator{\Ric}{Ric}
\DeclareMathOperator{\Vol}{Vol}
\DeclareMathOperator{\tr}{tr}
\DeclareMathOperator{\Scal}{Scal}
\DeclareMathOperator{\Rm}{Rm}
\newcommand{\rd}{\mathrm{rd}}
\newcommand{\HS}{\mathrm{HS}}

\title{Twistor sections and chiral Weyl curvature\\on homotopy four-spheres}
\author{Edward B. Baker III}
\date{Working draft, 24 November 2026}
% Prepared for Edward B. Baker III with OpenAI Codex assistance.
% The author-selected lead was retained. Its exact model identifier and effort
% were unavailable. Bounded analytical work and editorial assistance were
% configured as gpt-6.2-sol/ultra; backend settings were not independently attested.
% Manuscript revision based on 3fb565b70247a57e4fcc7e65886c57ce0ca0e1ab.
% Analytical results through SEC-E58. Input digests in the manuscript README.
\begin{document}
\maketitle
\thispagestyle{plain}
\begin{abstract}
The self-dual Weyl tensor \(W^+\) obstructs integrability of the positive
twistor space of an oriented Riemannian four-manifold. For a fixed metric,
we bound the optimized differentiated Frobenius defect of local twistor
sections between \(2/9\) and \(4\) times the chiral Weyl energy \(E_+\).
On a smooth homotopy four-sphere with a round marked ball, an arbitrarily
small exterior change of metric can preserve \(E_+\) exactly and make the
full linearization of \(W^+\) surjective. An exact local radial contraction
then starts at the new metric. A sign condition on the adjoint kernel
gives a separate contraction through the original metric. At a round
spherical interface we determine the opposite-chirality response and
the two moments that govern compactly supported radial variations.

On the standard sphere we construct regular marked metrics approaching
the sharp, unattained Yamabe-zero infimum \(E_+=8\pi^2\) for which the
Bach first variation forces the strain of unit radial contraction to
diverge. Below this energy, with a fixed
margin and positive class Yamabe constant, we consider sequences with
vanishing total Bach norm. A supplied complete nonflat asymptotically
locally Euclidean blow-up with smoothly embedded full core must have
one Euclidean end. It carries more than \(4\pi^2\) in each chirality
and leaves a uniformly positive energy remainder. For exactly
Bach-flat positive-Yamabe metrics across smooth homotopy four-spheres,
the set of positive critical values below \(8\pi^2\), if nonempty, has
an attained minimum strictly above \(4\pi^2\). The remaining global
problem is to construct and control a deformation or a convergent
sequence on the given smooth manifold.
\end{abstract}

\clearpage
\tableofcontents

\section{Introduction}\label{sec:introduction}

In two real dimensions an oriented conformal metric determines rotation
by a right angle in every tangent plane. The corresponding complex line
distribution is involutive, and its local holomorphic coordinates are
isothermal coordinates for the metric. In four real dimensions the
choice of an orthogonal complex structure of positive orientation at a
point is a two-sphere. These spheres form the positive twistor bundle.
Atiyah, Hitchin and Singer showed that the natural almost complex
structure on its total space is integrable precisely when one chiral
block of the Weyl tensor vanishes~\cite[Theorem 4.1]{AHS1978}.
We choose the orientation convention in which this block is \(W^+\)
and call \(W^+=0\) anti-self-duality (ASD). The resulting question about
coordinates has two parts. How much of \(W^+\) can be recovered from
the failure of locally chosen complex structures to be integrable?
Can the underlying conformal metric itself be changed so that this
failure disappears?

The smooth four-dimensional Poincar\'e conjecture gives a setting in
which this distinction matters. Let \(M\) be a closed smooth manifold
homotopy equivalent to \(S^4\). Its topological type is known by
Freedman's theorem~\cite{Freedman1982}. To recognize its smooth
structure by twistor geometry it suffices to produce one smooth
Riemannian ASD metric on the given \(M\). Simple connectivity and
the signature formula then give a global conformal developing map to
the round sphere. A coordinate ball in \(M\) gives a place to begin
this construction. The smooth structure and the curvature of its
complement remain part of the problem.

\subsection{Normalization and smooth recognition}

Let \(g\) be an oriented positive-definite metric on a real
four-manifold. Its Hodge operator splits
\(\Lambda^2=\Lambda_g^+\oplus\Lambda_g^-\), with each summand of real
rank three. We regard \(W^\pm\) as symmetric trace-free endomorphisms
of these summands. Set
\begin{equation}\label{eq:energies}
 E_\pm(g)=\int_M |W^\pm_g|_{\mathrm{HS}}^2\,dV_g,\qquad
 Y([g])=\inf_{u\in C^\infty(M),\ u>0}
 \frac{\int_M(6|du|_g^2+s_g u^2)\,dV_g}
      {(\int_M u^4\,dV_g)^{1/2}},
\end{equation}
where \(s_g\) is scalar curvature. The first definition uses the
Hilbert--Schmidt norm of the curvature operator. In the full
four-index tensor norm each unordered pair appears twice in each
of the two skew slots. Consequently
\(|W|_{\mathrm{tens}}^2=4(|W^+|_{\mathrm{HS}}^2+
|W^-|_{\mathrm{HS}}^2)\). Under \(g\mapsto u^2g\), the splitting
of two-forms is fixed, the Weyl endomorphism has weight \(-2\), and
the volume form has weight \(4\). Both energies and the Yamabe
constant in \eqref{eq:energies} depend only on the conformal class.

The signature and Gauss--Bonnet formulas in these conventions
are~\cite{Hirzebruch1971Signature,ChangGurskyYang2003}
\begin{align}
12\pi^2\tau(M)&=E_+(g)-E_-(g),\label{eq:sign}\\
8\pi^2\chi(M)&=E_+(g)+E_-(g)
 -\tfrac12\int_M|\Ric_0|^2\,dV_g
 +\tfrac1{24}\int_M s_g^2\,dV_g .\label{eq:GB}
\end{align}
Here \(\Ric_0=\Ric-\tfrac14 s_g g\). On a homotopy four-sphere
\(\pi_1(M)=0\), \(H^2(M;\mathbb Z)=0\), \(\chi=2\), and \(\tau=0\).

\begin{proposition}[Two sufficient conditions for smooth recognition]\label{prop:recognition}
Let \(M\) be an oriented smooth homotopy four-sphere.
If \(E_+(g)=0\), then \((M,[g])\) is conformally diffeomorphic to the
round sphere. If \(Y([g])>0\) and \(E_+(g)<4\pi^2\), then \(M\)
is diffeomorphic to the standard \(S^4\).
\end{proposition}
\begin{proof}
In the first case \eqref{eq:sign} gives \(E_-=0\). Smoothness and
nonnegativity give \(W=0\) everywhere. The locally conformally flat
charts have M\"obius transition maps. Development along paths is
single-valued by simple connectivity~\cite{SchoenYau1988Conformal}.
The developing map is a local diffeomorphism to \(S^4\). Its image is
open and compact, hence the whole sphere. The map is proper and is a
covering, so simple connectivity of the target gives a diffeomorphism.
In the second case the full four-index integral is \(8E_+<32\pi^2
=16\pi^2\chi(M)\). With the positive Yamabe hypothesis these are the
assumptions of Chang, Gursky and Yang~\cite[Theorem A and Remark 2]{ChangGurskyYang2003}.
Their two possibilities for the smooth type are \(S^4\) and
\(\mathbb{RP}^4\). Simple connectivity selects \(S^4\).
\end{proof}

\subsection{Local data, global metric}

A positive orthogonal almost complex structure \(J\) defines a two-form
\(\omega_J(X,Y)=g(JX,Y)\) of norm \(\sqrt2\). Its unit axis is
\(\omega_J/\sqrt2\in\Lambda_g^+\). Where spin frames are used, the
same sphere is the complex projective line \(\mathbb P(S^+_x)\).
The two descriptions agree under the covering
\(\mathrm{SU}(2)\to\mathrm{SO}(3)\).
In particular a projective spinor representative is free up to
nonzero complex scale. On overlaps the Hodge and spin frames transform
by their bundle transition maps. All local sections used below belong to
the bundle of a \emph{single} metric. A global positively oriented almost complex structure
on a homotopy four-sphere would give
\(c_1=0\) and \(\langle c_2,[M]\rangle=2\), while
\(p_1=c_1^2-2c_2\) and the signature theorem give
\(\langle p_1,[M]\rangle=0\). Local sections therefore carry exactly
the kind of data one can choose without a topological existence
assumption.

Section~\ref{sec:sections} turns the twistor integrability condition
into a measurement of curvature. Section~\ref{sec:lift} then formulates
the prescription of \(W^+\) in a metric chart of fixed volume density.
The linearization has a finite-dimensional adjoint kernel. Its moments
are the obstruction to prescribing the entire target, including its
values on the exterior of the marked ball.

Two local theorems address these moments. One may perturb the
reference metric arbitrarily little in the exterior at exactly the
original energy. Small curvature targets can then be prescribed
through this new reference. If the original metric must be kept,
paired kernel variations supply a quadratic correction under a sign
condition on the original adjoint fields. These constructions appear
in Sections~\ref{msreg:section} and~\ref{mscone:section}.
Sections~\ref{sec:transmission} and~\ref{sec:collars} explain which
constraints return when a variation must be supported in a prescribed
region. A radial source inside a round ball fixes five angular
projections of the opposite linearized curvature in a surrounding
annulus. Fixing both ends of a collar imposes further compatibility
conditions.

To compare local prescription with global recognition, consider the
chiral energy along a radial contraction. A deformation realizing
\(W^+(t)=(1-t)W^+(0)\) in the fixed-density chart has
\(E_+(t)=(1-t)^2E_+(0)\). Its initial metric strain is bounded below
by \(4E_+/\int |B|\,dV\), where \(B\) is the Bach tensor. The doubled
Taub--NUT construction on the standard sphere makes this ratio
diverge even at regular reference metrics. It identifies \(8\pi^2\)
as the sharp Yamabe-zero infimum for classes with a round marked ball
(Sections~\ref{sec:variation}--\ref{mscost:section}).

This leads to the positive-Yamabe band
\(4\pi^2\le E_+\le8\pi^2-\delta\), with \(\delta>0\).
The lower endpoint is the boundary of the strict recognition criterion
in Proposition~\ref{prop:recognition}. In this band the closed
Gauss--Bonnet identity controls the scalar and trace-free Ricci terms
in Yamabe gauge. Section~\ref{msale:section} compares those bounds to
the Euler and signature contributions of a smoothly embedded blow-up.
Section~\ref{msgap:section} uses the resulting compactification to prove
a gap in the closed critical energies. Section~\ref{msgauge:section}
treats convergence on the original smooth manifold and the converse
construction of shrinking marks from a supplied critical class.

\section{The Weyl energy measured by twistor sections}
\label{sec:sections}

Fix a smooth closed oriented Riemannian four-manifold \((M,g)\).
Write \(n_J=g(J\,\cdot,\cdot)/\sqrt 2\) for the unit self-dual axis of
a positive orthogonal almost complex structure \(J\). For each smooth local section
put
\[
 P=\tfrac12(1+\mathrm iJ),\qquad Q=1-P,\qquad
 \mathcal D=P T_{\mathbb C}M=T^{0,1}_J,\qquad F_J(X,Y)=Q[PX,PY].
\]
We take Hermitian norms on complexified tensors. In the skew input pair
of \(F_J\) and \(\nabla F_J\), we use exterior norms, counting each
unordered pair once. The real norm of \(\nabla J\) is the usual tensor
norm.

Choose a finite open cover with a smooth nonnegative subordinate
partition of unity \((\rho_a)\), and on each \(U_a\) choose three smooth
sections \(J_{ai}\) whose axes are orthonormal. All sections are for the
same metric \(g\). Define
\begin{equation}
 \label{secjet:energy}
 I(g)=\inf\sum_{a,i}\int_{U_a}\rho_a
   \bigl(|\nabla F_{J_{ai}}|^2+|\nabla J_{ai}|^2|F_{J_{ai}}|^2\bigr)
 \,dV_g ,
\end{equation}
where the cover, partition, and sections may all vary. The sections on
different sets are chosen independently. Let \(I_0(g)\) be the
corresponding infimum with density \(|F_{J_{ai}}|^4\).

\begin{theorem}[Energy measured by three local sections]
 \label{secjet:comparison}
 Under these conventions,
 \[
      \tfrac29 E_+(g)\le I(g)\le 4E_+(g),\qquad I_0(g)=0 .
 \]
\end{theorem}

\subsection{The null curvature component and the lower bound}

The bracket in \(F_J\) is tensorial after projection because the
terms differentiating scalar multipliers lie in \(\mathcal D\). With
the Levi--Civita connection \(\nabla\), torsion freeness gives
\begin{equation}
 \label{secjet:bracket}
 F_J(X,Y)=Q\bigl((\nabla_{PX}P)Y-(\nabla_{PY}P)X\bigr).
\end{equation}
Its inputs are projected by \(P\wedge P\), its value lies in
\(\overline{\mathcal D}\), and it vanishes precisely when
\(\mathcal D\) is involutive. Denote the complex-bilinear extension of
the metric by \(g_{\mathbb C}\). Both \(\mathcal D\) and
\(\overline{\mathcal D}\) are totally null for \(g_{\mathbb C}\), and
they pair nondegenerately with one another.

\begin{lemma}
 \label{secjet:null}
 Let \(X,Y\) be a Hermitian-orthonormal basis of \(\mathcal D\) at a
 point. Lower the indices of \(X\wedge Y\) to obtain the Hermitian-unit
 self-dual null form \(\xi\). Then
 \begin{equation}
  \label{secjet:nullestimate}
  |\langle W_g^+\xi,\xi\rangle_{\mathbb C}|^2
       \le \tfrac94\bigl(|\nabla F_J|^2+|\nabla J|^2|F_J|^2\bigr),
 \end{equation}
 where the Weyl pairing is complex bilinear.
\end{lemma}

\begin{proof}
 Write \(F=F_J\) and \(B(X,Y)=Q\nabla_XY\) for \(X,Y\) in
 \(\mathcal D\). The metric-derivative terms in Koszul vanish on the
 null distribution, giving
 \begin{equation}
  \label{secjet:koszul}
  2g_{\mathbb C}(B(X,Y),Z)=g_{\mathbb C}(F(X,Y),Z)
    -g_{\mathbb C}(F(Y,Z),X)+g_{\mathbb C}(F(Z,X),Y).
 \end{equation}
 Use the metric connection
 \(\nabla^0_X s=P\nabla_X(Ps)+Q\nabla_X(Qs)\), which preserves the
 splitting. The difference
 \(K_X=\nabla_X-\nabla^0_X=(\nabla_XP)(P-Q)\) interchanges the
 summands, and \(|K|=\tfrac12|\nabla J|\). Its contribution to
 the torsion of \(\nabla^0\) on \(\mathcal D\times\mathcal D\) is
 \(T^0(X,Y)=P[X,Y]-[X,Y]=-F(X,Y)\). In the convention
 \(R(X,Y)=[\nabla_X,\nabla_Y]-\nabla_{[X,Y]}\), the curvature
 comparison reads
 \[
 R(X,Y)=R^0(X,Y)+[K_X,K_Y]
        +(\nabla_X^0K)_Y-(\nabla_Y^0K)_X+K_{T^0(X,Y)} .
 \]
 The last term comes from the difference between the alternating
 tensor derivative and the exterior covariant derivative of \(K\).
 Both \(R^0\) and the commutator preserve \(\mathcal D\), so they
 vanish in the null pairing \(g_{\mathbb C}(R(X,Y)Y,X)\).
 On \(\mathcal D\), \(K_YZ=B(Y,Z)\). Moreover, for \(U,V\in
 \mathcal D\) and any \(L\),
 \[
    (\nabla^0_LF)(U,V)=Q(\nabla_LF)(U,V).
 \]
 Indeed the input correction terms contain a \(K_LU\) or \(K_LV\)
 in \(\overline{\mathcal D}\), and the output correction
 \(K_L F(U,V)\) lies in \(\mathcal D\).

 In a positive orthonormal frame with \(e_2=Je_1\) and \(e_4=Je_3\),
 a unit representative of \(\Lambda^2\mathcal D\) is
 \((e^{13}-e^{24}+\mathrm i(e^{14}+e^{23}))/2\), which is self-dual.
 The Schouten part of the curvature vanishes on the four null
 arguments. Hence the null curvature pairing is the stated Weyl
 evaluation, with the curvature operator normalized by
 \(\langle\mathcal R(X\wedge Y),U\wedge V\rangle_{\mathbb C}
 =g_{\mathbb C}(R(X,Y)V,U)\). Differentiating
 \eqref{secjet:koszul} with \(\nabla^0\), which preserves the null
 pairing, and substituting the projected derivative gives
 \begin{align*}
 \langle W_g^+\xi,\xi\rangle_{\mathbb C}
  &=g_{\mathbb C}((\nabla_XF)(X,Y),Y)
    -g_{\mathbb C}((\nabla_YF)(X,Y),X)\\
  &\quad-g_{\mathbb C}(K_{F(X,Y)}Y,X).
 \end{align*}
 The derivative terms have orthogonal Hermitian unit inputs
 \(X\otimes(X\wedge Y)\) and \(Y\otimes(X\wedge Y)\). Their
 combined magnitude is bounded by \(\sqrt2|\nabla F|\). The torsion
 term is bounded by \(\tfrac12|\nabla J||F|\). Applying
 Cauchy--Schwarz to these two bounds proves \eqref{secjet:nullestimate}.
 This null curvature identity is the intrinsic version of the
 derivative and connection terms in the Newman--Penrose calculation
 \cite[eq.~(13) and Appendix eq.~(74)]{GoverHillNurowski2011GoldbergSachs}.
\end{proof}

For an orthonormal real basis \(n_1,n_2,n_3\) of \(\Lambda^+\),
use the cyclic null representatives
\(\xi_i=(n_j+\mathrm i n_k)/\sqrt2\) for the three axes.
Conjugation or a unit phase leaves the magnitudes
\(q_i=|\langle W^+\xi_i,\xi_i\rangle_{\mathbb C}|\) unchanged.
Write the diagonal entries of \(W^+\) as \(d_i\) and its off-diagonal
entries as \(w_{ij}\). Direct expansion gives
\begin{equation}
 \label{secjet:orthogonal}
 \sum_{i=1}^3 q_i^2=\tfrac34\sum_i d_i^2+\sum_{i<j}w_{ij}^2
       \ge \tfrac12 |W^+|_{\mathrm{HS}}^2 .
\end{equation}
Here the differences of the three diagonal entries contribute
\(\tfrac34\sum_i d_i^2\) because \(\sum_i d_i=0\).
Applying \eqref{secjet:nullestimate} to the triple proves the pointwise
lower bound for its density in \eqref{secjet:energy}. Multiplication by
the partition weights and summation give
\(E_+(g)\le \tfrac92 I(g)\).

\subsection{Conformal normal jets and free refinement}

Fix \(p\in M\). Let \(\mathsf P_g=\tfrac12(\mathrm{Ric}_g-\tfrac16s_g g)\)
be the Schouten tensor. Choose \(f\) with \(f(p)=0\), \(df(p)=0\)
and \(\nabla_g^2 f(p)=\mathsf P_g(p)\), for instance a quadratic
in \(g\)-normal coordinates. For \(h=e^{2f}g\), the Schouten
transformation law
\[
 \mathsf P_h=\mathsf P_g-\nabla_g^2 f+df\otimes df-\tfrac12|df|_g^2g
\]
gives \(\mathsf P_h(p)=0\) \cite[(6.7) and proof of Proposition 8.4]{FeffermanGraham2008Ambient}.
Thus the curvature operator of \(h\) at \(p\) is \(W_h^+\oplus W_h^-\).
The Hodge splitting is conformally invariant, and \(W_h^+=W_g^+\)
at \(p\) as an operator since \(f(p)=0\).

Radially parallel transport an orthonormal basis \((\omega_i)\) of
\(\Lambda^+_h\) along the \(h\)-normal rays. This gives a smooth frame
near \(p\). Define \(h(J_i\,\cdot,\cdot)=\sqrt2\,\omega_i\). The
\(J_i\) are also orthogonal for \(g\), with axes \(e^{-2f}\omega_i\)
still orthonormal. Let \(A_b\) be the connection matrices in the radial
frame, and let \(\Omega_{ab}\) be their curvature matrix components.
The radial identity \(x^b A_b=0\) implies
\(A_b(p)=0\) and
\(\partial_a A_b(p)+\partial_b A_a(p)=0\).
Consequently \(\partial_a A_b(p)=\tfrac12\Omega_{ab}(p)\).
The \(J_i\) are mutually orthogonal for the endomorphism pairing
and each has squared norm \(4\). Therefore
\begin{equation}
 \label{secjet:radial}
 \nabla^h J_i(p)=0,\qquad
 \sum_i |(\nabla^h)^2J_i|^2(p)
       =2\sum_{a<b}|\Omega_{ab}|_{\mathrm{HS}}^2(p)
       =8|W_g^+|_{\mathrm{HS}}^2(p).
\end{equation}
For the last equality, identify a two-form \(\alpha\) with the skew
endomorphism \(K_\alpha\) given by \(h(K_\alpha X,Y)=\alpha(X,Y)\).
The curvature on \(\Lambda^+\) acts by the commutator with
\(K_\alpha\), where \(\alpha\) is the curvature row in the
\(e_a\wedge e_b\) direction, up to its overall sign. The
anti-self-dual ideal acts trivially. A unit self-dual generator
corresponds to left multiplication by a unit imaginary quaternion
divided by \(\sqrt2\). Its commutator acts on the other two unit
generators with coefficients of magnitude \(\sqrt2\). Thus the
action on \(\Lambda^+\) has squared norm \(4|\pi_+\alpha|^2\),
where \(\pi_+\) is the self-dual projection. Summing in \(a<b\)
gives \(4|W_h^+|^2\) since the positive curvature row is
\(W_h^+\pi_+\) at \(p\)
\cite[\S1, pp.~427--429]{AHS1978}.

Since \(\nabla^hP_i(p)=0\), \eqref{secjet:bracket} gives \(F_{J_i}(p)=0\).
Suppressing the index \(i\), differentiation at the center gives
\[
 (\nabla^h_ZF)(X,Y)
   =Q((\nabla^h)^2P)(Z,PX)Y-Q((\nabla^h)^2P)(Z,PY)X .
\]
With ordered skew input slots, each term has norm at most
\(\tfrac12|(\nabla^h)^2J|\). Passing to the exterior-pair
norm divides the squared norm of the alternating tensor by two.
The metrics give the same \(F\) because its bracket definition uses
only \(J\), and they have the same connections and norms at \(p\)
because \(f(p)=df(p)=0\). By \eqref{secjet:radial}, the density for
these sections at \(p\), in the fixed metric \(g\), satisfies
\begin{equation}
 \label{secjet:center}
 \sum_i\bigl(|\nabla F_{J_i}|^2+|\nabla J_i|^2|F_{J_i}|^2\bigr)(p)
      \le 4|W_g^+|_{\mathrm{HS}}^2(p).
\end{equation}

For any \(\epsilon>0\), choose the sections at each center as
above and shrink its open neighborhood until the density in
\eqref{secjet:center} is at most
\(4|W_g^+|_{\mathrm{HS}}^2+\epsilon/\operatorname{Vol}_g(M)\) throughout the
neighborhood. At the same time arrange
\(\sum_i|F_{J_i}|^4\le\epsilon/\operatorname{Vol}_g(M)\) by the
vanishing at the center. Use a finite subcover and a subordinate
partition of unity. Both pointwise inequalities persist under
the weighted sum and integrate to the claimed upper bounds within
\(\epsilon\). Derivatives are taken on each patch before weighting.
Taking the infima and letting \(\epsilon\) tend to zero completes
the proof of Theorem~\ref{secjet:comparison}.

In particular \(I(g)=0\) is equivalent to \(W_g^+=0\). For a smooth
homotopy four-sphere, if the same conformal class has positive Yamabe
constant then \(I(g)<8\pi^2/9\) gives the strict conformal sphere
criterion. Indeed it yields \(E_+(g)<4\pi^2\), and the vanishing
signature gives \(\int_M|W_g|_{(0,4)}^2\,dV_g=8E_+(g)\)
\cite{Hirzebruch1971Signature,ChangGurskyYang2003}. The relaxation in
\eqref{secjet:energy} concerns the auxiliary sections of a supplied
metric. Its value can be approached by refining the cover without
any assertion of attainment.

\subsection{Separated axes and the sharp loss}

For a real Euclidean three-space, write \(d([n],[m])=|n\times m|\)
for the sine distance between unoriented unit axes. For any three
axes with minimum pairwise distance \(\delta>0\) and any real
symmetric trace-free \(W\), the Hermitian-unit null evaluations satisfy
\begin{equation}
 \label{secjet:separation}
 |W|_{\mathrm{HS}}^2
       \le \tfrac{128}{3}\delta^{-4}
               \sum_i|\langle W\xi_i,\xi_i\rangle_{\mathbb C}|^2 .
\end{equation}
To prove this, diagonalize \(W\ne0\) with eigenvalues
\(\lambda_1\ge\lambda_2\ge\lambda_3\). Put \(D=\lambda_1-\lambda_3\),
\(c=-3\lambda_2/D\), so \(|c|\le1\). The two real unit vectors
\[
 a=\sqrt{\tfrac{1+c}{2}}e_1+\sqrt{\tfrac{1-c}{2}}e_3,\qquad
 b=\sqrt{\tfrac{1+c}{2}}e_1-\sqrt{\tfrac{1-c}{2}}e_3
\]
satisfy
\[
 W=D\bigl(\tfrac12(a\otimes b+b\otimes a)-\tfrac c3 I\bigr),
 \qquad |W|_{\mathrm{HS}}^2\le\tfrac23D^2 .
\]
For the null line of the axis \(n\), take an orthonormal basis
\(u,v\) of \(n^\perp\) and \(\xi=(u+\mathrm i v)/\sqrt2\).
Nullity removes the scalar term and gives
\[
 |\langle W\xi,\xi\rangle_{\mathbb C}|
        =\tfrac D2 d([a],[n])\,d([b],[n]).
\]
Equivalently the Weyl evaluation is a quartic on the null conic because
its representative \((z^2-w^2,\mathrm i(z^2+w^2),2zw)\) is quadratic
in the projective parameter \([z:w]\). For a real \(W\) the
roots include both null lines of each zero axis
\cite[Prop.~5.3]{GoverHillNurowski2011GoldbergSachs}.
The function \(d\) is a distance because it equals
\(|n\otimes n-m\otimes m|_{\mathrm{HS}}/\sqrt2\).
Each open ball of radius \(\delta/2\) about either factor axis
contains at most one of the three sample axes. At least one
sample therefore has evaluation of magnitude at least
\(D\delta^2/8\). This proves \eqref{secjet:separation}.

For sharpness take \(W=e_3\otimes e_3-I/3\) and three axes
at a common angle \(t\) from \(e_3\), with azimuths
\(0,2\pi/3,4\pi/3\). Their squared separation is
\(3\sin^2 t-\tfrac94\sin^4t\), and the sum of squared evaluations is
\(\tfrac34\sin^4t\). Thus the fourth power in
\eqref{secjet:separation} is sharp as \(t\to0\). Combining
\eqref{secjet:separation} and \eqref{secjet:nullestimate}
gives the sufficient factor \(96\delta^{-4}\) for the differential
density. Sharpness here concerns reconstruction from the null
evaluations.

\subsection{The quaternionic spinor description}

On a smooth spin four-manifold, choose the untwisted positive
half-spin bundle \(S^+\) so that self-dual two-forms act on it.
Use the parallel Hermitian identification \(T_{\mathbb C}M
\simeq S^-\otimes S^+\) and let \(j\) be the parallel antiunitary
quaternionic structure on \(S^+\)
\cite[\S1]{AHS1978}. For a smooth positive-chirality spinor
\(\psi\), on \(\{\psi\ne0\}\) consider the lines of
\[
 \pi_1=\psi,\qquad
 \pi_2=(\psi+j\psi)/\sqrt2,\qquad
 \pi_3=(\psi+\mathrm i j\psi)/\sqrt2 .
\]
A line \(\mathbb C\pi\) defines \(\mathcal D=S^-\otimes\mathbb C\pi\),
with conjugate plane \(S^-\otimes\mathbb C j\pi\). The
projectors of unit representatives \(u_i\) give twistor axes
with \(\langle n_i,n_k\rangle=2|\langle u_i,u_k\rangle|^2-1\).
The pairwise squared inner products of these representatives
are \(1/2\), so the axes are orthogonal.

Write \(P_{\mathrm{tw},X}\pi=\nabla_X\pi+\tfrac14 c(X)D\pi\),
where \(D\) is the Levi--Civita Dirac operator and
\(c(X)^2=-|X|^2\) for real \(X\)
\cite[eq.~(1.1)]{Friedrich1990TwistorSpinors}. For a single line
write \(u=\pi/|\pi|\), with \(p\) the orthogonal projector onto
\(\mathbb C u\) and \(q=1-p\). For \(X\in \mathcal D\) the
Clifford term \(c(X)D\pi\) belongs to \(\mathbb C u\), since
the action of \(a\otimes u\) on \(S^-\) contracts the
negative factor by its alternating form. For \(X=a\otimes u\)
and \(Y=b\otimes u\), torsion freeness gives
\[
 F_J(X,Y)=b\otimes v_X-a\otimes v_Y,\qquad
 v_X=q\nabla_Xu=|\pi|^{-1}q P_{\mathrm{tw},X}\pi .
\]
Only one exterior pair contributes when \(a,b\) form a unitary
basis. The two output terms are orthogonal, so
\(|F_J|^2\le |P_{\mathrm{tw}}\pi|^2/|\pi|^2\).
The real even Clifford action and the connection commute with
\(j\), and each spinor is orthogonal to its quaternionic image.
For each \(i\), \(|\pi_i|=|\psi|\) and
\(|P_{\mathrm{tw}}\pi_i|=|P_{\mathrm{tw}}\psi|\). Summing the
three projective bounds gives
\[
 \sum_i|F_{J_i}|^2\le 3\,\frac{|P_{\mathrm{tw}}\psi|^2}{|\psi|^2}
       \qquad\text{on }\{\psi\ne0\}.
\]
Each Frobenius tensor depends only on its projective line. The two
auxiliary lines are chosen using the supplied spinor \(\psi\).
An exact twistor spinor supplies involutive directions on its
nonzero locus. Control of \(\nabla F_{J_i}\) or of the sections
across spinor zeros requires further hypotheses.

\subsection{Projective variations on a round patch}
\label{secjet:germ-projective}

Write the round metric on a stereographic patch as \(r=\Omega^2\delta\)
and a metric variation as \(g_t=r+th+O(t^2)\). Use the fixed flat
representative, so the variation there is \(\widehat h=\Omega^{-2}h\).
Choose three constant positive orthogonal complex structures for
\(\delta\), with distinct axes modulo sign. The metric variation is
common to all three structures, while their projective lines may vary
independently.

For one structure \(J\), choose linear complex coordinates with
\(\delta=\sum_j|dz_j|^2\). Write \(U_j=\partial_{\bar z_j}\),
\(Z_j=\partial_{z_j}\), so that \(\delta(U_a,Z_b)=\delta_{ab}/2\).
The varying null plane is spanned to first order by
\(U_a+t\sum_b B_a{}^b Z_b\). Nullity gives
\((B_a{}^b+B_b{}^a)/2=-\widehat h(U_a,U_b)=H_{ab}\).
This symmetric part is the one furnished by the positive symmetric
frame transport, whose velocity is \(-\widehat h^\sharp/2\).
The remaining skew entries are \(B_1{}^2-H_{12}=v\) and
\(B_2{}^1-H_{21}=-v\), with \(v\) complex. These are the two real
degrees of freedom of the line, and reality fixes the conjugate
plane variation.

The coordinate fields commute at the baseline. The two \(Z_b\)
coefficients of the first variation of \(F_J(U_1,U_2)\), viewed as
a \((0,1)\)-form in the \(d\bar z_j\) basis, are
\begin{equation}
 \label{secjet:germ-graph}
 \mathcal F_J(h,v)=M_J(h)-\bar\partial_Jv,\qquad
 M_J(h)=(U_1H_{21}-U_2H_{11},\,U_1H_{22}-U_2H_{12}).
\end{equation}
Let \(m=U_1\wedge U_2\), of Hermitian norm \(1/2\). The coefficient
of \(d\bar z_1\wedge d\bar z_2\) in \(\bar\partial_J\mathcal F_J\) is
\begin{equation}
 \label{secjet:germ-compatibility}
 U_1^2H_{22}-2U_1U_2H_{12}+U_2^2H_{11}
 =2\langle (W^+)'_\delta[\widehat h]m,m\rangle_{\mathbb C}.
\end{equation}
Indeed the flat variation of \(g(R(U_1,U_2)U_2,U_1)\) is
\(U_1U_2\widehat h(U_1,U_2)
-\tfrac12\bigl(U_1^2\widehat h(U_2,U_2)+
U_2^2\widehat h(U_1,U_1)\bigr)\). The Schouten terms vanish
in these null slots. For the unit representative \(2m\), the right
side of \eqref{secjet:germ-compatibility} is half the corresponding
Weyl evaluation. The flat and round Weyl operator variations are
related by the nonzero factor \(\Omega^2\) in the common Hodge
splitting.

For a smooth metric germ with \((W^+)'_r[h]=0\), each \(M_{J_i}(h)\)
is \(\bar\partial_{J_i}\)-closed. The local \(\bar\partial\) Poincar\'e
lemma solves \(\bar\partial_{J_i}v_i=M_{J_i}(h)\) on small polydisks,
which can be restricted to a common neighborhood. Conversely, the
three homogeneous equations make the three null evaluations vanish
by \eqref{secjet:germ-compatibility}. Reconstruction from distinct
axes gives \((W^+)'_r[h]=0\). Each lift is determined up to an
arbitrary \(J_i\)-holomorphic function, including functions vanishing
at a prescribed base point if the value of \(v_i\) is fixed there.

For a real nonzero covector \(\xi\), the pair \(p_j=\xi(U_j)\)
is nonzero. The vertical symbol \(-vp\) is injective, with complex
line image. The quotient map \(c\mapsto p_1c_2-p_2c_1\) sends the
metric symbol to \(p_1^2H_{22}-2p_1p_2H_{12}+p_2^2H_{11}\), the
symbol of \eqref{secjet:germ-compatibility}. On the transverse
trace-free symbol slice the metric Weyl symbol identifies its five
real components with the trace-free self-dual Weyl block. The three
null evaluations embed that block into \(\mathbb C^3\), so the full
section symbol has rank \(5+6=11\) in a real twelve-dimensional
target. In cyclic unit null frames for orthogonal axes, the one
real relation on the three normalized quotient components is that
their real parts sum to zero. These statements concern smooth local
germs and their symbols. A global inverse would also have to
address the holomorphic freedom, patch overlaps, and any boundary
or support requirements.

\section{Prescribing chiral curvature}\label{sec:lift}

Let \(M\) be an oriented smooth closed connected four-manifold with a
smooth Riemannian reference metric \(r\). We use curvature-operator
Hilbert--Schmidt norms and the convention \(W^+=0\) for ASD. Set
\(\nu=dV_r\) and \(E=\Lambda_r^+\). The curvature source will be specified
in this fixed bundle. The exterior correction below applies to a reference
which is round on one coordinate patch, with its geometry elsewhere
supplied independently.

\subsection{Metric two-forms and compatibility}

An \(E\)-valued two-form \(\Sigma\) is normalized when in every oriented
orthonormal \(E\)-frame its components satisfy
\begin{equation}\label{lift:normalization}
 \Sigma_i\wedge\Sigma_j=2\delta_{ij}\nu .
\end{equation}
Choose the component in which \((\Sigma_i/\sqrt2)\) is oriented
by its projection to \(\Lambda_r^+\). This convention agrees
with the natural orientation of the reconstructed self-dual bundle.
To see the reconstruction, divide the wedge pairing on two-forms by
\(\nu\). It has signature \((3,3)\), and the span of the triple is a
maximal positive plane. The action of \(\mathrm{SL}(4,\mathbb R)\) on
two-forms surjects onto \(\mathrm{SO}_0(3,3)\), since its differential is
injective and the groups have the same dimension. This action is
transitive on maximal positive planes. The stabilizer in
\(\mathrm{SL}(4,\mathbb R)\) of the standard self-dual plane is
\(\mathrm{SO}(4)\). Polar decomposition reduces the assertion to a positive
diagonal matrix of determinant one, and preservation of the plane
equates the products of diagonal entries in every complementary pair.
Those equalities force all four entries to be one. Thus a positive
plane determines a unique metric of volume \(\nu\), smoothly in the
plane, and \(\Sigma/\sqrt2\) identifies \(E\) isometrically with
\(\Lambda_g^+\). Since oriented coframe changes act on \(\Lambda_g^+\)
through \(\mathrm{SO}(3)\) surjectively, choose a local coframe with
\(\Sigma_i=e^i\wedge e^4+\tfrac12\epsilon_{ijk}e^j\wedge e^k\).
Changing the oriented internal frame rotates the triple and changes the
connection by its usual gauge transformation. The triple itself is a
bundle-valued field~\cite[Sections~1.3 and~1.6]{Fine2011}.

Use \([T_i,T_j]=\epsilon_{ijk}T_k\) for the internal generators. For a
metric connection \(A\) on \(E\), define the matrix \(C\) by
\begin{equation}\label{lift:projection}
 C_{ij}\nu=-\tfrac12 F_A^i\wedge\Sigma_j .
\end{equation}
The equation \(d_A\Sigma=0\) makes \(A\) the transported chiral
Levi--Civita connection \(\Gamma^+(g)\). Indeed, writing
\(b=A-\Gamma^+(g)\) gives \(d_A\Sigma=[b,\Sigma]\), since
Levi--Civita compatibility gives \(d_{\Gamma^+}\Sigma=0\).
In the adapted coframe let \(\eta_{i,ab}\) be the components of
\(\Sigma_i\). Then the map \(\mathscr C b=*[b,\Sigma]\) has components
\((\mathscr C b)_{ic}=\epsilon_{ijk}\eta_{k,dc}b_{jd}\).
Contracting with
\(\eta_{i,ab}\eta_{j,cb}=\delta_{ij}\delta_{ac}
+\epsilon_{ijk}\eta_{k,ac}\) gives
\(\mathscr C^2+\mathscr C-2I=0\), so its inverse is
\((\mathscr C+I)/2\). This computation uses the normalized coframe
alone~\cite[compatibility proposition]{Baker2026ChiralMetricEvolution}.

For the sign in \eqref{lift:projection}, let \(K_\alpha\) be defined by
\(g(K_\alpha X,Y)=\alpha(X,Y)\). The generators
\(K_{\Sigma_i}/2\) have the stated brackets, and the \(\Lambda^-\)
summand commutes with their action on \(\Lambda^+\). With positive
sectional curvature on the round sphere, the two-form of \(R(X,Y)\)
is \(-\mathcal R(X^\flat\wedge Y^\flat)\), where \(\mathcal R\) is the
curvature operator. Its projection gives
\(F_{\Gamma^+}^i=-\mathcal R\Sigma_i\), since \(|\Sigma_i|^2=2\)
and the generators are \(K_{\Sigma_i}/2\). Consequently
\[
 F_{\Gamma^+}=-(W^++s_g I/12)\Sigma-B_{\mathrm{Ric}}\Sigma^- ,
\]
with \(B_{\mathrm{Ric}}\) the off-diagonal trace-free Ricci block
in this row convention and \(\Sigma^-\) an equally normalized
anti-self-dual frame. The \(\Sigma^-\) term drops out of
\eqref{lift:projection}, giving \(C=W^++s_g I/12\).
Thus \(d_A\Sigma=0\) and a prescribed \(\operatorname{Sym}_0 C\)
give a first-order mixed system whose elimination is the second-order
metric equation for \(W^+\). Positivity is imposed on the metric
triple, and the construction admits all smooth Riemannian metrics,
including ones with degenerate chiral curvature~\cite[Section~1]{GoverGursky2024}
\cite[chiral curvature decomposition]{Baker2026ChiralMetricEvolution}.

\subsection{Fixed target and the linearized range}

Write \(V=S^2_0T^*M\), \(B=S^2_0E\), and \(h^\sharp=r^{-1}h\).
Use the fixed-volume chart and symmetric comparison
\begin{equation}\label{lift:chart}
 g_h=r(e^{h^\sharp}\cdot,\cdot),\quad
 S_h=e^{-h^\sharp/2}:(TM,r)\longrightarrow(TM,g_h),\quad
 F_r(h)=I_hW^+_{g_h}I_h^{-1},\qquad
 I_h=\Lambda^2S_h^*\big|_{\Lambda_{g_h}^+}.
\end{equation}
The metric is positive and \(dV_{g_h}=\nu\) because
\(\operatorname{tr}_r h=0\).
The isometry \(S_h\) preserves orientation and makes \(I_h\) a fiber
isometry. All local expressions agree across the fixed bundle
transitions. Put \(W=F_r(0)=W_r^+\) and \(L=dF_r(0)\). In a fixed atlas,
\(F_r\) is a smooth function of the two-jet of \(h\), with background
coefficients determined by \(r\). In particular it is smooth from
\(H^{s+2}(V)\) to \(H^s(B)\) for any fixed integer \(s\ge5\), and from
\(C^{k+2,\alpha}(V)\) to \(C^{k,\alpha}(B)\) for \(k\ge2\),
\(0<\alpha<1\), by multiplication and composition in those spaces.

Let \(\mathcal KX=(\mathcal L_X r)_0\). For a flow
\(\varphi_\epsilon\) generated by \(X\), the metric
\(g_\epsilon=\rho_\epsilon\varphi_\epsilon^*r\),
\(\rho_\epsilon=(\nu/\varphi_\epsilon^*\nu)^{1/2}\), has volume
\(\nu\) and derivative \(\mathcal KX\). The comparison from \(r\) to
itself at \(\varphi_\epsilon(x)\) is the isometry
\(O_\epsilon=\rho_\epsilon^{1/2}d\varphi_\epsilon S_\epsilon\),
where \(S_\epsilon\) is the map in \eqref{lift:chart} for \(g_\epsilon\).
Conformal covariance of the Weyl endomorphism gives
\(\rho_\epsilon^{-1} O_\epsilon^*W(\varphi_\epsilon(x))(O_\epsilon^*)^{-1}\)
for the transported field. The derivative of \(S_\epsilon\) cancels
the symmetric trace-free part of \(\nabla X\), and that of
\(\rho_\epsilon^{1/2}\) cancels its trace. Thus if \(\Theta_X\) is the
infinitesimal pullback on two-forms of \(\operatorname{skew}\nabla X\),
\begin{equation}\label{lift:gauge}
 L\mathcal KX=\nabla_X W+[\Theta_X,W]+\tfrac12(\operatorname{div}_r X)W .
\end{equation}
The normalized pullback also transforms a transported source. A source
fixed as a section of \(B\) can therefore constrain these directions
off ASD~\cite[discussion after (6.7)]{FeffermanGraham2008Ambient}.

The principal symbol is computed by freezing a flat orthonormal
coframe. In the real symbol convention \(\partial\mapsto\xi\), take
\(\xi=|\xi|e^0\). A symbol transverse-traceless tensor has
\(h_{0a}=0\) and a trace-free \(3\times3\) spatial block. The curvature
variation then has only entries of type
\(\mathcal R'_{0i,0j}=-|\xi|^2h_{ij}/2\). Projection to the orthonormal
self-dual frame gives \(-|\xi|^2(h_{ij})/4\), and the scalar symbol is zero.
The orthogonal complement of this five-dimensional domain fiber is
\(\operatorname{im}\sigma_\xi(\mathcal K)\), since
\(\mathcal K^*=-2\delta_r\), with
\((\delta_r h)_j=\nabla^i h_{ij}\). The symbol of \(\mathcal K\) is
injective for \(\xi\ne0\), and \eqref{lift:gauge} shows that \(L\)
annihilates its image at principal order. In particular
\begin{equation}\label{lift:normal}
 P=LL^*,\qquad \sigma_\xi(P)=|\xi|_r^4I_B/16 .
\end{equation}
This is an elliptic self-adjoint operator on the closed manifold.
Let \(\mathcal C=\ker L^*=\ker P\), the equality following by pairing
\(Pu\) against \(u\). Elliptic regularity makes \(\mathcal C\) finite
dimensional and smooth. The generalized inverse \(G_P\) maps
\(H^s(B)\) to \(H^{s+4}(B)\cap\mathcal C^\perp\) and satisfies
\(PG_P=I-\pi_{\mathcal C}\), with the reference \(L^2\) projection.
Thus the full range of \(L\) is \(H^s(B)\cap\mathcal C^\perp\), and
\begin{equation}\label{lift:split}
 H^{s+2}(V)=\ker L\ \oplus\
 L^*(H^{s+4}(B)\cap\mathcal C^\perp).
\end{equation}
The projection onto the second summand is \(L^*G_P L\).
When \(\mathcal C=0\), \(Q=L^*P^{-1}\) is a bounded right inverse
gaining two derivatives. At any reference a differentiable lift of
\(F_r(h(t))=(1-t)W\) through \(h(0)=0\) requires \(W\perp\mathcal C\),
by its tangent equation \(Lh'(0)=-W\).

\paragraph{The fixed transverse-traceless slice.}
Write \(\mathrm{TT}_r=\ker\delta_r\) in \(V\). On a homotopy
four-sphere the Fredholm index of \(L|_{\mathrm{TT}_r}\) is
\begin{equation}\label{lift:index}
 \operatorname{ind}(L|_{\mathrm{TT}_r})
     =\dim\ker\mathcal K-15 .
\end{equation}
For the index, the operator \((\delta_r,L)\) is elliptic with row
orders one and two, as is seen by multiplying the divergence row
by an elliptic isomorphism of order one with symbol \(|\xi|I\).
The adjoint of this augmented symbol folds the universal conformal
deformation symbol from the even bundles to the odd bundle. Its
index is \((15\chi+29\tau)/2\). For clarity, this number uses the
symbol even at an off-ASD reference. On the spin lift, let \(x,y\)
be the half-spin representations, \(x\) corresponding to
\(\Lambda^+\). The even-minus-odd representation is
\[
 xy+\operatorname{Sym}^4x-\operatorname{Sym}^2x\,\operatorname{Sym}^2y
 =(x-y)(\operatorname{Sym}^3x+\operatorname{Sym}^2x\,y).
\]
With half-spin weights \(\pm u,\pm v\), the oriented Euler and
Pontryagin classes are \(e=u^2-v^2\), \(p_1=2(u^2+v^2)\).
The Atiyah--Singer character formula divides the character of this
difference by \(e\) before evaluating on \(M\) and multiplies by
\(\operatorname{Td}(TM_{\mathbb C})\). The degree-four term is
\(18u^2+3v^2-10p_1/24=(15/2)e+(29/6)p_1\). The augmented operator
has the opposite index, hence \(-15\) at \(\chi=2,\tau=0\).
The divergence has range \((\ker\mathcal K)^\perp\), using the elliptic
inverse of \(\delta_r\mathcal K=-\mathcal K^*\mathcal K/2\) on that
complement. If \(R_\delta\) is its resulting right inverse for
\(\delta_r\), remove the \(\ker\mathcal K\) part of the first target
and subtract \(LR_\delta\) of the first output from the second.
The resulting operator splits into the identity on the divergence
range and \(L|_{\mathrm{TT}_r}\). This proves \eqref{lift:index}~\cite[Section~2]{AtiyahSinger1968IndexIII}
\cite[Section~1.1]{GoverGursky2024}.

The conformal Killing equation determines a field by its value,
skew derivative, divergence and derivative of the divergence at
one point. Indeed the variations of the Levi--Civita connection
and the Schouten tensor under \(\mathcal L_X r=2fr\) express
\(\nabla^2 X\) in these data and \(\nabla^2 f\) in \(X,\nabla X\).
They give a homogeneous first-order system along curves.
Thus \(\dim\ker\mathcal K\le4+6+1+4=15\). Equality would realize
every skew first jet at a point with \(X=0\) as a conformal Killing
field. By \eqref{lift:gauge}, this would make \(W\) commute with all
rotations of the positive bundle, forcing \(W=0\). Hence the fixed
slice has nontrivial cokernel at every off-ASD reference on a
homotopy sphere. The index leaves undecided membership of the
particular target \(-W\) in that range.

\paragraph{Round inverse and gap.}
On the standard round sphere \(r_0\) of curvature \(K>0\), write
\(D_0=L_{r_0}\) and \(\Delta_\nabla=\nabla^*\nabla\). The round
factorization on TT tensors in the operator norm convention is
\[
 D_0^*D_0=\tfrac1{16}(\Delta_\nabla+2K)(\Delta_\nabla+4K).
\]
The rough Laplacian preserves TT. For a TT tensor, integrating the
cross terms in the square of
\(S_{abc}=\nabla_a h_{bc}+\nabla_b h_{ca}+\nabla_c h_{ab}\)
and using \([\nabla^a,\nabla_b]h_{ca}=4Kh_{bc}\) gives
\(\|S\|_2^2=3\|\nabla h\|_2^2-24K\|h\|_2^2\).
Thus \(\Delta_\nabla\ge8K\) and
\(\|D_0h\|_2^2\ge(15/2)K^2\|h\|_2^2\).
The round ASD cokernel vanishes and \(D_0\mathcal K=0\). Thus \(Q_0=D_0^*(D_0D_0^*)^{-1}\) is the inverse on the round
TT slice with gain two~\cite[round metric Hessian and TT gap]{Baker2026ChiralMetricEvolution}
\cite[Section~1.1, round unobstructedness]{GoverGursky2024}.

\paragraph{The two sourced equations on the round sphere.}
At the round metric write \(S(h)=\dot\Sigma\), \(\mathcal B a=[a,\Sigma]\),
and \(\mathcal P u=\operatorname{Sym}_0((u^i\wedge\Sigma_j)/(2\nu))_{ij}\).
In the symmetric frame,
\(S_i(h)(X,Y)=\tfrac12(\Sigma_i(h^\sharp X,Y)+\Sigma_i(X,h^\sharp Y))\)
belongs to \(\Lambda^-\). Indeed the self-dual coefficient matrix is
symmetric by the frame choice and skew by differentiated wedge normalization.
This map from trace-free symmetric tensors to \(E\otimes\Lambda^-\)
is an isomorphism: in a diagonal eigenframe, the \(i4\) coefficient
is \((h_{ii}+h_{44})/2\), so vanishing of all three forces \(h=0\);
both spaces have dimension nine. For \(\Gamma=\Gamma^+(r)\) set
\(\mathcal Hh=-\mathcal B^{-1}d_\Gamma S(h)\).
Given compatibility and chiral curvature sources \(u,f\), respectively,
the independent connection variation \(a\) satisfies
\[
 d_\Gamma S(h)+\mathcal B a=u,\qquad -\mathcal P d_\Gamma a=f .
\]
The term \(F_\Gamma^i\wedge S_j(h)\) in the differentiated curvature
projection vanishes at the round metric. With
\(\beta=a-\mathcal Hh\) and \(D_0=-\mathcal P d_\Gamma\mathcal H\), the
TT solution and the connection variation are
\[
 \beta=\mathcal B^{-1}u,\qquad
 h=Q_0(f+\mathcal P d_\Gamma\mathcal B^{-1}u),\qquad
 a=\mathcal Hh+\mathcal B^{-1}u .
\]
Thus \(u\in H^{s+1}\) and \(f\in H^s\) give \(h\in H^{s+2}\) and
\(a\in H^{s+1}\) continuously. At the corresponding mixed orders
the symbol is an isomorphism on the TT slice, with twelve algebraic
connection components and five metric components. Before imposing TT,
its kernel is the four-dimensional graph of
\(\sigma_\xi(\mathcal H)\) over \(\operatorname{im}\sigma_\xi(\mathcal K)\).
Assigning equal order to both fields drops the algebraic term
\(\mathcal B a\). For \(\xi\parallel e^4\), \(\xi\wedge S_i(h)=0\)
forces the anti-self-dual form \(S_i(h)\) to vanish, hence \(h=0\).
The remaining connection symbol has a seven-dimensional kernel:
the three \(a^i_4\) and the trace and skew parts of \((a^i_j)_{i,j\le3}\).

\paragraph{The missing transverse-traceless rows.}
At a general reference let \(\mathcal C_T\) be the smooth
annihilator of the TT image, and \(\Pi_T\) the \(L^2\) projection
of metric tensors onto TT. The operator \(L\Pi_TL^*\) is elliptic
with the symbol of \(LL^*\), since \(\Pi_T\) projects onto the TT
symbol slice. Its generalized inverse therefore gives TT solutions
for precisely the targets orthogonal to \(\mathcal C_T=\ker(\Pi_T L^*)\).
The York splitting yields the finite-dimensional criterion
\[
 L\ {\rm onto}\quad\Longleftrightarrow\quad
 \operatorname{im}(\pi_{\mathcal C_T}L\mathcal K)=\mathcal C_T .
\]
When this holds, at most \(\dim\mathcal C_T\) smooth Lie variations
supply all the missing rows: take a smooth preimage for each member
of a basis and decompose it into its TT and Lie parts. Its action
on curvature is the one in \eqref{lift:gauge}.

\subsection{A quantitative one-chart lift}

\begin{proposition}[A whole curvature chord in one metric chart]\label{lift:radial}
Let \(X,Y\) be the metric and target Banach spaces at either of the
grades stated above. Suppose \(L\) has a bounded right inverse
\(Q:Y\to X\), and \(DF_r\) is \(H\)-Lipschitz in the radius-\(\rho\)
ball in \(X\). Set \(\eta=\|QW\|\), \(q=\|Q\|H\rho\).
If \(q<1\) and \(\eta+q\rho/2\le\rho\), there is a unique
path in the closed radius-\(\rho\) ball of \(\operatorname{im}Q\)
with \(F_r(h(t))=(1-t)W\), \(0\le t\le1\). It is smooth in the
fixed Banach grade, and
\[
 \|h(t)-h(u)\|\le\eta|t-u|/(1-q),\qquad
 \|h(t)\|\le t\eta/(1-q/2).
\]
\end{proposition}

\begin{proof}
For \(N(h)=F_r(h)-W-Lh\), integration on segments gives
\(\|N(h)\|\le H\|h\|^2/2\) and Lipschitz constant \(H\rho\)
on the ball. The map \(h\mapsto-tQW-QN(h)\) is a
\(q\)-contraction of the stated ball into itself. The ball is
complete because \(QL\) is a bounded projection onto
\(\operatorname{im}Q\). Its fixed point solves the equation, and
conversely \(QL=I\) on this range. The two bounds follow by subtracting
fixed point equations and by estimating the quadratic remainder,
respectively. The derivative of the fixed point equation on
this range is the identity plus an operator of norm at most \(q\).
The implicit function theorem gives parameter smoothness.
The chart bound gives two-sided uniform comparison to \(r\) by
exponentiating the pointwise eigenvalue bounds on \(h^\sharp\),
as well as the coefficient bounds in the given grade. These bounds also
give a positive lower injectivity radius on this chart ball. The
curvature is bounded by the coefficient estimates, and the two-sided
metric comparison gives a uniform lower volume bound at a fixed small
radius. The injectivity estimate applies at that scale
\cite[Theorem 4.7]{CheegerGromovTaylor1982}.
\end{proof}

On the known standard \(S^4\), this criterion holds at every sufficiently
close smooth reference to the unit round metric in \(C^{k+2,\alpha}\).
Identify the metric and Weyl targets with those of \(r_0\) by the
positive symmetric isometry between \(r_0\) and \(r\), and put tildes
on the compared operators. The smooth finite-jet expression gives
\(\|\widetilde L_r-D_0\|+\|\widetilde W\|
\le C\|r-r_0\|_{C^{k+2,\alpha}}\) and a common local Lipschitz bound
for the derivative. The inverse
\(Q_0(\widetilde L_rQ_0)^{-1}\) exists by a Neumann series,
has bounded norm and a fixed range. Choose a common radius with
\(q\le1/2\), then shrink the reference neighborhood to make
\(\eta\le\rho/2\). Smooth spatial regularity in this differential
range is verified below. An abstract right inverse in
Proposition~\ref{lift:radial} gives the stated Banach regularity.

\subsection{Exterior corrections on the given manifold}

\begin{theorem}[Local prescription with exterior corrections]\label{lift:exterior}
Suppose \(r\) is isometric, orientation-preservingly, to the round
model on a coordinate ball \(O\subset M\), and let \(K\Subset O\)
be a closed coordinate subball. Take a reference \(L^2\)-orthonormal
basis \(Z_1,\ldots,Z_m\) of \(\mathcal C=\ker L^*\) and a smooth
function \(w\) zero on \(K\) and positive on \(M\setminus K\).
Let
\[
 \mathcal Wc=\sum_i c_i wZ_i,\quad
 b(y)_i=\langle Z_i,y\rangle_{L^2(r)},\quad
 \Gamma_{ij}=\langle Z_i,wZ_j\rangle_{L^2(r)} .
\]
Then \(\Gamma\) is positive definite. For any fixed \(s\ge5\), every
sufficiently small \(y\in H^s(B)\) has a unique small pair \((u,c)\)
in \((H^{s+4}(B)\cap\mathcal C^\perp)\times\mathbb R^m\) giving
\begin{equation}\label{lift:corrected}
 h=L^*u,\qquad F_r(h)=W+y-\mathcal Wc,\qquad
 c=\Gamma^{-1}b(y)+O(\|y\|_{H^s}^2).
\end{equation}
For smooth \(y\), the metric and its compatible chiral connection
are smooth. In particular \(F_r(h)=y\) on \(K\).
\end{theorem}

\begin{proof}
Any \(Z\in\mathcal C\) vanishing on \(M\setminus K\) can be
transplanted from the round patch to the standard round sphere
and extended by zero. The extension is smooth because it vanishes
in a collar before the patch boundary, and it solves the round
adjoint equation since the differential operator and pairings
agree there. Round adjoint injectivity gives \(Z=0\).
Positive definiteness of \(\Gamma\) follows by writing its quadratic
form as \(\int_M w|\sum_i c_i Z_i|^2\,dV_r\). If it vanishes,
the combination vanishes on the exterior. For example \(w\) can
be chosen flat at \(\partial K\) by using \(e^{-1/d^2}\) for
positive outward collar parameter \(d\) and joining to a strictly
positive function farther out. Its support may touch the boundary.
A slightly larger subball in \(O\) gives the same argument with an
intervening collar of zero correction.

The derivative at zero of
\((u,c)\mapsto F_r(L^*u)-W+\mathcal Wc\) is \(P\dot u+\mathcal W\dot c\).
For output \(z\) its inverse is
\[
 \dot c=\Gamma^{-1}b(z),\qquad
 \dot u=G_P\bigl(z-\mathcal W\Gamma^{-1}b(z)\bigr).
\]
The target for \(G_P\) is orthogonal to \(\mathcal C\), and the
inverse is bounded in the displayed grades. The Banach inverse
function theorem gives the local lift and the expansion of the
coefficient. The sign of the realized exterior change is
\(-\mathcal Wc\). If \(\mathcal C=0\), the finite moment and correction
rows are empty. The metric has volume \(\nu\), and the restrictions
of the global metric and compatible connection have matching jets
in common coordinates across the interface. The prescribed datum
on \(K\) is the curvature source. Its realization can change the
metric on \(K\).
\end{proof}

\paragraph{Spatial regularity.}
For the exterior theorem, work on the full potential space and
adjoin the finite row \(b(u)=0\). The derivative at zero sends
\((v,d)\) to \((Pv+\mathcal Wd,b(v))\). Its inverse is the one
above with the addition of \(\sum_i a_i Z_i\) to \(v\) for a
prescribed moment output \(a\). This derivative remains invertible
at the small solution. The fourth-order spatial row is a
finite-jet differential map. Conjugate it by local component
pullbacks \(T_\epsilon\) that are translations on a smaller
framed ball and the identity outside a larger one. Conjugation
makes the coefficients smooth in the translation parameter at
the same Sobolev grades. The moment row is \(b(T_\epsilon^{-1}v)\)
and is smooth in \(\epsilon\) by transferring the pullback to
the smooth basis and reference density. For smooth \(y\) the translated
source is smooth in the parameter. The implicit function theorem
and strong continuity of \(T_\epsilon\) identify \(T_\epsilon u\)
with the smooth solution branch. Differentiating translations yields
all local spatial derivatives in the frozen Sobolev grade.

The round neighborhood uses the same argument without finite rows
in the fixed potential range \(h=D_0^*u\) in compared coordinates.
The derivative of its potential equation, composed with
\((D_0D_0^*)^{-1}\), equals
\(D\widetilde F_r(h)\widetilde Q_r\,\widetilde L_rQ_0\), where
\(\widetilde Q_r=Q_0(\widetilde L_rQ_0)^{-1}\). Both factors
\(D\widetilde F_r(h)\widetilde Q_r\) and \(\widetilde L_rQ_0\)
are invertible by the two Neumann series.
For the H\"older proof restrict to the closures of smooth sections
in those grades. Elliptic inverses, the Neumann series and the
contraction iterates from zero preserve these closures, on which
translations are strongly continuous. This gives spatial smoothness
for smooth references and sources. The same differential-potential
argument applies to \(Q=L^*(LL^*)^{-1}\) at a smooth reference
where it is defined and satisfies Proposition~\ref{lift:radial}.

\paragraph{The exact balance.}
On the slice of Theorem~\ref{lift:exterior}, taking moments gives
\begin{equation}\label{lift:balance}
 \Gamma c(y)=b\bigl(y-N(h(y))\bigr),\qquad
 N(h)=F_r(h)-W-Lh .
\end{equation}
The corrected solution is an uncorrected lift exactly when this
right-hand side vanishes. For \(y=-tW\) the linear coefficient is
\(-t\Gamma^{-1}b(W)\), with \(W\) including the original exterior
Weyl curvature. The splitting \eqref{lift:split} also permits
\(h=k+L^*u\) with small \(k\in\ker L\) as a parameter in the augmented
implicit function theorem. The corresponding uncorrected condition
uses its own \(N(k+L^*u)\). Thus a failure of balance on the chosen
slice leaves the other kernel parameters open.

If \(M\) is a homotopy four-sphere, the actual exterior
\(M\setminus\operatorname{int}K\) is contractible with boundary \(S^3\).
Van Kampen gives simple connectivity by using a collar along the
boundary. Mayer--Vietoris gives reduced acyclicity, since the
connecting map sends the fundamental class of \(M\) to the
boundary orientation class. The smooth exterior has CW type,
so the simply connected homology form of Whitehead's theorem
gives contractibility. The proof of Theorem~\ref{lift:exterior}
uses its given smooth structure and a fixed reference. The next two sections treat the exact balance, first at a nearby
reference of the same energy and then at the original reference.

\section{Exact prescription after changing the exterior reference}
\label{msreg:section}

The finite corrections in Theorem~\ref{lift:exterior} compensate
for the cokernel at one fixed reference. Changing the reference in
the exterior makes the full curvature derivative surjective. The
change can preserve the initial Weyl energy exactly. We first show
how to detect the missing directions and then estimate the short
radial step in the new chart.

Let \(M\) be a smooth closed oriented homotopy four-sphere and \(r\)
a smooth positive Riemannian metric. Let \(K\) be a smooth closed
coordinate ball with a marked round metric, and suppose \(r\)
agrees with that metric on \(K\). The agreement is allowed to
hold only on the closed ball. Use the density \(\nu=dV_r\) and
the bundles and chart of \eqref{lift:chart}. Fix an integer \(s\ge5\)
and write \(X=H^{s+2}(V)\), \(Y=H^s(B)\), and
\[
 U=H^{s+4}(B)\cap\mathcal C^\perp,\qquad
 m=\dim\mathcal C,\qquad b(f)_i=\langle Z_i,f\rangle_{L^2(r)},
\]
where \(Z_1,\ldots,Z_m\) is an \(L^2\)-orthonormal basis of
\(\mathcal C=\ker L^*\). Call a reference regular if its full
trace-free derivative is onto. By \eqref{lift:normal}, this is
equivalent to triviality of its smooth adjoint kernel and is
independent of \(s\).

\begin{theorem}[Exact prescription at the same reference energy]\label{msreg:density}
Given \(r\) and \(K\) as above, arbitrarily close to \(r\) in
\(C^\infty\) there is a regular reference \(\tilde r\) agreeing with
\(r\) on a neighborhood of \(K\), with
\(dV_{\tilde r}=dV_r\) and \(E_+(\tilde r)=E_+(r)\).
In the fixed-density chart centered at \(\tilde r\), every sufficiently
small \(H^s\) target increment is attained. In particular a smooth
radial chord of positive duration starts at \(\tilde r\), with energy
\((1-t)^2E_+(r)\) and \(W^+=0\) on \(K\).
\end{theorem}

\subsection{Comparing the derivatives}

For a smooth logarithmic displacement \(p\in\Gamma(V)\), define
\(r_p=r(e^{p^\sharp}\cdot,\cdot)\). The positive symmetric
isometry \(S_p=e^{-p^\sharp/2}:(TM,r)\to(TM,r_p)\)
induces \(J_p h=(S_p^{-1})^*h\) on trace-free metric tensors
and \(I_p=\Lambda^2S_p^*|_{\Lambda^+_{r_p}}\) on positive two-forms.
Define the compared map and derivative by
\begin{equation}
 \widehat F(p,h)=I_pF_{r_p}(J_ph)I_p^{-1},
 \qquad A(p)=D_h\widehat F(p,0).
 \label{msreg:comparison}
\end{equation}
The metric in this expression is
\[
 g(p,h)=r\bigl(e^{p^\sharp/2}e^{h^\sharp}e^{p^\sharp/2}\cdot,\cdot\bigr).
\]
Its volume form is \(\nu\). The metric and target comparisons are
isometric for the fixed-density \(L^2\) pairings. The map
\(\widehat F\) is smooth from a neighborhood of \((0,0)\) in
\(X\times X\) to \(Y\). In a fixed atlas the curvature is linear in
second metric derivatives and otherwise quadratic in first
derivatives, with coefficients smooth in the metric and its inverse.
The projections, exponentials and isometries are smooth pointwise
operations. Sobolev multiplication and composition give the stated
smoothness with two derivatives lost. In particular \(A(p)\) is
smooth into \(\mathcal L(X,Y)\), \(A(0)=L\), and
\(\widehat F(p,0)=F_r(p)\). For a smooth path \(p(a)\) through zero
write \(A_a=A(p(a))\) and \(W_a=\widehat F(p(a),0)\).

Suppose smooth \(k_1,\ldots,k_m\in\ker L\) give an invertible matrix
\begin{equation}
 \mathsf M_{ij}=\langle Z_i,A'_0 k_j\rangle_{L^2(r)} .
 \label{msreg:mixedmatrix}
\end{equation}
Let \(\mathcal J c=\sum_j c_j k_j\). The operator
\begin{equation}
 T_a(u,c)=A_a L^*u+\frac{A_a\mathcal J}{a}c,
 \qquad (u,c)\in U\times\mathbb R^m ,
 \label{msreg:columns}
\end{equation}
extends smoothly across \(a=0\), since
\(A_a\mathcal J/a=\int_0^1 A'_{\theta a}\mathcal J\,d\theta\).
Write \(D_0=A'_0\mathcal J\). At \(a=0\) the inverse for output \(f\)
is given by
\[
 c=\mathsf M^{-1}b(f),\qquad u=G_P(f-D_0c).
\]
The argument of \(G_P\) lies in \(\mathcal C^\perp\), where the
normal operator in \eqref{lift:normal} is invertible, with gain
four~\cite[(1)--(2)]{Melrose2017Lecture22}. Smoothness gives
\(\|T_a-T_0\|\le C_0|a|\). In particular
\(\|T_a^{-1}\|\le2\|T_0^{-1}\|\) once
\(C_0|a|\|T_0^{-1}\|\le 1/2\). For \(a\ne0\) this yields the
right inverse
\begin{equation}
 Q_a f=L^*u+a^{-1}\mathcal J c,\qquad (u,c)=T_a^{-1}f,\qquad
 \|Q_a\|\le 2\|T_0^{-1}\|\bigl(\|L^*\|+|a|^{-1}\|\mathcal J\|\bigr).
 \label{msreg:inverse}
\end{equation}
Thus the new reference is regular. If \(m>0\) this power of
\(|a|\) is optimal over all target directions. For any \(A_a h=f\),
\begin{equation}
 |b(f)|=|b((A_a-L)h)|\le C_1|a|\,\|h\|_X.
 \label{msreg:lower}
\end{equation}
Taking \(f=Z_1\) bounds the norm of every right inverse below
by a positive multiple of \(|a|^{-1}\). When \(m=0\) the finite
columns are empty and the inverse is bounded through \(a=0\).

\subsection{An exterior perturbation of full rank}

Restriction of \(\mathcal C\) to the open exterior \(M\setminus K\)
is injective. If \(Z\in\mathcal C\) vanishes there, all its jets
vanish at \(\partial K\) by smoothness. Transplant \(Z|_K\) by
the round marking and extend by zero on the standard round sphere.
The extension is smooth by the vanishing jets. The adjoint
equation on the interior is the round equation because the
metrics, coefficients and pairings agree there. The equation also
holds outside and hence at the boundary by continuity. The
round adjoint has zero kernel, so the extension vanishes.
This uses round unobstructedness only on the standard
sphere~\cite[\S1.1]{GoverGursky2024}.

Suppose \(m>0\). The polarized evaluations
\(Z\mapsto\langle Z(x),Y\rangle_r\) for exterior \(x\) and
\(Y\in B_x\) separate \(\mathcal C\), so \(m\) of them give
a basis of its dual. The centers can be made distinct by
small displacements and smooth extension of their polarizations.
Fix disjoint relatively compact exterior coordinate balls
\(U_j\) around these centers. Extend the polarizations to smooth
local sections \(Y_j\). Choose real \(\beta_j\in C_c^\infty(U_j)\)
with \(\int \beta_j^2\nu=1\) such that the matrix
\begin{equation}
 D_{ij}=\int_{U_j}\beta_j^2\langle Z_i,Y_j\rangle_r\,\nu
 \label{msreg:evalmatrix}
\end{equation}
is invertible. This is possible because shrinking the supports
makes the entries arbitrarily close to those of the selected
evaluation matrix.

Let \(\phi_j\) be real coordinate phases with \(d\phi_j\ne0\)
on \(U_j\). On the symbol transverse-traceless subspace described
before \eqref{lift:normal}, \(-\sigma_{d\phi_j}(L)\) is an
isomorphism onto \(B\). Choose the smooth inverse images
\(H_j\) of \(Y_j\) under this symbol. For integers \(N\to\infty\)
define smooth fields by extension by zero from these balls:
\begin{equation}
 v_N=N^{-2}\sum_j\beta_j\cos(N\phi_j)H_j,\qquad
 X_{j,N}=N^{-1}\beta_j\sin(N\phi_j)
                \frac{\operatorname{grad}_r\phi_j}{|d\phi_j|_r^2}.
 \label{msreg:probes}
\end{equation}
Writing \(w_N=Lv_N\), the fixed-support estimates are
\begin{align}
 \|v_N\|_{C^\ell}&=O(N^{\ell-2})\quad(0\le\ell\le2),
 &w_N|_{U_j}&=\beta_j\cos(N\phi_j)Y_j+O(N^{-1}),\notag\\
 \|X_{j,N}\|_{C^0}&=O(N^{-1}),
 &\|\nabla X_{j,N}\|_{C^0}&=O(1),\notag\\
 \operatorname{div}_\nu X_{j,N}&=\beta_j\cos(N\phi_j)+O(N^{-1}),
 &\operatorname{skew}\nabla X_{j,N}&=O(N^{-1}).
 \label{msreg:oscest}
\end{align}
Here the tensor errors are in \(C^0\). The sign in the expression
for \(w_N\) comes from the two derivatives of the cosine in the
real symbol convention. The leading derivative of \(X_{j,N}\)
is a symmetric rank-one endomorphism with the displayed trace.

Take \(p(a)=av_N\) for the moment. In the compared bundles let \(\mathcal K_a X=J_{p(a)}^{-1}(\mathcal L_Xr_{p(a)})_0\),
with trace taken in \(r_{p(a)}\). Let \(\nabla^a\) denote the compared connection on \(B\), and
let \(\Theta_X^a\) denote the compared infinitesimal skew action
on two-forms. The identity \eqref{lift:gauge} at the moved
reference gives
\[
 A_a\mathcal K_a X=\nabla_X^a W_a+[\Theta_X^a,W_a]
                         +\tfrac12(\operatorname{div}_\nu X)W_a .
\]
The divergence is independent of \(a\) because the density
is fixed. Also \(W'_0=w_N\). Differentiating at zero and pairing
against \(Z\in\mathcal C\) eliminates \(L\dot{\mathcal K}X\),
and gives
\begin{align}
 \langle Z,A'_0\mathcal KX\rangle_{L^2}
  =\big\langle Z,\nabla_X w_N+[\Theta_X,w_N]
          +\tfrac12(\operatorname{div}_\nu X)w_N
          +\dot\nabla_X W+[\dot\Theta_X,W]\big\rangle_{L^2}.
 \label{msreg:differentiate}
\end{align}
The dots on the connection and skew action are zeroth-order
operators on \(W\). Varying the connection and the pointwise
frame gives
\[
 |\dot\nabla_X W|\le C |X|(|v_N|+|\nabla v_N|),\qquad
 |\dot\Theta_X|\le C\bigl(|v_N||\nabla X|
                      +|X|(|v_N|+|\nabla v_N|)\bigr).
\]
For \(X=X_{j,N}\) these terms are \(O(N^{-2})\), and
\([\Theta_X,w_N]=O(N^{-1})\) by \eqref{msreg:oscest}.
For the transport term, integration by parts gives
\[
 \langle Z,\nabla_X w_N\rangle_{L^2}
  =-\int_M(\operatorname{div}_\nu X)\langle Z,w_N\rangle_r\,\nu
     -\int_M\langle \nabla_X Z,w_N\rangle_r\,\nu .
\]
The last integral is \(O(N^{-1})\). The transport term
therefore leaves a coefficient \(-1/2\) after adding the
half-divergence in \eqref{msreg:differentiate}. The supports
are disjoint, so \(X_{j,N}\) sees only the \(j\)-th probe.
Consequently
\[
 \langle Z,A'_0\mathcal KX_{j,N}\rangle_{L^2}
 =-\tfrac12\int_{U_j}\beta_j^2\cos^2(N\phi_j)
                       \langle Z,Y_j\rangle_r\,\nu+O(N^{-1}).
\]
The oscillatory part of \(\cos^2(N\phi_j)\) tends to zero
after integration by parts along
\(\operatorname{grad}_r\phi_j/|d\phi_j|_r^2\).
The resulting limit for \(Z_i\) is \(-D_{ij}/4\).

The vectors \(\mathcal K X_{j,N}\) must be projected to the
exact kernel of \(L\). The identity \eqref{lift:gauge} and
\eqref{msreg:oscest} show that
\[
 L\mathcal K X_{j,N}
     =\tfrac12\beta_j\cos(N\phi_j)W+O(N^{-1})
       \rightharpoonup0 \quad\text{in }L^2 .
\]
Weak convergence follows first against smooth sections by
integration by parts, then against all \(L^2\) sections by
density and the uniform bound. The map \(L^*G_P:L^2(B)\to H^2(V)\)
is bounded and hence compact into \(L^2(V)\)~\cite[Theorem 16.10]{Dyatlov2026Elliptic}.
Thus the correction
\(q_{j,N}=L^*G_P L\mathcal K X_{j,N}\) tends to zero in \(L^2\).
Set \(k_{j,N}=\mathcal KX_{j,N}-q_{j,N}\).
The tensors are smooth by elliptic regularity, have uniformly
bounded \(L^2\) norm, and satisfy \(Lk_{j,N}=0\) exactly by
\eqref{lift:split}.

To control the correction in the mixed matrix, the coefficient
of order \(\ell\) in \(h\) in \(A'_0 h\) is linear in derivatives
of \(v_N\) of order at most \(2-\ell\) for \(0\le\ell\le2\).
This follows from the two-jet expression for curvature in
\eqref{msreg:comparison}: second metric derivatives occur
linearly with algebraic coefficients, and first derivatives
occur at most quadratically. Differentiating the pointwise
comparisons preserves this total derivative order. Taking the
formal adjoint differentiates each coefficient at most
\(\ell\) more times, so for the fixed \(Z_i\)
\[
 \|(A'_0)^*Z_i\|_{L^2}\le C_i\|v_N\|_{C^2}\le C_i'.
\]
It follows that \(\langle Z_i,A'_0 q_{j,N}\rangle_{L^2}=o(1)\).
Writing \(A'(0)[v]\) for the derivative of \(A(p)\) in direction
\(v\), we have proved
\begin{equation}
 \langle Z_i,A'(0)[v_N]k_{j,N}\rangle_{L^2}
       \longrightarrow -\tfrac14D_{ij}.
 \label{msreg:ranklimit}
\end{equation}
The right-hand side is invertible. Only \(v_N\) is required
to be supported away from \(K\). The projected kernel fields
can have global support.

\subsection{Preserving the energy of the reference}

Assume \(E_+(r)>0\). Proposition~\ref{var:stationarity} applied
to the ASD cap supplies a smooth trace-free exterior direction
\(u_0\) compactly supported in \(M\setminus K\), with
\(dE_+(r)[u_0]\ne0\). Equivalently, by \eqref{var:bach},
a bump times the Bach tensor where it is nonzero gives
such a direction. From the fixed-volume isometric curvature
comparison we also have
\[
 dE_+(r)[v]=2\langle W,Lv\rangle_{L^2}
           =2\langle L^*W,v\rangle_{L^2}.
\]
In particular \(dE_+(r)[v_N]=O(N^{-2})\). Change \(v_N\) to
\[
 v_N^\parallel=v_N-\frac{dE_+(r)[v_N]}{dE_+(r)[u_0]}\,u_0 .
\]
The new direction is tangent to the energy level. Linearity
of \(A'(0)\) changes each entry in \eqref{msreg:ranklimit}
by at most \(O(N^{-2})\), because \((A'(0)[u_0])^*Z_i\)
is fixed and smooth and \(\|k_{j,N}\|_{L^2}=O(1)\).
Thus the new matrix still has the invertible limit
\(-D/4\). Choose one sufficiently large finite \(N\), put
\(v=v_N^\parallel\), and retain its \(k_j=k_{j,N}\).

Apply the scalar implicit function theorem to
\[
 (a,d)\longmapsto
 E_+\bigl(r(e^{(av+d u_0)^\sharp}\cdot,\cdot)\bigr)-E_+(r).
\]
The derivative in \(d\) at the origin is nonzero and the
derivative in \(a\) is zero. This gives a smooth \(\beta(a)=O(a^2)\)
such that, for \(p(a)=av+\beta(a)u_0\),
\begin{equation}
 E_+(r_{p(a)})=E_+(r),\qquad dV_{r_{p(a)}}=\nu .
 \label{msreg:level}
\end{equation}
The support of the reference change stays in a fixed compact
subset of the open exterior. Thus \(r_{p(a)}=r\) on a neighborhood
of \(K\), also when the round agreement of \(r\) itself is given
only on \(K\). Its logarithmic tangent is \(v\), so the mixed
matrix for \(A_a=A(p(a))\) is precisely the adjusted matrix
above. Equations \eqref{msreg:columns}--\eqref{msreg:inverse}
give regularity for all sufficiently small \(a\ne0\).

When the reference is regular already, take it without an
exterior change. If \(E_+(r)=0\), Proposition~\ref{prop:recognition}
gives conformal equivalence to the round sphere. In the canonical
identification of chiral two-forms under conformal rescaling,
\(F_{\lambda r}(\lambda h)=\lambda^{-1}F_r(h)\) for a positive
smooth function \(\lambda\), by the Weyl endomorphism weight
and \eqref{lift:chart}~\cite[discussion after (6.7)]{FeffermanGraham2008Ambient}.
Diffeomorphisms also transport the derivative, and the round
derivative is onto. Thus the zero-energy reference is regular,
and its radial prescription is constant. We have proved that
regular references are dense in the smooth fixed-density marked
class even on each Weyl energy level. They are relatively open:
a bounded right inverse persists under small operator-norm
changes, and \eqref{msreg:comparison} yields those changes from
\(C^{s+2}\) closeness of smooth references.

\subsection{Exact prescription and radial step size}

For each sufficiently small \(a\ne0\) on the path
\eqref{msreg:level}, abbreviate \(r_a=r_{p(a)}\).
The closed slice
\(\mathcal S=L^*U\oplus\mathcal J\mathbb R^m\) is the image of
\(Q_a\). The sum is direct by \(LL^*|_U\) injective and the
invertible matrix \eqref{msreg:mixedmatrix}. The operator \(Q_a A_a\)
is a bounded projection onto it. The inverse function theorem
on this slice gives exact realization of every sufficiently
small \(H^s\) target increment in the new chart.

The radial right-hand side has a more precise estimate.
Put \((u_a,c_a)=T_a^{-1}W_a\). This pair is smooth in \(a\)
through zero, and \(c_0=\mathsf M^{-1}b(W)\). If \(b(W)=0\),
then \(c_a/a\) extends smoothly and \(Q_aW_a\) is bounded in \(X\).
For \(E_+(r)>0\) its norm is bounded away from zero since
\(A_aQ_aW_a=W_a\to W\ne0\) and \(\|A_a\|\) stays bounded.
If \(b(W)\ne0\), then
\begin{equation}
 aQ_aW_a\longrightarrow \mathcal J\mathsf M^{-1}b(W)\ne0
                         \quad\text{in }X .
 \label{msreg:radialpole}
\end{equation}
In this case every radial preimage has norm at least a
positive multiple of \(|a|^{-1}\), by \eqref{msreg:lower}
with \(f=W_a\).

Hereafter suppose \(m>0\) and \(E_+(r)>0\). Along the fixed smooth
reference path the two-jet map has a common local Hessian bound \(H\)
in \(h\). Let
\(N_a(h)=\widehat F(p(a),h)-W_a-A_ah\).
Choose \(\rho_a=\gamma|a|\) with \(\gamma>0\) sufficiently small
that \(\|Q_a\|H\rho_a\le1/2\). On the closed ball of radius
\(\rho_a\) in \(\mathcal S\), the map
\(h\mapsto Q_af-Q_aN_a(h)\) is a contraction into the ball
when \(\|Q_af\|\le\rho_a/2\). Indeed the remainder is quadratic,
with \(\|N_a(h)\|_Y\le H\|h\|_X^2/2\), and it is
\(H\rho_a\)-Lipschitz there. This gives a sufficient target
radius of order \(a^2\) in arbitrary directions. For the radial
increment \(f=-tW_a\) it gives a solution on the whole interval
\begin{equation}
 \widehat F(p(a),h_a(t))=(1-t)W_a,\qquad
 h_a(0)=0,\qquad
 0\le t\le
 \begin{cases}
 C_2|a|,& b(W)=0,\\
 C_2 a^2,& b(W)\ne0 ,
 \end{cases}
 \label{msreg:radialstep}
\end{equation}
for some \(C_2>0\) independent of small \(a\) on the chosen path.
This is the contraction of Proposition~\ref{lift:radial}
with the actual directional norm \(\|Q_aW_a\|\).

Smooth sources yield spatially smooth solutions at the same
low-norm threshold. Write \(h=L^*u+\mathcal J d\) and adjoin
\(b(u)=0\) to the curvature equation on
\(H^{s+4}(B)\times\mathbb R^m\). The derivative is invertible on
\(U\times\mathbb R^m\) by the contraction bound, and the moment
row is an isomorphism on the added \(\mathcal C\) summand, which
\(L^*\) annihilates. The local translation proof following
Theorem~\ref{lift:exterior} applies to this fourth-order potential
equation, with \(\mathcal J d\) a smooth finite-dimensional
term. For the radial source it also gives smooth dependence
jointly on space and \(t\) for each fixed \(a\ne0\). At an
already regular initial reference the same argument works
with the fixed normal inverse and a fixed positive radius
and duration.

The powers in \eqref{msreg:radialstep} are optimal for solutions
with \(\|h\|_X\le C_3|a|\), for any fixed \(C_3\).
If \(b(W)\ne0\), applying \(b\) to the equation and using
\(bL=0\) gives
\[
 t|b(W_a)|\le C_4\bigl(|a|\|h\|_X+\|h\|_X^2\bigr)=O(a^2),
\]
and \(|b(W_a)|\) is bounded below. If \(b(W)=0\), use the
unprojected bound \(t\|W_a\|_Y\le C_4\|h\|_X=O(|a|)\)
and \(W\ne0\). These necessary bounds hold in the full chart
on the indicated metric scale, regardless of the choice of slice.

Let \(g_{a,t}=g(p(a),h_a(t))\). The identity in
\eqref{msreg:radialstep}, isometric target comparison and fixed
volume give
\[
 E_+(g_{a,t})=(1-t)^2 E_+(r_a)=(1-t)^2 E_+(r).
\]
The prescribed Weyl block vanishes on \(K\), so \(W^+_{g_{a,t}}=0\)
there. The metric itself can deform on \(K\). If \(r\) was
round on a neighborhood of \(K\), the block vanishes on a
possibly smaller common neighborhood. A positive initial
Yamabe sign persists for small \(|a|\) on the constructed
interval because the metrics are uniformly \(C^2\)-close to \(r\).
Here the sign equals that of the first conformal-Laplacian
eigenvalue, which is continuous by comparison of the Rayleigh
forms~\cite[(2.6)--(2.7)]{LeeParker1987Yamabe}.
Recognition by Proposition~\ref{prop:recognition} needs
positive Yamabe sign and a supplied time with
\((1-t)^2E_+(r)<4\pi^2\), or a supplied zero-energy metric.

Finally, the prescribed tensor in \eqref{msreg:radialstep}
uses the new reference. In the old chart let
\[
 \ell^\sharp=\log(e^{p^\sharp/2}e^{h^\sharp}e^{p^\sharp/2}),\qquad
 O=e^{\ell^\sharp/2}e^{-p^\sharp/2}e^{-h^\sharp/2}.
\]
The orthogonal map \(O\) is the ratio of the two successive
tangent isometries to the direct symmetric isometry from \(r\)
to the solved metric. For \(R=\Lambda^2O^*|_E\) this gives
\(\widehat F(p,h)=R F_r(\ell)R^{-1}\). Thus the old fixed
target would require the moving target
\((1-t)RWR^{-1}\), which depends on the unknown \(h\).
Indeed, expanding the exponentials at \(p=av+O(a^2)\) and
\(h=\epsilon k\) gives mixed coefficient
\([v^\sharp,k^\sharp]/8\) in \(O\). At \(p=0\) a differentiable
lift through the original
metric still requires \(b(W)=0\), by
\(Lh'(0)=-W\). The estimates above are local and their constants
depend on the chosen exterior variation. Continuing the
contraction would require control of further reference
changes and of the geometry of the deformed cap.

\section{Quadratic balance at the original reference}\label{mscone:section}

The deformation of Section~\ref{msreg:section} starts at a nearby
reference. Here the prescribed Weyl field and the center of the chart
are both those of the original metric \(r\).
Assume that \(r\) is round on an oriented round ball \(O\), with a
closed concentric ball \(K\Subset O\). Use the chart
\eqref{lift:chart}, the volume form \(\nu=dV_r\), and the real
\(L^2(r)\) pairings. Write \(W=F_r(0)\), \(L=dF_r(0)\),
\(H=d^2F_r(0)\), \(\mathcal C=\ker L^*\), and \(Q=L^*G_P\) for the
generalized inverse from \eqref{lift:normal}--\eqref{lift:split}.
We identify the moments \(b(f)\) of Theorem~\ref{lift:exterior}
with the functional \(Z\mapsto\langle Z,f\rangle_{L^2(r)}\)
on \(\mathcal C\).

The sign \(\sigma=1\) gives contraction and \(\sigma=-1\) expansion
in \(F_r(h_\tau)=(1-\sigma\tau^2)W\). In the chart domain
\(H^{s+2}(V)\), \(s\ge5\), a branch with \(h_\tau/\tau\to k\)
necessarily obeys
\begin{equation}\label{mscone:balance}
 Lk=0,\qquad q(k):=\tfrac12 b(H[k,k])=-\sigma b(W).
\end{equation}
Indeed, divide the equation by \(\tau\) for the first condition,
then take moments and divide by \(\tau^2\) for the second.
If \(b(W)\ne0\), the paired Taylor estimate
\(\|b(F_r(h)-W-Lh)\|\le C\|h\|_{H^{s+2}}^2\) gives
\(\|h\|_{H^{s+2}}\ge c\sqrt{|t|}\) for all sufficiently small
solutions with \(F_r(h)=(1-t)W\). A differentiable start in \(t\)
would require \(b(W)=0\).

\subsection{Round cap directions and the fixed-chart Hessian}

\begin{proposition}[Kernel directions supported in the round cap]
\label{mscone:cap}
Let \(U\Subset O\) be a concentric round ball. The smooth kernel
tensors compactly supported in \(U\) are precisely
\(\mathcal K C_c^\infty(U,TM)\), where
\(\mathcal K X=(\mathcal L_Xr)_0\). Each of them satisfies
\[
 b(H[\mathcal K X,j])=0
 \qquad\text{for every smooth }j\in\ker L.
\]
For smooth kernel tensors supported in a closed concentric
subball of \(O\), the field \(X\) can be taken supported in that
same closed subball.
\end{proposition}

To prove the support assertion, transplant the tensor to the
standard round sphere and extend it by zero. The extension
\(\widehat k\) satisfies \(D_0\widehat k=0\) by locality. Solve
\(\mathcal K_0^*\mathcal K_0 Y=\mathcal K_0^*\widehat k\)
orthogonally to \(\ker\mathcal K_0\). The right side annihilates
this kernel, and \(\mathcal K_0^*\mathcal K_0\) is elliptic by
injectivity of the conformal Killing symbol. The remainder
\(\widehat k-\mathcal K_0Y\) is TT. Since \(D_0\mathcal K_0=0\),
the round TT factorization in Section~\ref{sec:lift} kills the
remainder, giving \(\widehat k=\mathcal K_0Y\)
~\cite[\S3]{York1974Covariant}
\cite[Theorem 3.10 and \S7.1]{GurskyViaclovsky2015}.

On the connected round annulus outside the support, \(Y\) is
conformal Killing. In oriented stereographic coordinates centered
on the ball, the same vector field is conformal Killing for the
Euclidean metric. If \(f=\operatorname{div}Y/4\) in these coordinates,
differentiating the conformal Killing equation gives
\[
 \partial_i\partial_jY_k
 =\delta_{jk}\partial_i f+\delta_{ik}\partial_j f-\delta_{ij}\partial_k f.
\]
Thus \(\Delta Y=-2\,df\), \(\Delta f=0\), and applying \(\Delta\)
to the first-order equation yields \(\partial_i\partial_j f=0\).
Integration on the connected annulus gives
\(Y=a+A_0x+\lambda x+2(d_0\cdot x)x-|x|^2d_0\), with \(A_0\)
skew and all coefficients constant. This polynomial extends as a
conformal Killing field across the whole ball. Subtract the extension
and continue by zero to obtain \(X\) with the required support.
Conversely, \eqref{lift:gauge} places such \(\mathcal K X\) in
the kernel, since \(W=0\) throughout the round patch.

For the Hessian calculation take any smooth \(X\) and write
\(\varphi_a\) for its flow and \(J_a=\varphi_a^*\nu/\nu\).
Define the chart action \(\Psi_a\) by
\(g_{\Psi_a(h)}=J_a^{-1/2}\varphi_a^*g_h\).
This metric has volume form \(\nu\). At \(y=\varphi_a(x)\), set
\[
 R_h=e^{h^\sharp/2},\qquad B_{a,h}=J_a^{-1/4}R_h(y)d\varphi_a(x),
 \qquad O_{a,h}=B_{a,h}(B_{a,h}^*B_{a,h})^{-1/2}.
\]
The tangent adjoint is taken with respect to the reference
metrics at \(x,y\). We have \(B_{a,h}^*B_{a,h}=e^{\Psi_a(h)^\sharp}\),
so \(O_{a,h}\) compares the symmetric frames by an oriented
\(r\)-isometry. The conformal weight of the Weyl endomorphism
is the inverse metric scale, and the middle-degree Hodge star
is unchanged by conformal rescaling. Thus
\begin{equation}\label{mscone:naturality}
 F_r(\Psi_a(h))(x)=J_a(x)^{1/2}\,
 O_{a,h}^* F_r(h)(y)(O_{a,h}^*)^{-1},
\end{equation}
where \(O_{a,h}^*\) on the right denotes pullback on self-dual forms
~\cite[discussion after (6.7)]{FeffermanGraham2008Ambient}.

Compare tangent spaces by parallel transport along the flow.
Put \(D=\nabla X\), with \(Dv=\nabla_vX\),
\(S=\operatorname{sym}D\), \(\Omega_0=\operatorname{skew}D\),
and \(\theta=\operatorname{div}_rX\).
If \(\Omega(h)=\partial_a O_{a,h}|_{a=0}\), differentiating the
polar product and subtracting its adjoint gives
\[
 \Omega(h)R_h+R_h\Omega(h)=R_hD-D^*R_h .
\]
Here \(\partial_a B_{a,h}|_0=\nabla_XR_h+R_hD-\theta R_h/4\);
the first and last terms are self-adjoint. Differentiating
the displayed equation in a metric direction \(j\) at zero gives
\begin{equation}\label{mscone:polar}
 \begin{aligned}
 2(d\Omega)_0[j]+\tfrac12(\Omega_0j^\sharp+j^\sharp\Omega_0)
   &=\tfrac12(j^\sharp D-D^*j^\sharp),\\
 (d\Omega)_0[j]&=\tfrac14[j^\sharp,S]
                  =\tfrac18[j^\sharp,e^\sharp],
                 \qquad e=\mathcal K X.
 \end{aligned}
\end{equation}
In the last equality \(e^\sharp=2S-\theta I/2\).
Denote by \(\iota\) the infinitesimal pullback on \(\Lambda^+\)
induced by a skew tangent operator, obtained by inserting that
operator in each argument of a two-form. On target sections define
\[
 T_X f=\nabla_Xf+\tfrac12\theta f+[\iota(\Omega_0),f],
 \qquad C_X(j)=[\iota([j^\sharp,e^\sharp]/8),W].
\]
Differentiating \eqref{mscone:naturality} first in \(a\) and then in
\(h\) gives \(Le=T_XW\) and the exact identity
\begin{equation}\label{mscone:mixed}
 H[e,j]+L\mathcal S_X j=T_XLj+C_X(j),
 \qquad \mathcal S_X j=d_h\big(\partial_a\Psi_a(h)|_{a=0}\big)_0[j].
\end{equation}
The mixed derivative of the logarithm in \(\Psi_a\) subtracts
\((j^\sharp e^\sharp+e^\sharp j^\sharp)/2\) from the mixed metric
derivative, giving
\[
 (\mathcal S_X j)^\sharp=\nabla_Xj^\sharp+[j^\sharp,\Omega_0].
\]
Pairing \eqref{mscone:mixed} with \(Z\in\mathcal C\) kills the
\(L\mathcal S_Xj\) term. For \(X\) supported in
the round patch, \(C_X(j)=0\) everywhere, and for \(Lj=0\)
the right side of \eqref{mscone:mixed} vanishes. This proves
the mixed assertion of Proposition~\ref{mscone:cap}.

The flow of such an \(X\) preserves the cap region where it
acts and is the identity off that region. Equation
\eqref{mscone:naturality} therefore preserves the entire prescribed
field \((1-t)W\) at any exact lift in the chart, since that field is
zero where the flow acts. For the quadratic balance, adding a compact
cap kernel direction changes neither \(q(k)\) nor its kernel
polarization. In particular the projected Hessian on the smooth kernel
descends to its exterior restriction outside a closed concentric
subball of \(O\). A family consisting only of such cap directions
cannot solve \eqref{mscone:balance} when \(b(W)\ne0\).

\subsection{Exterior polarizations}

\begin{proposition}[Oscillatory kernel directions]\label{mscone:probe}
In a relatively compact patch of \(M\setminus K\), let
\(n=\nabla_r\phi\) with \(|n|_r=1\), let \(\beta\) be a real smooth
compactly supported cutoff, and let \(u\) be a smooth real tangent
field on the patch. Put
\[
 X_N=N^{-1}\beta\sin(N\phi)u,\qquad e_N=\mathcal K X_N,\qquad
 p_N=QLe_N,\qquad k_N=e_N-p_N ,
\]
extending \(X_N\) by zero. Then \(k_N\in\ker L\) is smooth and
global. For every \(Z\in\mathcal C\), with
\(d=\langle u,n\rangle_r\) and
\(A=\iota(\operatorname{skew}(u\otimes n^\flat))\), one has
\begin{equation}\label{mscone:probeformula}
 \lim_{N\to\infty} b(H[k_N,k_N])(Z)
 =\tfrac12\int_M\beta^2\!
       \left(\tfrac{d^2}{4}\langle Z,W\rangle
             +\langle[A,Z],[A,W]\rangle\right)\nu .
\end{equation}
For a finite family with pairwise disjoint compact supports of the
cutoffs, the off-diagonal projected Hessians tend to zero as the
frequencies tend independently to infinity. The longitudinal
choice \(u=n\) gives \(+\tfrac18\int_M\beta^2\langle Z,W\rangle\nu\)
for the unhalved Hessian.
\end{proposition}

\begin{proof}
For any smooth \(X\), let \(e=\mathcal K X\), \(p=QLe\) and
\(k=e-p\). Since \(Le\perp\mathcal C\), \(Lp=Le\) and \(Lk=0\).
Elliptic regularity gives smoothness. The divergence theorem and
skew-adjointness of commutation by \(\iota(\Omega_0)\) in
Hilbert--Schmidt norm give \(T_X^*=-T_X\). Also \(C_X(e)=0\)
exactly. Consequently \eqref{mscone:mixed} yields
\begin{equation}\label{mscone:projected}
 b(H[k,k])(Z)
 =\langle T_XZ,T_XW\rangle_{L^2(r)}
   -2\langle Z,C_X(p)\rangle_{L^2(r)}+b(H[p,p])(Z).
\end{equation}
The positive first term comes from the uncorrected pure pairing
\(\langle Z,T_X^2W\rangle=-\langle T_XZ,T_XW\rangle\)
minus twice its copy in the mixed pairing.

For \(X=X_N\) we have \(X_N=O(N^{-1})\) and
\(\nabla X_N=\beta\cos(N\phi)u\otimes n^\flat+O(N^{-1})\)
uniformly. Hence, for any fixed smooth target \(f\),
\[
 T_{X_N}f=\beta\cos(N\phi)\big(\tfrac d2 f+[A,f]\big)+O_f(N^{-1}).
\]
In particular \(Le_N\) is bounded in \(L^2\) and is \(O(N^{-1})\)
in \(H^{-1}\). For the latter, integrate
\(\cos(N\phi)=N^{-1}n(\sin(N\phi))\) against an \(H^1\) test on
the compact cutoff support. The elliptic normal inverse gives
\(Q:H^\ell(B)\to H^{\ell+2}(V)\) for \(\ell=0,-1\).
Thus \(p_N\) is bounded in \(H^2\) and is \(O(N^{-1})\) in \(H^1\)
~\cite[(1)--(2)]{Melrose2017Lecture22}. The rotation term in
\eqref{mscone:projected} is bounded by
\(C_Z\|\nabla X_N\|_\infty\|p_N\|_2\). The paired Hessian is
bounded on \(H^1\times H^1\): its second-order terms have the
forms \(p\nabla^2j\), \(j\nabla^2p\), and \(\nabla p\nabla j\)
with smooth coefficients, since curvature is affine in second
metric derivatives. One integration by parts against \(Z\) gives
the bound. Thus both errors in \eqref{mscone:projected} tend to zero.

The cross terms between \(d/2\) and \([A,\cdot]\) cancel by
skew-adjointness. Integrating \(\cos^2(N\phi)=(1+\cos(2N\phi))/2\)
then proves \eqref{mscone:probeformula}, with the oscillating term
vanishing by the same integration by parts. Polarizing
\eqref{mscone:projected} proves the disjoint statement. The leading
mixed terms are products of fields supported at the original
cutoffs, and the mixed corrections are estimated by the same
\(L^\infty\), \(L^2\), and \(H^1\) bounds. Because \(\mathcal C\)
is finite dimensional, these are limits in \(\mathcal C^*\).
In the longitudinal case \(\Omega_0=O(N^{-1})\), so also directly
the pure and mixed unhalved pairings each tend to
\(-\tfrac18\int_M\beta^2\langle Z,W\rangle\nu\), and the
correction-correction pairing tends to zero.
\end{proof}

At a point take the oriented orthonormal tangent frame
\((n,e_1,e_2,e_3)\) and the positive frame
\(\omega_i=(n^\flat\wedge e_i^\flat+*(n^\flat\wedge e_i^\flat))/\sqrt2\).
Write \(J_i\omega_j=\varepsilon_{ijk}\omega_k\) for the generators of
unit angular speed, with \(\|J_i\|_{\HS}=\sqrt2\).
For \(u=e_1\), the tangent skew part sends \(n\) to \(e_1/2\)
and \(e_1\) to \(-n/2\). Its pullback fixes \(\omega_1\)
infinitesimally and sends \(\omega_2\) to \(-\omega_3/2\) and
\(\omega_3\) to \(\omega_2/2\). Cyclically,
\[
 u=dn+\sum_i u_i e_i\quad\Longrightarrow\quad
 A=-\tfrac12\sum_i u_iJ_i .
\]
Thus \(d\) and \(A\) in \eqref{mscone:probeformula} are independent
pointwise, and every rotation generator on \(\Lambda_r^+\) occurs
at a fixed \(n\).

\subsection{The dual quadratic form and an exact signed lift}

For \(Z\in\mathcal C\), define a symmetric form \(D_Z(x)\) on
\(\mathfrak{so}(\Lambda_{r,x}^+)\) by its quadratic evaluation
\[
 D_Z(x)(A,A)=\langle[A,Z(x)],[A,W(x)]\rangle .
\]
In the \(J_i\) basis its matrix is the symmetrization of
\(\langle[J_i,Z],[J_j,W]\rangle\). The two alternating-symbol
contractions \(\sum_iJ_i^2=-2I\) and
\(\sum_iJ_iM J_i=M^{\mathsf T}-(\tr M)I\) give
\(-\sum_i[J_i,[J_i,M]]=6M\) on trace-free symmetric
endomorphisms. Thus
\begin{equation}\label{mscone:trace}
 \sum_iD_Z(x)(J_i,J_i)=6\langle Z(x),W(x)\rangle .
\end{equation}
In an oriented Weyl eigenframe \(W=\mathrm{diag}(w_1,w_2,w_3)\),
\(Z=(z_{ij})\), the diagonal entries are
\(2(w_j-w_k)(z_{jj}-z_{kk})\) for cyclic \((i,j,k)\).
The \(ij\)-entry for distinct \(i,j,k\) is \(3w_kz_{ij}\),
by trace-freeness and the displayed commutation rule for \(J_i\).

\begin{theorem}[A criterion at the fixed reference]\label{mscone:chord}
The positive hull in \(\mathcal C^*\) of the diagonal limits
in \eqref{mscone:probeformula} is the whole space exactly when
there is no \(0\ne Z\in\mathcal C\) such that
\begin{equation}\label{mscone:witness}
 D_Z(x)\succeq0\qquad\text{for every }x\in M\setminus K.
\end{equation}
Under this absence condition, for either fixed sign
\(\sigma\in\{1,-1\}\) there is an \(\epsilon>0\) and a family
\(h_\tau\), smooth in \((x,\tau)\) for \(|\tau|<\epsilon\),
with \(h_0=0\) and
\[
 F_r(h_\tau)=(1-\sigma\tau^2)W .
\]
In particular there is a one-sided contracting chord for
\(0\le t<\epsilon^2\). The volume form is \(\nu\) and the resulting
metrics are ASD on \(K\).
\end{theorem}

\begin{proof}
If \eqref{mscone:witness} holds, the trace identity
\eqref{mscone:trace} makes the longitudinal term nonnegative in
the \(Z\) row and the definition of \(D_Z\) does the same for the
transverse term. Conversely, if that row is negative for some
rotation at an exterior point, realize the rotation by a transverse
\(u\). A signed-distance phase exists locally with any prescribed
unit normal, and the transverse field extends smoothly in the patch.
A nonzero cutoff of sufficiently small support gives a strictly
negative moment in this row.

Suppose \(\mathcal C\ne0\) and there is no witness. For each unit
\(Z\) select a strict negative exterior point evaluation with a
transverse polarization. The corresponding open subsets cover the
unit sphere of \(\mathcal C\). A finite subcover gives a uniform
strict margin. Moving repeated centers slightly and replacing the
point evaluations by cutoffs with \(\int\beta_j^2\nu=1\) in small
disjoint patches preserves a strict margin. Testing also \(-Z\)
shows that the convex hull of the resulting limiting vectors contains
a ball about zero, by finite-dimensional convex separation.
They therefore positively span \(\mathcal C^*\). The finite
margin establishes the claim directly for the positive hull. In particular the
strict both-signs condition for \(x\mapsto\langle Z(x),W(x)\rangle\)
for each \(Z\ne0\) already suffices using longitudinal probes alone.

Let \(v_j\) be half the limiting Hessian vectors for the selected finite
family. For the desired target \(z=-\sigma b(W)\), express
\(z-\sum_jv_j\) as a nonnegative combination of the \(v_j\).
Then \(\sum_j (a_j^0)^2v_j=z\) for strictly positive amplitudes
\(a_j^0\). The derivative columns \(2a_j^0v_j\) span
\(\mathcal C^*\). Fix all but an independent set of columns.
By \eqref{mscone:probeformula} and the disjoint mixed limits,
the quadratic coefficients of
\[
 q_N(a)=\tfrac12 b\big(H[k_N(a),k_N(a)]\big),
 \qquad k_N(a)=\sum_j a_j k_{j,N},
\]
converge to those of \(\sum_j a_j^2v_j\). The implicit function
theorem with the coefficients as parameters gives a regular
solution \(a^N\to a^0\) of \(q_N(a^N)=z\). Fix one sufficiently
large finite \(N\) and put \(k(a)=k_N(a)\), \(k_*=k(a^N)\).
Regularity implies \(k_*\ne0\). The nonlinear step uses
this fixed frequency and a parameter interval that can depend on it.

Use \(\mathcal W\) and \(\Gamma=b\mathcal W\) from
Theorem~\ref{lift:exterior}. At the fixed frequency, solve
\begin{equation}\label{mscone:augmented}
 F_r\big(\tau k(a)+L^*v\big)-W+\mathcal Wc+\sigma\tau^2W=0,
 \qquad v\in H^{s+4}(B)\cap\mathcal C^\perp ,
\end{equation}
with \((\tau,a)\) as parameters. The derivative in \((v,c)\)
at \(\tau=0\) is the isomorphism
\((\dot v,\dot c)\mapsto P\dot v+\mathcal W\dot c\).
The implicit solution has \(v(0,a)=c(0,a)=0\) and zero first
\(\tau\)-derivatives there, since \(Lk(a)=0\). Put
\(h(\tau,a)=\tau k(a)+L^*v(\tau,a)\), and
\(\widehat h(\tau,a)=\int_0^1\partial_\tau h(\rho\tau,a)\,d\rho\).
Then \(h=\tau\widehat h\) with \(\widehat h(0,a)=k(a)\). Taking
moments in \eqref{mscone:augmented} and using integral Taylor
expansion gives the smooth divided equation
\[
 \begin{aligned}
 \Gamma c(\tau,a)&=-\tau^2\mathcal R_\sigma(\tau,a),\\
 \mathcal R_\sigma(\tau,a)
 &=\sigma b(W)+\int_0^1(1-\lambda)\,
    b\big(d^2F_r(\lambda h)[\widehat h,\widehat h]\big)\,d\lambda,\\
 \mathcal R_\sigma(0,a)&=\sigma b(W)+q_N(a).
 \end{aligned}
\]
The chosen finite derivative is invertible at \(a^N\), so a
second implicit function theorem gives \(a(\tau)\) with
\(\mathcal R_\sigma(\tau,a(\tau))=0\) and \(c=0\). This argument
also applies to any supplied finite-dimensional smooth kernel
family with a regular signed zero of \eqref{mscone:balance}.
When \(\mathcal C=0\), the inverse of \(P\) gives a branch
directly in \(t\), as in Theorem~\ref{lift:exterior}.

For spatial regularity in \eqref{mscone:augmented}, work on
the full potential space with the adjoined row \(b(v)=0\).
The derivative in the two unknowns remains invertible at
the small solutions. Conjugate the differential row by local
component pullbacks that are translations on smaller framed balls.
The translated coefficients and the translated smooth finite kernel
family are smooth in the translation parameter. In the moment row
a change of variables transfers the translation to the smooth
test sections and density. Strong continuity and uniqueness identify
the translated solution with the smooth implicit solution at the
same Sobolev grades. Translation derivatives on the smaller balls
give all spatial derivatives with smooth parameter dependence.
A finite cover gives joint smoothness on the same local interval.
The cap conclusion is the prescribed vanishing of the positive Weyl
block there.
\end{proof}

For the constructed branch with \(\mathcal C\ne0\),
\(h_\tau=\tau k_*+O(\tau^2)\) in each fixed finite spatial
norm. In \(t=\tau^2\) on the contracting side this branch is
\(1/2\)-H\"older at zero, smooth for \(t>0\), and has no right
derivative at zero. The metric on the cap can vary.

A witness \eqref{mscone:witness} that is nonzero somewhere as a
quadratic form has strictly positive \(\langle Z,W\rangle\) there
by \eqref{mscone:trace} and nonnegative pairing everywhere in the
exterior. Since \(W=0\) on \(K\), it strictly separates the
contracting target \(-b(W)\) from the limiting cone. For a fixed
disjoint configuration the balance then fails on bounded amplitude
sets at sufficiently high frequency. If every chosen diagonal
evaluation in the witness row is strictly positive, the finite
row remains positive definite and the exclusion holds for all
amplitudes of that configuration. A witness with \(D_Z\equiv0\) on
the exterior instead gives a zero limiting equation and zero
limiting amplitude derivative in that row. These restrictions
concern the cone of one-phase limits; they leave general kernel
directions and finite-frequency corrections open. In particular,
whether the absence condition holds at an arbitrary marked
reference is unresolved.

For a separate normalization check, the equal-curvature Einstein
product of unit round two-spheres has \(Z=W\in\ker L^*\)
by the first variation of \(E_+\) at a closed Einstein metric.
Oppositely rescaling the factors by \(e^\tau,e^{-\tau}\) gives the
fixed-volume field \(\cosh(\tau)W\), so its second derivative points
toward increasing Weyl curvature. Its \(D_W\) is nonnegative, as
\(\langle[A,W],[A,W]\rangle\) shows directly. This example checks
the signs of the moment and the radial target on a manifold with
\(b_2=2\).

\section{Transmission across a round spherical interface}\label{sec:transmission}

Let \(r\) be the unit round metric on the oriented standard \(S^4\).
Use oriented stereographic coordinates \(x\) from the south pole,
\(\rho=|x|\), and \(r=\Omega^2\delta\), where
\(\Omega=2/(1+\rho^2)\) and \(\delta\) is Euclidean. The northern
hemisphere is \(B=\{\rho<1\}\). Denote the closed southern hemisphere,
including the removed pole, by \(S_-\). We use the curvature-operator
Hilbert--Schmidt norms and the convention that \(W^+=0\) is ASD.

At a reference metric \(g\) use trace-free covariant logarithmic
variations \(h\), with \(h^\sharp=g^{-1}h\), and write
\[
 g_h=g(e^{h^\sharp}\,\cdot,\cdot),\qquad
 F_{g,\pm}(h)=I_{h,\pm}W_{g_h}^{\pm}I_{h,\pm}^{-1}.
\]
Here \(I_{h,\pm}\) is the pullback on chiral two-forms induced by
the positive symmetric isometry \(e^{-h^\sharp/2}\) from \(g\) to
\(g_h\). These identifications preserve the chiral norms, and
\(dV_{g_h}=dV_g\). Write \(D_{\pm,g}=dF_{g,\pm}(0)\) for \(g=r\)
or \(\delta\). At these conformally flat backgrounds the derivative is
independent of the infinitesimal target frame change.

\subsection{The relative chiral operator}

In an oriented flat orthonormal coframe \(e^0,e^1,e^2,e^3\), put
\(\omega_i^\pm=(e^0\wedge e^i\pm e^j\wedge e^k)/\sqrt2\) for cyclic
\((i,j,k)\) in \((1,2,3)\). The round coframes on chiral two-forms
are \(\Omega^2\omega_i^\pm\). For a unit covector \(n\), let
\(R_n=I-2n\otimes n^\sharp\) and
\(C_n=\Lambda^2R_n|_{\Lambda_\delta^+}:\Lambda_\delta^+\to\Lambda_\delta^-\).
Orientation reversal exchanges the two Hodge eigenspaces. Define the
isometric conjugation \(M(n)V=C_nVC_n^{-1}\) on the Weyl targets.

\begin{proposition}\label{xfer:multiplier}
Use the Fourier transform with phase \(e^{-ix\cdot\xi}\) and inverse
measure \((2\pi)^{-4}d\xi\). Let \(T\) be the flat multiplier with
symbol \(M(\xi/|\xi|)\). If
\(Q_r=(D_{+,r}|_{\mathrm{TT}_r})^{-1}\) is the round inverse, then
\[
 D_{-,r}Q_r=\mathcal C_-^{-1}T\mathcal C_+,\qquad
 \mathcal C_\pm f=\Omega^2 f
\]
in the induced orthonormal frames. Each \(\mathcal C_\pm\) is an
onto \(L^2\) isometry, and so is the relative round operator.
\end{proposition}

\begin{proof}
The inverse \(Q_r\) from Section~\ref{sec:lift} uses round ASD
unobstructedness~\cite[\S1.1]{GoverGursky2024} and the round TT
kernel statement. For the flat symbols first substitute real
derivatives \(\partial\mapsto\xi=\ell e^0\). The TT symbol
fiber consists of the trace-free spatial components \(h_{ij}\).
The linearized curvature operator has entries
\(\mathcal R'_{0i,0j}=-\ell^2h_{ij}/2\), with entries involving
purely spatial pairs zero. Its two diagonal chiral compressions
are therefore both \(-\ell^2(h_{ij})/4\) in the respective
\(\omega_i^\pm\) frames. Their traces vanish, so the scalar
subtraction in \(W^\pm\) does not affect these matrices.
Fourier differentiation changes both signs together.
The symbol of the trace-free Lie derivative spans the
orthogonal complement of the TT symbol fiber in trace-free
tensors. Both chiral symbols vanish there by diffeomorphism and
conformal invariance, and both vanish on pure trace variations.
At \(\xi=\ell e^0\), \(C_n\omega_i^+=-\omega_i^-\), and the two minus
signs cancel in conjugation. Rotation equivariance consequently gives,
for the full Fourier symbols \(d_\pm\) and \(\xi\ne0\),
\begin{equation}\label{xfer:symbol}
 \begin{gathered}
 d_-=M d_+,\qquad d_-^*M=d_+^*,\\
 d_-^*d_-=d_+^*d_+,\qquad d_\pm d_\pm^*=\tfrac{|\xi|^4}{16}I .
 \end{gathered}
\end{equation}

The Hodge star on two-forms is conformally invariant in four
dimensions. The covariant Weyl tensor scales with weight \(2\),
so its endomorphism in the induced orthonormal frame has
weight \(-2\)~\cite[Ch.~6]{FeffermanGraham2008Ambient}. For a
round covariant variation \(h\) this gives
\begin{equation}\label{xfer:weights}
 u=\Omega^{-2}h,\qquad D_{\pm,\delta}u=\Omega^2D_{\pm,r}h,
 \qquad dV_r=\Omega^4\,dx .
\end{equation}
No Euclidean divergence condition on \(u\) is required. For smooth
round \(h=Q_r f\), \(u\) is bounded in the flat frame, and the two
fields in \eqref{xfer:weights} are in \(L^2(\mathbb R^4)\).
Apply \eqref{xfer:symbol} to their Fourier transforms away from
zero. The difference is an \(L^2\) field supported at a single
frequency and therefore vanishes. This determines the relative
operator on smooth round fields. The volume transformation gives
the isometries \(\mathcal C_\pm\), and Plancherel and the
pointwise orthogonality of \(M\) give the asserted extension.
\end{proof}

\subsection{Angular harmonicity, radial moments, and the pole}

\begin{lemma}\label{xfer:harmonic}
For any constant self-dual Weyl endomorphism \(V\),
\[
 Y_V(n)=M(n)V,\qquad H_V(x)=\rho^4Y_V(x/\rho)
\]
extend in the second formula to a componentwise homogeneous
harmonic quartic on \(\mathbb R^4\).
\end{lemma}

\begin{proof}
Identify the oriented Euclidean space with the quaternions in
the frame \((1,\mathbf i,\mathbf j,\mathbf k)\).
Left and right multiplication by an imaginary quaternion correspond
to the indicated positive and negative two-form bases with the same
normalization. For example, multiplication by \(\mathbf i\) takes
\(1\) to \(\mathbf i\) in both cases and takes \(\mathbf j\)
to \(\mathbf k\) on the left and \(-\mathbf k\) on the right.
Reflection acts by \(R_n(q)=-n\bar qn\), so conjugating
left multiplication by \(v\) with \(R_n\) gives right multiplication
by \(-\bar nvn\). Thus \(C_n=-O(n)\) in the indexed bases, where
\(O(n)v=\bar nvn\), and \(M(n)V=O(n)VO(n)^{\mathsf T}\).
Writing \(x=x_0+b\) with \(b\) imaginary, define the quadratic
matrix \(L(x)v=\bar xvx\). With \([b]_\times v=b\times v\),
\[
 L(x)=(x_0^2-|b|^2)I+2bb^{\mathsf T}-2x_0[b]_\times.
\]
It satisfies \(\Delta L=0\) for \(\Delta=\sum_j\partial_j^2\),
and \(H_V=L V L^{\mathsf T}\).
At \(x=1\) the first derivatives of \(L\) are \(2I,-2J_i\),
where \(J_i v=e_i\times v\). Contracting the alternating symbols
gives \(\sum_i J_i VJ_i^{\mathsf T}=(\operatorname{tr}V)I-V^{\mathsf T}\).
Thus \(\Delta H_V(1)=8(V+(\operatorname{tr}V)I-V^{\mathsf T})=0\).
For unit \(a\), \(L(xa)=O(a)L(x)\), and right multiplication by
\(a\) is orthogonal. The Laplacian has the same equivariance, so
it vanishes on the unit sphere and hence by homogeneity everywhere.
\end{proof}

Write \(H_4(\xi)V=H_V(\xi)\) for the polynomial numerator of \(M\).
The fundamental solution of \(\Delta^2\) in these conventions is
\(G(x)=-\log\rho/(8\pi^2)\). Indeed
\(\Delta\log\rho=2\rho^{-2}\) distributionally with zero atom,
and the flux of \(\nabla\rho^{-2}\) is \(-2|S^3|=-4\pi^2\),
giving \(\Delta^2G=\delta_0\). The order-four differentiation
has phase factor \(i^4=1\). Thus the Fourier transform of
\(H_4(\partial)G\) is \(M\) away from zero. Exact homogeneity of
\(H_4(\partial)G\) of degree \(-4\) gives homogeneity of degree
zero for its transform, excluding a discrepancy supported at the
origin~\cite[Theorem~4.19]{Dyatlov2022Elliptic}.
For a scalar function \(\Phi\) of \(\rho^2\), the trace-free
quartic coefficient tensor gives
\begin{equation}\label{xfer:radial-derivative}
 H_V(\partial)\Phi(\rho^2)=16H_V(x)\Phi^{(4)}(\rho^2).
\end{equation}
Indeed all second derivatives of the radial argument pair two
indices of the coefficient tensor and drop out. The derivatives
on \(\log\rho\) and \(\rho^{-2}\) are therefore
\(-48H_V(x)\rho^{-8}\) and \(384H_V(x)\rho^{-10}\).
The physical kernel acting on \(V\) is the spherical principal
value of \(6H_V(x)/(\pi^2\rho^8)\). The angular mean vanishes
by harmonicity. The possible point mass in a homogeneous
extension of degree \(-4\) is absent, since both this principal
value and \(H_V(\partial)G\) pair to zero with radial tests.

For a smooth signed radial \(\chi\) supported in \(\{\rho\le R\}\),
write \(m_j=\int_{\mathbb R^4}|x|^j\chi(|x|)\,dx\) for \(j=0,2\).
For \(s<\rho\), the normalized sphere average in \(n\) of
\(\log|x-sn|\) is \(\log\rho+s^2/(4\rho^2)\). To see this,
its radial Laplacian in \(s\) is the mean of
\(2|x-sn|^{-2}\), which is \(2/\rho^2\) by harmonicity, and
the initial value and derivative are \(\log\rho\) and zero.
Integration and \eqref{xfer:radial-derivative} then yield
\begin{equation}\label{xfer:radial}
 \begin{gathered}
 G*\chi=-\frac1{8\pi^2}\left(m_0\log\rho+\frac{m_2}{4\rho^2}\right),
 \qquad T(\chi V)=a(\rho)Y_V(n),\\
 a(\rho)=\frac6{\pi^2}\left(m_0\rho^{-4}-2m_2\rho^{-6}\right)
 \qquad(\rho>R).
 \end{gathered}
\end{equation}
Both identities extend continuously to \(\rho=R>0\) for the
smooth source. For \(V\ne0\) the exterior field vanishes
identically exactly when both moments vanish.

From now on take \(f=\chi Z\) with \(Z=\operatorname{diag}(1,1,-2)\)
and the positive seed
\[
 \chi(\rho)=
 \begin{cases}
 \exp\bigl(-1/(1-4\rho^2)\bigr),&\rho<1/2,\\
 0,&\rho\ge1/2 .
 \end{cases}
 \qquad f_r=\Omega^{-2}f .
\]
The source \(f_r\) extends by zero at the south pole.
Here \(a(\rho)>0\) for \(\rho\ge1/2\). Indeed, with \(v=4s^2\),
\(m_2/m_0\) is one quarter of the mean of \(v\) for density
\(v e^{-1/(1-v)}\,dv\) on \((0,1)\). For \(v>1/2\) the logarithm
of the ratio of this density to its value at \(1-v\) is
\(\log b-b+1/b<0\), with \(b=v/(1-v)>1\). The mean is therefore
less than \(1/2\), so \(2m_2/m_0<1/4\).

The first-order round field \(U=D_{-,r}Q_r f_r\) has components
\(\Omega^{-2}a(\rho)Y_Z(x/\rho)\) for \(\rho\ge1/2\). In the
opposite oriented pole chart \(y=Jx/\rho^2\), \(J=R_{e^0}\), the
orthonormal frame transition has orthogonal factor \(JR_n\),
since \(dy=\rho^{-2}JR_n\,dx\) and \(\Omega(y)=\rho^2\Omega(x)\).
The composition of \(\Lambda^2(JR_n)\) on the negative frame
with \(C_n\) equals \(\Lambda^2J|_{\Lambda_\delta^+}\), represented
by \(-I\). Thus the angular conjugation cancels, and the
components of \(U\) in the \(y\) frame are
\begin{equation}\label{xfer:pole}
 \frac{3}{2\pi^2}(1+w)^2(m_0-2m_2w)Z,\qquad w=|y|^2 .
\end{equation}
In particular the equatorial and southern pole norms are
\(6\sqrt6(m_0-2m_2)/\pi^2\) and \(3\sqrt6\,m_0/(2\pi^2)\),
respectively. For \(0\le w\le1\) the scalar derivative in
\eqref{xfer:pole} is proportional to
\((1+w)(m_0-m_2-3m_2w)>0\), since \(m_2<m_0/4\).
Hence the pole norm is the positive minimum of \(|U|_r\) on \(S_-\).

\subsection{The sharp observation inside a supplied collar}

Let \(A=\{2/3<\rho<1\}\) and let \(\mathcal Y\) be the
five-dimensional angular subspace
\(\{Y_V:V\in S^2_0(\Lambda_\delta^+)\}\) in
\(L^2(S^3,dS)\) of negative Weyl endomorphisms. The angular
projection \(P_{\mathcal Y}\) is applied separately on every radius.
The isometry \(M(n)\) gives
\(\langle Y_V,Y_{V'}\rangle_{L^2(S^3)}=2\pi^2\langle V,V'\rangle\).

\begin{theorem}[Five angular projections independent of exterior data]\label{xfer:sharp}
On \(B\) suppose the same smooth flat metric variation \(h\)
satisfies \(D_{+,\delta}h=tf+e\) and \(q=D_{-,\delta}h\), with
\(e\in L^2(B)\). Let \(\widetilde e\) denote its zero extension.
On \(A\) one has
\begin{equation}\label{xfer:angular}
 P_{\mathcal Y}(q-T\widetilde e)=t a(\rho)Y_Z,\qquad
 \|q\|_{L^2(A)}+\|e\|_{L^2(B)}\ge |t|\,d_A,
 \quad d_A=\|aY_Z\|_{L^2(A)}>0 .
\end{equation}
The constant is attained by restricting and conformally flattening
the round TT lift of \(tf_r\), which has \(e=0\).
\end{theorem}

\begin{proof}
Write \(D_\pm=D_{\pm,\delta}\) in this proof and put
\(z=q-tTf-T\widetilde e\).
The identities \eqref{xfer:symbol} give \(D_-^*z=0\)
distributionally on \(B\), since \(D_-^*T=D_+^*\) and
derivatives of the zero extension agree with those of \(e\) on
compact tests in \(B\). The multiplier identities for \(L^2\) inputs
hold distributionally by polynomial multiplication on the Fourier
side. Consequently \(z\) is biharmonic by
\(D_-D_-^*=\Delta^2 I/16\) and is smooth in the ball by
elliptic regularity. This also gives the smooth representative
of \(q-T\widetilde e\) used for the all-radii projection.

For any \(\eta\in C_c^\infty((2/3,1))\) and any \(V\), set
\(\psi(s)=\eta(\sqrt{s})/s^2\) on the positive axis and extend
it by zero near zero. Define
\(\phi(s)=\tfrac16\int_s^\infty(v-s)^3\psi(v)\,dv\).
Then \(\phi^{(4)}=\psi\), and \(\phi\) is a polynomial of
degree at most three near zero and vanishes above the support of \(\psi\).
The metric test \(h_{\rm test}=D_+^*(\phi(\rho^2)V)\) is smooth
and compactly supported in \(B\). By \eqref{xfer:symbol},
\(D_-D_+^*=H_4(\partial)/16\). Using
\eqref{xfer:radial-derivative}, the factors of 16 cancel and
\(D_-h_{\rm test}=\rho^4\psi(\rho^2)Y_V=\eta(\rho)Y_V\).
Testing \(D_-^*z=0\) thus gives
\[
 \int_{2/3}^1 \eta(\rho)\rho^3
        \int_{S^3}\langle z(\rho n),Y_V(n)\rangle\,dS(n)\,d\rho=0 .
\]
Arbitrariness of the radial probe implies that each angular pairing
vanishes, which proves the first equality in \eqref{xfer:angular}.
Contractivity of the angular projection and
\(\|T\widetilde e\|_{L^2(\mathbb R^4)}=\|e\|_{L^2(B)}\) give the inequality.
For the restricted round lift, Proposition~\ref{xfer:multiplier}
gives \(q=tTf\) on \(B\), proving attainment.
\end{proof}

With \(a_0=6m_0/\pi^2\), \(a_2=12m_2/\pi^2\) and
\(I_\nu=\int_{2/3}^1\rho^{\nu-1}d\rho\), the constant in
\eqref{xfer:angular} satisfies
\[
 d_A^2=12\pi^2\bigl(a_0^2I_{-4}-2a_0a_2I_{-6}+a_2^2I_{-8}\bigr).
\]
For \(e=0\), the individual \(Y_V\) pairing equals
\(2\pi^2t a(\rho)\langle Z,V\rangle\) at every radius.
In particular the cross-pairings with \(V\perp Z\) vanish.
The other angular components remain unspecified, and equality
of the full norm at zero error requires their vanishing on \(A\).

\begin{remark}\label{xfer:relaxed}
The weaker equation \(\Delta^2 z=0\) alone allows both regular
degree-four radial coefficients \(\rho^4,\rho^6\), since
\(\mathcal Y\) has sphere Laplacian eigenvalue \(-24\). The usual
radial biharmonic equation has two other powers \(\rho^{-6},\rho^{-4}\),
excluded by regularity at zero. The relaxed norm bound would therefore
use the distance from \(a\) to
\(\operatorname{span}\{\rho^4,\rho^6\}\) in
\(L^2((2/3,1),\rho^3d\rho)\), times \(\sqrt{12\pi^2}\).
It is positive by independence of the powers and strictly smaller
than \(d_A\) because \(a\) has positive pairing with \(\rho^4\).
The compact tests use the full adjoint equation and rule out
both regular powers in the measured sector.
\end{remark}

\begin{corollary}\label{xfer:local-nonlinear}
On a supplied round patch, the conformal transfer
\eqref{xfer:weights} preserves \(d_A\). If a smooth trace-free
\(h\), sufficiently small in \(C^2\), in the fixed-volume chart satisfies
\(F_{r,+}(h)=tf_r\) there, then
\[
 \|W_{r_h}^-\|_{L^2(A,r_h)}\ge |t|d_A-C\|h\|_{C^2(B,r)}^2 .
\]
\end{corollary}

\begin{proof}
The Weyl and frame maps are smooth in the two-jet of \(h\), so
\(F_{r,\pm}(h)=D_{\pm,r}h+R_\pm\), with
\(\|R_\pm\|_{C^0}\le C\|h\|_{C^2}^2\) on the fixed patch.
Use \(-R_+\) for the linear source error and include \(R_-\)
in the opposite field. The frame and volume identities give
the claimed norm. The bound proves a nonzero response when
the stated loss is less than \(|t|d_A\).
\end{proof}

\subsection{Compact variation and the two center data}

\begin{proposition}[Compact support and the two radial moments]\label{xfer:compact}
For signed smooth radial \(\chi\) supported in \(\{\rho\le R\}\),
\(R\le1/2\), and a nonzero constant Weyl polarization \(V\),
\(D_{+,\delta}h=\chi V\) admits a smooth compactly supported
trace-free metric variation in \(B\) exactly when \(m_0=m_2=0\).
The same holds for the round source \(\Omega^{-2}\chi V\).
\end{proposition}

\begin{proof}
The flat adjoint kills \(V\) and \(\rho^2V\). In the algebraic
Weyl tensor realization \(T^z_{ij,pq}\) of a positive block
\(z\), the flat adjoint in operator normalization is
\((D_{+,\delta}^*z)_{jq}=-\tfrac12\partial_i\partial_p T^z_{ij,pq}\),
with repeated indices summed. This follows by contracting the
curvature derivative with \(T^z\) and integrating by parts twice. In
all ordered indices the operator pairing has factor \(1/4\),
the curvature symmetries combine the four second derivatives, and
the Ricci and scalar subtractions pair to zero.
For the quadratic field one obtains \(-\sum_iT^V_{ij,iq}=0\).
The two compact-support tests against the source give
\(|V|^2m_0\) and \(|V|^2m_2\), proving necessity.

For sufficiency take \(p=G*\chi\). It solves \(\Delta^2p=\chi\)
and vanishes outside the support radius by \eqref{xfer:radial}.
Thus \(h=16D_{+,\delta}^*(pV)\) is smooth and compactly supported
and solves the equation by \eqref{xfer:symbol}.
Equation~\eqref{xfer:weights} preserves support and
trace-freeness for the round lift.
\end{proof}

For an arbitrary smooth global background \(g\) with this round
patch, let \(L_g=dF_{g,+}(0)\). If \(\mathcal A\in\ker L_g^*\),
let \(\alpha\) be its round components and set \(z=\Omega^2\alpha\)
in the flat frame. Pairing against compactly supported round metric
variations and using \eqref{xfer:weights} shows \(D_{+,\delta}^*z=0\).
Then \(\Delta^2z=0\) by \eqref{xfer:symbol}. The sphere mean of \(z\)
is \(z(0)+\rho^2(\Delta z)(0)/8\), since its radial Laplacian
is the center value of the harmonic field \(\Delta z\).
Thus for the zero-extended round source the exact adjoint pairing is
\begin{equation}\label{xfer:center}
 \langle \Omega^{-2}\chi V,\mathcal A\rangle_{L^2(g)}
 =\bigl\langle V,m_0 z(0)+(m_2/8)(\Delta z)(0)\bigr\rangle .
\end{equation}
The weight on the adjoint incorporates both the Weyl weight and
the round volume. Across radial polarizations the pair of center
data lies in a space of dimension ten. Global adjoint solutions
may supply a smaller image. For the positive seed on the round sphere there is both nonzero
angular transmission and a vanishing global cokernel.

\subsection{An exact nonlinear lift on the standard sphere}

\begin{theorem}[Transmission to the opposite hemisphere]\label{xfer:nonlinear-round}
For the unit round \(r\) and the positive seed \(f_r\) above,
there is a smooth small fixed-volume family \(r_{h(t)}\) with
\(F_{r,+}(h(t))=tf_r\) and \(h(0)=0\).
For sufficiently small nonzero \(t\) the entire closed southern
hemisphere is ASD and nowhere Weyl-flat. Its squared negative
Weyl energy is
\begin{equation}\label{xfer:south-energy}
 \int_{S_-}|W_{r_{h(t)}}^-|^2\,dV_{r_{h(t)}}
 =\frac{432}{\pi^2}
   \left(\frac{m_0^2}{4}-\frac{2m_0m_2}{3}+\frac{m_2^2}{2}\right)t^2
   +O(|t|^3).
\end{equation}
The coefficient of \(t^2\) in \eqref{xfer:south-energy} is strictly
between zero and \(\|f_r\|_{L^2(r)}^2\). The family has positive
Yamabe constant.
\end{theorem}

\begin{proof}
On the fixed round TT slice the Banach implicit function theorem
in \(C^{k+2,\alpha}\to C^{k,\alpha}\) with \(k\ge4\) applies to
\(F_{r,+}\), whose derivative has inverse \(Q_r\). It gives
\(h(t)=tQ_rf_r+O(t^2)\) in that grade. For spatial smoothness
write \(h=D_{+,r}^*v\) using the inverse of
\(D_{+,r}D_{+,r}^*\). The resulting equation in \(v\) is a
fourth-order quasilinear elliptic system whose leading coefficients
are close to the round ones, and the source is smooth.
The elliptic regularity argument of Section~\ref{sec:lift}
gives a smooth family. Both chiral blocks and the compatible
chiral connection come from this one metric. Thus the two
hemispheres have the same interface jets in a common chart,
with opposite outward normals.

The positive block is zero on a neighborhood of \(S_-\).
Taylor expansion gives \(F_{r,-}(h(t))=tU+O(t^2)\)
uniformly on the sphere. By \eqref{xfer:pole},
\(|U|\ge 3\sqrt6\,m_0/(2\pi^2)>0\) on \(S_-\), and the
remainder then gives the claimed absence of Weyl-flat points.
For the squared energy use norm and volume preservation to
obtain \(t^2\|U\|_{L^2(S_-,r)}^2+O(|t|^3)\).
The squared conformal weight cancels the volume,
and \(|Y_Z|^2=6\). Hence the coefficient is
\[
 6\cdot2\pi^2(6/\pi^2)^2
 \int_1^\infty (m_0\rho^{-4}-2m_2\rho^{-6})^2\rho^3\,d\rho ,
\]
which is \eqref{xfer:south-energy}. The linear transfer is isometric on the full sphere. Its nonzero
opposite-chiral field at the equator remains nonzero on an adjacent
northern open set by continuity. Hence the southern integral of the
linearized field is strictly less than its full squared norm.
The scalar curvature
is positive by \(C^2\) closeness to the unit round metric,
where it is \(12\). The compact Sobolev inequality then bounds
the Yamabe quotient below by a positive constant
\cite[Remark~1]{ChangGurskyYang2003}.
\end{proof}

The exact prescribed positive block also gives
\(E_+=6t^2\int_{\mathbb R^4}\chi^2\,dx\) on this sphere.
In the operator normalization \(E_+-E_-=12\pi^2\tau\), so the two
global energies agree on any homotopy sphere~\cite{Hirzebruch1971Signature}.
In particular, any supplied opposite energy in a retained collar
contributes a lower bound to the global \(E_+\). In the region
where the positive block vanishes the corresponding twistor
almost complex structure is integrable, with orientation chosen
accordingly~\cite[Theorem~4.1]{AHS1978}. The nonzero opposite
block obstructs the opposite-chiral integrability condition.

On an unknown smooth homotopy sphere the local theorem applies
inside a supplied round chart, with no hypothesis on the smooth
geometry of its complement. It fixes the stated angular projection
there and gives no exterior pointwise formula. The balancing
construction of Section~\ref{sec:lift} can allow compensating
sources supported in the actual exterior. The weighted center
pairings in \eqref{xfer:center} test the initial linear range
condition for radial data. Recovering the unmodified prescription
on that slice requires cancellation of the exterior coefficients.
Section~\ref{msreg:section} changes the reference at fixed energy, and
Section~\ref{mscone:section} gives a sufficient sign condition at the original reference.
Continuation in the original smooth structure and the estimates
needed for a recognizing metric remain global requirements.

\section{What fixed boundary data can obstruct}
\label{sec:collars}

The spherical calculation allows the boundary metric and its normal
derivatives to respond to the source on either side. Fixing the metric
on both ends of a flat collar leads to further compatibility conditions.
To keep track of a change of reference, write
\(I_{r,g}:\Lambda_g^+\to\Lambda_r^+\) for pullback by the positive
symmetric tangent isometry \((r^{-1}g)^{-1/2}\), and write
\(\mathcal F_r(g)=I_{r,g}W_g^+I_{r,g}^{-1}\). A metric variation in the
pointwise-volume chart is \(g_h=r(e^h\cdot,\cdot)\), where \(h\) is
self-adjoint and trace-free. Write \(D_r\) for the derivative at \(h=0\) of \(h\mapsto \mathcal F_r(g_h)\),
the operator denoted by \(L\) in Section~\ref{sec:lift}.
There are distinct obstructions to a uniform right inverse, to a
particular prescribed source, and to a retained conformal structure.

\subsection{Small singular values and the prescribed direction}

Let \(M\) be a closed oriented smooth homotopy four-sphere. A product of two circles in an oriented \(\mathbb R^4=\mathbb C^2\)
chart has a normal framing given by the two radial directions. A shell
in its normal disk bundle embeds a separating \(T^3\times[-1,1]\).
Choose smooth metrics on both complementary caps that agree with the flat
product in the seam collars. Lengthen the cylinder to
\(T^3\times[-L,L]\) using an increasing radial coordinate of derivative
one near both ends. On a fixed smooth manifold this gives a family
\(r_L\) with fixed cap metrics, fixed \(E_+\), and volume
\(V_L=V_{\rm caps}+2L\,\operatorname{Vol}(T^3)\).

For a parallel unit tangential trace-free matrix \(H\) let \(Z\) be
its unit image in the parallel self-dual frame. At the flat product,
a variation \(a(s)H\) has radial curvature block \(-a''H/2\).
Its self-dual projection has an additional factor \(1/2\). All the
nonradial flat derivatives integrate to zero on tests constant in the
torus variables. Thus the flat formulas for their compactly supported
profiles are
\[
 D_{r_L}(aH)=-\tfrac14 a''Z,\qquad
 D_{r_L}^*(aZ)=-\tfrac14 a''H .
\]
The principal symbol of \(D_rD_r^*\) at any smooth reference is
\(|\xi|_r^4 I/16\), by aligning \(\xi\) with the normal and using the
five-dimensional trace-free tangential fiber. Its orthogonal complement
is the infinitesimal conformal Lie symbol, which the Weyl symbol kills.
The ASD deformation complex provides the corresponding comparison at
an ASD background \cite[\S1.1]{GoverGursky2024}.
For any fixed \(j\), choose \(j\) nonzero profiles
\(\beta_i\in C_c^\infty((-1,1))\) with disjoint supports. The fields
\(\beta_i''(s/L)Z\) have compact preimages
\(-4L^2\beta_i(s/L)H\), so they are orthogonal to \(\ker D_{r_L}^*\).
Their Rayleigh quotients, and that of any nonzero combination, are bounded
by \(C_j L^{-4}\). Min--max on the orthogonal complement of the kernel
gives this bound on the \(j\)-th strictly positive eigenvalue of
\(D_{r_L}D_{r_L}^*\). Under \(r\mapsto c r\) this normal operator scales
by \(c^{-2}\) in the isometric frames, so unit-volume normalization
multiplies its eigenvalues by \(V_L\). Every fixed positive eigenvalue
rank tends to zero at least as fast as \(L^{-3}\) in the normalized
family, still at fixed \(E_+\). A unit target eigensection with eigenvalue
\(\lambda>0\) requires an \(L^2\) preimage of norm at least
\(\lambda^{-1/2}\) by adjoint pairing. These particular small-Rayleigh
test fields themselves pair to zero with the radial Weyl source,
which vanishes on their supports.

The source can lie in the adjoint kernel in a different geometry.
At a closed Einstein reference the first variation of the integral
of squared scalar curvature in the pointwise-volume chart vanishes because
the scalar variation is a divergence for trace-free \(h\). The
squared trace-free Ricci integral has zero first variation as well.
The Euler and signature identities then give
\(D_r^*W_r^+=0\), since
\(dE_+(h)=2\langle D_rh,W_r^+\rangle_{L^2(r)}\)
\cite[(0.4) and Remark 2]{ChangGurskyYang2003}.
The product of unit round spheres \(S^2\times S^2\) is Einstein with
\(W^+\) having entries \((2/3,-1/3,-1/3)\).
It therefore cannot start a differentiable path with source
\((1-t)W_r^+\), whose initial tangent would require \(D_r h=-W_r^+\).
This example has zero signature and \(b_2=2\).

\subsection{A fixed two-ended chord and corrected sources}

Let \(I=[-1,1]\), give \(T^3\) its product angular coordinates, and
orient \(I\times T^3\) by \(ds\wedge dx_1\wedge dx_2\wedge dx_3\).
Both open end collars are required to retain the exact flat product metric. For
\(g_a=ds^2+\sum_i e^{2a_i(s)}dx_i^2\), with \(\sum_i a_i=0\) and
all \(a_i\) zero near the ends, put \(p_i=a_i'\).
The radial and tangential sectional curvatures are
\(-p_i'-p_i^2\) and \(-p_jp_k\), respectively. Taking the half-sum
in the normalized chiral frame and subtracting scalar curvature divided
by twelve gives the ordered Weyl entries
\[
 w_i=-\tfrac12p_i'-p_i^2+\tfrac13\sum_j p_j^2.
\]
The symmetric transport between diagonal metrics takes these indexed
frames to each other even at repeated eigenvalues.

\begin{proposition}\label{bdry:axial-chord}
For a nonconstant smooth interior seed \((f,f,-2f)\), the chord
\(\lambda W_{\rm seed}^+\) in the seed frame has no two-ended
diagonal synchronous lift at \(0<\lambda<1\).
\end{proposition}
\begin{proof}
Equality of the first two Weyl entries gives
\((p_1-p_2)'+2(p_1+p_2)(p_1-p_2)=0\).
The left flat end fixes \(p_1=p_2\) and \(a_1=a_2=a\) throughout.
Writing \(y=e^{-2a}\), \(y_0=e^{-2f}\) and \(V_0=y_0''/y_0\), the
remaining Weyl block is \((y''/4y)\operatorname{diag}(1,1,-2)\).
The prescribed equation is \(-y''+\lambda V_0y=0\) with Neumann
ends and \(y=1\) there. For the associated forms,
\[
 Q_0[v]=\int_I|v'|^2,\qquad
 Q_1[v]=\int_I y_0^2|(v/y_0)'|^2.
\]
The second equality follows by integration by parts using
\(y_0'=0\) at the ends. The forms are nonnegative with distinct
null lines \(\mathbb R\) and \(\mathbb R y_0\). Their strict convex
combination is positive on every nonzero test. A Neumann solution
would have \(Q_\lambda[y]=0\), contrary to \(y>0\).
\end{proof}

An even seed admits a one-profile correction at any amplitude. Fix any
smooth nonzero even \(p\ge0\) with interior support and replace the
potential by \(\lambda V_0-\beta p\). The Neumann ground eigenvalue is
continuous and strictly decreasing in \(\beta\) by testing against
its positive normalized eigenfunction. A test supported where
\(p>0\) makes it negative for sufficiently large \(\beta\). Proposition
\ref{bdry:axial-chord} gives positive ground eigenvalue at
\(\beta=0\) for interior \(\lambda\), so a unique \(\beta(\lambda)>0\)
sets it to zero. At both parameter endpoints \(\beta=0\). The
positive eigenfunction is simple by Neumann initial-value uniqueness
and is even by reflection. It is constant on each flat end and the
two constants agree, so it normalizes to \(1\) on both. To check smooth
dependence, shoot with \(y(-1)=1,y'(-1)=0\). At a positive zero of
\(y'(1)\), differentiating the equation in \(\beta\) and using Green's
identity gives \(\partial_\beta y'(1)=-\int_I p y^2/y(1)<0\).
The resulting metrics solve the corrected source
\((\lambda V_0-\beta p)\operatorname{diag}(1,1,-2)/4\).

The size of the correction can also be computed. For \(f=\varepsilon\phi\)
with \(\phi\) a fixed nonconstant even interior bump,
\(V_0=-2\varepsilon\phi''+4\varepsilon^2(\phi')^2\).
At \(\varepsilon=0\) the shooting derivative is \(-\int_Ip\),
uniformly in \(0\le\lambda\le1\). Integrating the first-order
equation with Neumann data gives \(\partial_\varepsilon\beta=0\)
and \(\partial_\varepsilon y=-2\lambda\phi\). At second order an
integration by parts gives
\[
 \beta(\lambda,\varepsilon)=
 4\varepsilon^2\lambda(1-\lambda)\frac{\int_I|\phi'|^2}{\int_Ip}
 +O\bigl(|\varepsilon|^3\lambda(1-\lambda)\bigr).
\]
The remainder before factoring is \(O(|\varepsilon|^3)\) in
\(C^1([0,1]_\lambda)\) by smooth shooting. It vanishes at both
endpoints, which supplies the stated endpoint factor.

For a matrix version let \(V=\operatorname{Sym}_0(3,\mathbb R)\),
fix \(K\Subset I^\circ\) an interval with nonempty interior, and
take \(A\in C_K^{k+2,\alpha}(I,V)\), \(k\ge2\), \(0<\alpha<1\).
The subscript indicates support in \(K\) as a function on \(I\).
Use the metric \(ds^2+dx^{\mathsf T}Gdx\), \(G=e^{2A}\),
and a seed \(A_0\) in the same space. In the seeded orbit frame
the positive matrix is \(H=e^{-A_0}G e^{-A_0}\), and the
transported tangent frame has columns \(Q=e^{-A_0}H^{-1/2}\).
This order matters for noncommuting matrices. Put
\[
 S=\tfrac12 Q^{\mathsf T}G'Q,\qquad
 R=-\tfrac12 Q^{\mathsf T}G''Q+
       \tfrac14 Q^{\mathsf T}G'G^{-1}G'Q .
\]
The normal Riccati formula gives \(R\), while Gauss and Codazzi give
tangential block \(-\operatorname{cof}S\) under the tangential Hodge
identification and zero mixed block. Hence the transported source map
is \(\mathcal F_{A_0}(A)=\tfrac12\operatorname{tf}(R-\operatorname{cof}S)\).
It depends smoothly on \((A_0,A,A',A'')\) and at the flat origin its
metric derivative is \(-\tfrac12\partial_s^2\).

On supported profiles, the range of this derivative is exactly the
kernel of \(\mu(y)=(\int_I y,\int_I s y)\). Indeed, integrating twice
from the left endpoint gives
\(-2\int_{-1}^s(s-t)y(t)\,dt\), which vanishes to the right of \(K\)
precisely at zero moments. The kernel is zero and the integral
operator is bounded in the indicated Hölder grades. Take two fixed
interior scalar bumps with independent moment columns and multiply by
a basis of \(V\), defining a source insertion \(\mathcal P\) with ten
coefficients. Augmenting by \(\mathcal P c\) now gives an isomorphism
of derivatives. The Banach inverse function theorem with seed as
parameter gives a smooth local inverse \((A,c)\) to the augmented map.
The unaugmented image in \(C_K^{k,\alpha}(I,V)\) is the split
codimension-ten submanifold specified by \(c=0\).
For a small smooth seed the whole segment
\(y_\lambda=\lambda\mathcal F_{A_0}(A_0)\) lies in this chart.
The correction coefficients are zero at both endpoints. The smooth
inverse has bounded second target derivative on a smaller neighborhood,
so \(\sup_\lambda|c''(\lambda)|\le C\|A_0\|_{k+2,\alpha}^2\).
The Dirichlet Green function on \([0,1]\) gives
\[
 |c(\lambda)|\le C\lambda(1-\lambda)\|A_0\|_{k+2,\alpha}^2,\qquad
 \mathcal F_{A_0}(A(\lambda))=y_\lambda-\mathcal P c(\lambda).
\]
The coefficient of \(A''\) is invertible near zero by its flat value
\(-I/2\), yielding smooth profiles by the ordinary differential
equation. This construction is within the synchronous torus-invariant
class, with the flat end data exactly preserved.

\subsection{Unrestricted flat-end obstructions}

Write \(\gamma=ds^2+\sum_i dx_i^2\) for the flat collar, taking angular
periods \(2\pi\). In the flat self-dual frame the axial seed transported
from its own frame gives
\[
 T=\left(-\tfrac12 f''+(f')^2\right)Z_0,\qquad
 Z_0=\operatorname{diag}(1,1,-2).
\]
For an arbitrary compactly supported trace-free variation \(h\),
the flat double-divergence formula for the adjoint annihilates every
parallel Weyl target. In particular
\(\int\langle D_\gamma h,Z_0\rangle=0\), but
\(\int\langle T,Z_0\rangle=6\operatorname{Vol}(T^3)\int_I(f')^2>0\).
This excludes a differentiable lift through the flat metric with
nonzero chord speed. The same pairing and the quadratic curvature
remainder imply $\lambda\le C\|g-\gamma\|_{C^2}^2$ for any nearby
fixed-end solution of the original chord. For general boundary values, the pairing
equals the flux
\(-\frac14[\int_{T^3}\partial_s(h_{11}+h_{22}-2h_{33})]_{-1}^1\),
as all terms with tangential derivatives integrate out.

There are also infinitely many tangential adjoint conditions.
Set \(z=ds+i\,dx_1\), \(w=dx_2+i\,dx_3\) and \(\omega=z\wedge w\).
For each positive integer \(m\),
\(e^{m(s+ix_1)}\omega\otimes\omega\) is a complex self-dual Weyl
tensor: \(\omega\) is simple and null, and its Ricci contraction is
zero because \(z,w\) span a totally null plane. The first divergence
vanishes by \(\iota_{z^\sharp}\omega=0\), so the flat Weyl adjoint
annihilates it by its double-divergence formula
\cite[\S1.1]{GoverGursky2024}. Its real part defines a real pairing on
compact target sources. Multiplying these real parts by a fixed
nonnegative interior cutoff gives compact source tests with a diagonal
positive pairing matrix by Fourier orthogonality in \(x_1\). The
functionals are independent and annihilate the image of compact
metric variations. Consequently no fixed finite-dimensional space
of interior source profiles repairs the full supported linear range.

\begin{theorem}\label{bdry:unrestricted-near-flat}
For each fixed nonconstant smooth axial seed as in Proposition
\ref{bdry:axial-chord}, a \(C^{10}\) neighborhood of the product flat
metric contains no smooth metric that equals the flat product on open
collars of both ends and realizes the uncorrected seed-frame chord
\(\lambda W_{\rm seed}^+\) for any \(0<\lambda\le1\). The candidate
may vary in all four coordinates and have arbitrary lapse, shift and
volume density.
\end{theorem}
\begin{proof}
Identify the end tori using their product coordinates. The metrics
and the equation extend smoothly onto the auxiliary rectangular
flat \(T^4\), since the data agree on open end collars. Write \(r\)
for the seed and keep \(\gamma\) for the periodically extended flat
metric. The prescribed flat target \(T=I_{r,\gamma}^{-1}W_r^+I_{r,\gamma}\)
has nonzero flat average. All Sobolev norms in this proof use the
flat torus, and \(C^{10}\) closeness on the collar controls
\(H^{10}\) closeness on the torus independently of the end widths.

For a mean-zero displacement \(X\in H^8\), use the small diffeomorphism
\(\phi(x)=x+X(x)\) and set \(P=\gamma^{-1}\phi^*g\),
\(\rho=(\det P)^{-1/4}\), \(\widetilde g=\rho\phi^*g\).
Its logarithm is \(h=\log P-\tfrac14(\operatorname{tr}\log P)I\).
Solve \(\delta_\gamma h=0\), with divergence defined by
\((\delta_\gamma h)_j=\partial^i h_{ij}\). Its mean vanishes by periodicity.
The gauge map \(H^{10}\times H^8_0\to H^6_0\) is \(C^1\). Indeed, composition uses one spare metric derivative into \(H^8\), the two
factors \(d\phi\) leave \(H^7\), and divergence costs one more derivative
\cite[Theorem 1.2]{InciKappelerTopalov2012Composition}.
Its displacement derivative is
\(\Delta_\gamma X^\flat+\tfrac12d\operatorname{div}_\gamma X\) for
\(\Delta_\gamma=\partial^i\partial_i\), invertible on nonzero Fourier
modes with eigenvalues \(-|\xi|^2\) and \(-3|\xi|^2/2\). The implicit
function theorem gives \(X\), \(\rho-1\) and \(h\) small in the stated
grades. These grades suffice for curvature and pullback. For smooth
inputs, translation equivariance
and the implicit function theorem with translated smooth input imply
smooth dependence of \(X(\cdot+a)\) in \(H^8\) on \(a\), giving all
spatial derivatives. The gauge takes place after completing the torus.

Write \(h=c+v\), with \(c\) constant and \(v\) mean-zero TT. The
constant matrix gives a flat metric, so its Weyl target vanishes.
On the nonzero Fourier modes the flat Weyl derivative maps TT tensors
isometrically up to the factor \(|\xi|^2/4\), hence has bounded
inverse from mean-zero Weyl \(H^5\) to mean-zero TT \(H^7\).
Write \(F(h)=F_\gamma(\widetilde g)\). Smoothness of \(F:H^7\to H^5\)
and \(F(c)=0\) give
\[
 \|F(c+v)-D_\gamma v\|_{H^5}
 \le C(\|c\|_{H^7}+\|v\|_{H^7})\|v\|_{H^7}.
\]
Integrate the derivative along \(c+\theta v\) to obtain this estimate.
Absorption using the flat inverse gives
\(\|v\|_{H^7}\le C\|F(h)\|_{H^5}\). Averaging then gives
\[
 |\overline{F(h)}|\le C\|h\|_{H^7}\|F(h)\|_{H^5}.
\]
Thus the estimate accounts for the constant flat moduli without
inverting a constant target.

The target moves under this gauge. The tangent map
\(O_x=\rho(x)^{1/2}d\phi_x(\gamma^{-1}\widetilde g)^{-1/2}_x\)
is an oriented isometry from \(\gamma_x\) to \(g_{\phi(x)}\).
Let \(J_O\) be its pullback on positive forms, and, at the
indicated basepoints in the parallel flat frames, set
\[
 U_g=J_O(I_{r,g}\circ\phi)^{-1}(I_{r,\gamma}\circ\phi).
\]
Naturality, Weyl-operator conformal weight, and the seed-frame chord
give \(\lambda^{-1}F(h)=\rho^{-1}U_g(T\circ\phi)U_g^{-1}\).
This order retains all three identifications. As \(g\to\gamma\),
the composition estimates and the matrix algebra in \(H^7\) give
\(U_g\to I\), \(\rho\to1\), and \(T\circ\phi\to T\) in \(H^7\).
The normalized target therefore tends to \(T\) there and stays
bounded in \(H^5\). Its mean is bounded away from zero, whereas
the averaged estimate divided by \(\lambda\) makes that mean
tend to zero. This excludes every sufficiently near-flat metric
uniformly over the indicated chord parameters.
\end{proof}

\subsection{Conformal holonomy of a retained flat product}

\begin{proposition}\label{bdry:holonomy}
A smooth oriented homotopy four-sphere cannot carry a global ASD metric
whose conformal structure on an embedded flat product collar equals
the flat product conformal structure.
\end{proposition}
\begin{proof}
The signature identity forces any global ASD metric on this manifold
to be Weyl-flat. Its locally conformally flat structure has a
single-valued developing map because the manifold is simply connected
\cite[\S1]{SchoenYau1988Conformal}. On the collar, developing its
universal cover by the Euclidean inclusion gives holonomy
\(x\mapsto x+\ell e\) for a primitive lattice vector of length \(\ell>0\)
and unit direction \(e\). This is nontrivial up to Möbius conjugacy.
A restriction from the global conformally flat structure would have
trivial holonomy on that loop. The flat Levi--Civita linear holonomy
of the product is the identity, since the parallel frame descends.
\end{proof}

A conditional quantitative version uses any smooth metric extension
equal to the product on the collar. In its metric splitting, use
the positive tractor norm
\(\|(\sigma,\mu,\varrho)\|^2=\sigma^2+|\mu|^2+\varrho^2\).
The normal tractor connection on the product has zero Schouten tensor
and the off-diagonal terms in direction \(e\) are
\((-\mu(e),\varrho e^\flat,0)\). Parallel transport around the loop is
\[
 T_\ell(\sigma,\mu,\varrho)=
 (\sigma+\ell\mu(e)-\tfrac12\ell^2\varrho,\,
  \mu-\ell\varrho e^\flat,\,\varrho).
\]
In particular its distance from the identity in operator norm is
at least \(\delta_\ell=(\ell^2+\ell^4/4)^{1/2}\).
Let \(F(u,t)\) be a smooth based contraction of the loop, and
\(S_u(t_2,t_1)\) the normal tractor transport along its \(t\)-segments.
Differentiating the parallel equation gives, with the commutator sign
for tractor curvature \(\mathcal R^{\mathcal T}\),
\[
 \partial_u S_u(1,0)
 =-\int_0^1 S_u(1,t)\mathcal R^{\mathcal T}(F_u,F_t)S_u(t,0)\,dt .
\]
If the segment operators are bounded by \(B_*\) in the endpoint
positive norms and the filling area with multiplicity outside the
product is at most \(A_*>0\), integration gives
\[
 \sup_{F([0,1]^2)\setminus U} (|W|_{\mathrm{op,HS}}+|C|_*)
 \ \ge\ \frac{\delta_\ell}{B_*^2 A_*}.
\]
Here \(U\) is the retained product region and \(|C|_*\) is the
Cotton comass on unit bivectors. The normal tractor curvature is
bounded by the displayed sum in the chosen splitting and vanishes on
\(U\) \cite[(22) and (32)]{CurryGover2015Introduction}. The area and
transport controls are scale-dependent filling hypotheses. The
estimate measures full tractor curvature along a two-dimensional
filling. The four-dimensional chiral energy threshold requires
separate control. The fixed-end and holonomy obstructions motivate
allowing the metric and its normal derivatives to respond on both
sides of an interface with trivial conformal holonomy.

\section{A marked ball and the Weyl variational problem}\label{sec:variation}

The preceding deformation problem prescribes an entire curvature field.
For smooth recognition it can suffice to control its energy. We now
vary the metric freely outside one coordinate ball and ask what a
stationary extension would imply. Throughout this section \(M\) is a
closed, connected, oriented smooth homotopy four-sphere.
Write \(Y([g])=Y(M,[g])\) for the Yamabe constant of one conformal
class. Fix an oriented extended chart \(\psi:B_{1+d}(0)\to M\),
\(d>0\), and set \(K=\psi(\overline B_1)\). A marked metric has
\(\psi^*g=\Omega^2\delta\) on \(\overline B_1\), where
\(\Omega(x)=2/(1+|x|^2)\) and \(\delta\) is Euclidean. Let
\(\mathcal G\) be the smooth metrics on \(M\) and \(\mathcal G_K\) its
marked subclass. Superscripts \(+\) on these two spaces impose \(Y>0\).

\begin{proposition}[Approximation with a prescribed round ball]\label{var:marking}
For every \(g\in\mathcal G\) and \(\varepsilon>0\) there is
\(g_K\in\mathcal G_K\) such that
\[
 |E_+(g_K)-E_+(g)|<\varepsilon,\qquad
 |Y([g_K])-Y([g])|<\varepsilon .
\]
Thus \(\inf_{\mathcal G_K}E_+=\inf_{\mathcal G}E_+\). If
\(\mathcal G^+\) is nonempty, so is \(\mathcal G_K^+\), and
\(\inf_{\mathcal G_K^+}E_+=\inf_{\mathcal G^+}E_+\).
\end{proposition}

\begin{proof}
Apply the approximation theorem of Akutagawa, Botvinnik, Kobayashi
and Seshadri at \(o=\psi(0)\). It produces a smooth metric \(\widetilde g\)
that is conformally flat near \(o\), with arbitrarily small changes
in \(Y\) and in the full Weyl functional. In dimension four their
functional is \(\mathcal W(g)=\int_M|W_g|_{(0,4)}^2\,dV_g\).
Signature and the operator norm convention give
\(\mathcal W=4(E_++E_-)=8E_+\) on \(M\). Thus the approximation
controls both quantities in the assertion~\cite[\S1.2 and Proposition 4.1]{AkutagawaEtAl2002Weyl}.

Choose oriented coordinates \(y\) in the flat conformal germ,
with \(y(o)=0\) and \(y_*\widetilde g=e^{2v}\delta\). The transition
germ \(F=\psi^{-1}\circ y^{-1}\) fixes \(0\). It can be realized by
a diffeomorphism \(\Xi\) of \(M\) supported in the extended chart.
To see this, join \(I\) to \(DF(0)\) in \(\mathrm{GL}^+(4)\) and
join \(DF(0)\) to \(F\) by \(t^{-1}F(tx)\), extended smoothly to
\(t=0\). On a common small domain these are isotopies of local
diffeomorphisms. Cut off the time-dependent velocity field by a
bump equal to one near \(0\) and integrate it. Trajectories starting
sufficiently near \(0\) stay where the cutoff is one, so the time-one
map agrees with the prescribed germ there.

On a neighborhood of \(\overline B_1\), realize \(x\mapsto cx\),
\(0<c<1\), by the flow of a cut-off radial vector field in the
extended chart. Call the resulting diffeomorphism \(D_c\).
For sufficiently small \(c\), \(\phi=\Xi\circ D_c\) satisfies
\(y(\phi(\psi(x)))=cx\) on a neighborhood of \(\overline B_1\).
There \(\psi^*\phi^*\widetilde g=c^2e^{2v(cx)}\delta\).
Multiply by the square of the positive factor
\(\Omega(x)/(c e^{v(cx)})\), extended to \(M\) using a cutoff
of its logarithm. This gives the prescribed round representative
on \(K\). The two invariants are preserved by these diffeomorphism
and conformal operations. Taking infima proves the first equality.
For the positive subclass choose the approximation error smaller
than the starting \(Y\) before taking its infimum.
\end{proof}

\subsection{Exterior stationarity}

Let \(B_g\) be the covariant Bach tensor normalized so that the
gradient of the full four-index Weyl functional is \(-4B_g\).
It is symmetric and trace free, and
\(B_{u^2g}=u^{-2}B_g\) for smooth \(u>0\). The preceding norm
conversion yields
\begin{equation}\label{var:bach}
 dE_+(g)[h]=-\tfrac12\int_M\langle B_g,h\rangle_g\,dV_g .
\end{equation}
The conformal factors in \(|B_{u^2g}|_{u^2g}=u^{-4}|B_g|_g\)
and \(dV_{u^2g}=u^4dV_g\) cancel. In particular the total
Bach norm \(\int_M |B_g|_g\,dV_g\) is invariant under conformal
change and diffeomorphism~\cite[\S1 and (1.1)]{BahuaudGuentherIsenbergMazzeo2025Bach}.

\begin{lemma}[Continuation of a vanishing Weyl block]\label{var:continuation}
A positive \(C^4\) Bach-flat metric on \(M\) with \(W^+=0\) on a
nonempty open set satisfies \(W^+=0\) everywhere.
\end{lemma}

\begin{proof}
In dimension four the ambient obstruction is a fixed nonzero
multiple of the Bach tensor. For the vanishing equation the
normalizations can also be compared by varying
\(\int_M Q\,dV=8\pi^2\chi(M)-\mathcal W/4\). The first variation of
\(\int Q\) is a nonzero multiple of the obstruction pairing, and
the identity follows from the four-dimensional formula for \(Q\)
and Chern--Gauss--Bonnet~\cite[\S1 and Theorem 1.1]{GrahamHirachi2004Ambient}.
This comparison extends to \(C^4\) metrics by smooth approximation
in \(C^4\).

Fix \(0<\beta<1\). In a conformal harmonic chart the coordinate
map is \(C^{4,\beta}\) for a \(C^{3,\beta}\) metric. To transfer
the obstruction equation at this grade, approximate the inverse
chart locally in \(C^4\) by smooth local diffeomorphisms and the
original metric in \(C^4\) by smooth positive metrics. Their
pullbacks converge in \(C^3\) in the new chart. The Bach formula
is continuous into distributions under this convergence. Its
highest terms have two derivatives of Schouten, which leave
only one such derivative after integration by parts against a
test function. The other terms involve at most three derivatives
of the metric. Naturality for the smooth approximations therefore
transfers the zero equation to the chart. Theorem 1.3(c) of
Lassas and Liimatainen then makes the determinant-normalized
representative \((\det g)^{-1/4}g\) real analytic. Its hypothesis
\(r>3\) is met by \(r=3+\beta\), and the weighted obstruction is
identically zero~\cite[Theorems 1.3(c), 2.2 and 3.7(c), and \S3.2]{LassasLiimatainen2019ConformalHarmonic}.

The Hodge star and the Weyl endomorphism of this representative
have analytic coefficients. In particular \(P_+WP_+\), where
\(P_+=(I+*)/2\), is analytic in each chart. Positive conformal
rescaling preserves \(P_+\) on two-forms and preserves the zero set
of the block. Let \(V\) be the set of points with a neighborhood
on which \(W^+=0\). Any point in \(\overline V\) has a connected
conformal harmonic chart meeting \(V\) in an open subset. The
analytic identity theorem for the projected tensor implies that
this whole chart lies in \(V\), proving the result by connectedness.
This uses the local analytic regularity theorem. The authors'
Theorem 1.4 has a hypothesis on the full Weyl tensor.
\end{proof}

\begin{proposition}[Stationarity with a fixed ASD cap]\label{var:stationarity}
Suppose the metric on \(K\) is fixed to a smooth ASD metric with
a smooth positive metric germ across \(\partial K\).
A smooth extension is stationary for \(E_+\) against all
compactly supported trace-free variations in \(M\setminus K\)
if and only if it is conformally diffeomorphic to the round sphere.
\end{proposition}

\begin{proof}
The test \(h=\chi B_g\) for nonnegative smooth bumps supported in
\(M\setminus K\) shows by \eqref{var:bach} that \(B_g=0\) there.
For auxiliary variations supported in \(K^\circ\), the derivative
of the squared \(W^+\) density vanishes because \(W^+=0\) on the
support. The same variational formula on the closed \(M\) gives
\(B_g=0\) in \(K^\circ\). Smoothness of \(g\) gives the zero
equation across the interface. Lemma~\ref{var:continuation}
and signature now give both Weyl blocks zero. The conformal
developing map of the closed simply connected \(M\) gives a
conformal diffeomorphism to \(S^4\). Conversely a conformally
round extension has energy zero and is stationary.
\end{proof}

If the ASD cap has \(W^-\ne0\) somewhere, every smooth extension
has a strictly decreasing exterior direction. More generally
at any nonstationary extension choose \(\chi\ge0\) supported
in the exterior with \(A_g=\int_M\chi|B_g|^2\,dV_g>0\).
The variation \(h=\chi B_g\) is realized by
\(g_t=g(e^{t g^{-1}h}\cdot,\cdot)\), fixed on \(K\) and
preserving volume pointwise. Its energy is at most
\(E_+(g)-tA_g/4\) for all sufficiently small positive \(t\), since
its derivative at zero is \(-A_g/2\).

\subsection{Conditional compactness on the same manifold}

\begin{theorem}[Recognition from a surviving open cap]\label{var:compactness}
Let \(g_j\in\mathcal G_K\). Suppose that after smooth
diffeomorphisms \(\phi_j:M\to M\) and positive conformal changes
a subsequence of the metrics \(G_j=u_j^2\phi_j^*g_j\) has a uniform
positive lower bound and converges strongly in \(C^4\) in a fixed smooth atlas.
Assume some nonempty open set \(U\) lies in
\(\phi_j^{-1}(K^\circ)\) along this subsequence.
If either \(g_j\) minimizes \(E_+\) asymptotically over
\(\mathcal G_K\) or \(\int_M |B_{g_j}|_{g_j}\,dV_{g_j}\to0\), then
the limit is Weyl-flat and \(M\) is smoothly diffeomorphic to
\(S^4\). In the minimizing case the marked and unrestricted infima are zero.
\end{theorem}

\begin{proof}
A uniform \(C^{4,\alpha}\) bound for \(0<\alpha<1\) gives the
stated convergence after extraction, provided the positive
lower bound and common open set are retained. Write \(G\) for
the positive limit. The Weyl energies are invariant under the
given gauges. Orientation reversal merely exchanges the equal
chiral energies on this \(M\). The round cap gives \(W^+=0\) on
\(U\) for \(G_j\) and hence for \(G\) by \(C^2\) convergence.

In the minimizing case, Proposition~\ref{var:marking} gives a
common infimum \(I\) for \(\mathcal G_K\) and \(\mathcal G\).
For any smooth symmetric \(h\), \(G_j+sh\) is positive for
sufficiently small \(|s|\) independently of \(j\). Using the
global infimum at each fixed \(s\) and passing to the limit gives
\(I\le E_+(G+sh)\) and \(I=E_+(G)\). The density is a smooth
function of the positive metric and its two-jet. Thus the
first variation against fixed \(h\) is continuous under \(C^2\)
convergence. The Bach tensor is continuous under \(C^4\)
convergence, so \eqref{var:bach} passes to \(G\) and the two-sided
minimum gives \(B_G=0\). In the small-Bach case the gauge invariance
and \(C^4\) convergence directly give \(\int |B_G|_G\,dV_G=0\).

Lemma~\ref{var:continuation} applies to \(G\). Signature on the
smooth \(G_j\) gives \(E_-(G_j)=E_+(G_j)\to0\). Passing to the limit
gives \(W^-=0\) as well. To justify smooth recognition at the
finite metric regularity, the contracted Bianchi identity gives
zero Cotton tensor and hence the standard real tractor
connection is flat. This connection has \(C^{1,\beta}\)
coefficients for any \(0<\beta<1\) at the stated grade, so its
parallel frames are \(C^{2,\beta}\). They give a global
parallel trivialization by simple connectivity. In this
trivialization the canonical null ray develops into the
projective light cone \(S^4\). If \(v\) is the null lift in
the scale \(G\), its differential is the middle slot
\(X\mapsto(0,G(X,\cdot),0)\) transported by the parallel frame.
Normalizing by \(a=-\langle v,t\rangle>0\) for a fixed unit
timelike \(t\) gives a \(C^{3,\beta}\) map whose pullback round
metric is \(a^{-2}G\)~\cite[\S3.1, (10) and (13)]{GoverSilhan2009HigherSymmetries}.
It is a local diffeomorphism and hence a covering by compactness.
The target is simply connected, so it has one sheet. Approximate
the map in \(C^1\) by smooth maps into \(\mathbb R^5\) and
radially project onto \(S^4\). A sufficiently close approximation
still has invertible derivative everywhere and is again a
one-sheeted covering. This yields a smooth diffeomorphism on the
original smooth \(M\).
\end{proof}

The common open set supplies the initial vanishing for analytic
continuation. A compactness result for unmarked metrics alone may
lose that set. The argument also needs the positive lower metric
bound and the strong convergence that passes the Bach equation.

\begin{proposition}[Completion at zero chiral energy]\label{var:zeroenergy}
Let \(g_j\) be smooth metrics on a closed oriented homotopy four-sphere
with \(E_+(g_j)\to0\). Suppose that after orientation-preserving
pullbacks they have uniformly bounded \(C^{3,\alpha}\) coefficients in
a fixed finite smooth atlas, \(0<\alpha<1\), and satisfy
\(dV_{g_j}\ge c\,dV_r\) for a fixed smooth metric \(r\) and \(c>0\).
Then the manifold is diffeomorphic to \(S^4\) and
\(Y([g_j])>0\) for all sufficiently large \(j\).
\end{proposition}
\begin{proof}
Indeed the upper eigenvalue bound and the determinant
lower bound give uniform ellipticity. Compact embedding gives a
positive \(C^{3,\beta}\) subsequential limit \(G\) for
\(0<\beta<\alpha\). Curvature convergence and \eqref{eq:sign}
give \(W_G^+=W_G^-=0\). The development argument in Theorem~\ref{var:compactness} applies at
this lower regularity. The differential Bianchi identity gives zero
Cotton tensor, so the normal tractor connection is flat and still has
\(C^{1,\beta}\) coefficients. Its parallel frame develops the canonical
null ray into a conformal local diffeomorphism, and the same covering
and smoothing argument gives smooth recognition. Moreover, write
\(u^2G=\operatorname{dev}^*g_{\rm rd}\). The scalar transformation
law gives \((6\Delta_G+s_G)u=12u^3\), hence
\[
 Q_G(f)=6\int_M u^2|d(f/u)|_G^2\,dV_G+
             12\int_M u^2 f^2\,dV_G.
\]
The positive factor \(u\) and compactness make this coercive on
\(H^1\). Coercivity persists under \(C^2\) convergence, and Sobolev
embedding yields positive Yamabe sign for nearby smooth metrics.
Applying the argument to any purported nonpositive subsequence gives
eventual positivity for the full sequence. This criterion uses a
pointwise lower volume-density bound in the coordinates that control
the derivatives.
\end{proof}

In particular a radial tail with \(E_+\to0\) and negative Yamabe sign
cannot satisfy the fixed-atlas hypotheses of Proposition~\ref{var:zeroenergy}
along any subsequence. In a pointwise fixed-volume chart a uniform
\(C^{3,\alpha}\) bound alone would supply the density hypothesis.

\subsection{The scalar-flat threshold and two conformal gauges}

Put \(E=E_+(g)\) and \(T_g=\operatorname{Ric}_g-R_g g/4\). The
Yamabe theorem supplies a smooth unit-volume minimizing
representative \(g_Y\), with \(R_{g_Y}=Y([g])\)~\cite[\S3, Theorem 3]{Schoen1984Yamabe}.
Chern--Gauss--Bonnet and signature yield
\begin{equation}\label{var:wall}
 \begin{aligned}
 16\pi^2&=2E-\tfrac12\|T_{g_Y}\|_2^2+\tfrac1{24}Y([g])^2,\\
 Y([g])^2&=48(8\pi^2-E)+12\|T_{g_Y}\|_2^2 .
 \end{aligned}
\end{equation}
Here the Euler integrand uses one quarter of the full
four-index Weyl norm squared~\cite[p.~107, (0.4) and Remark 2]{ChangGurskyYang2003}.
At \(Y=0\) this gives \(E\ge8\pi^2\), with equality precisely
when the class has a Ricci-flat representative. If the class
admits the marked round cap, equality would give a Bach-flat
marked metric by conformal covariance, hence a conformally
round extension by Proposition~\ref{var:stationarity}. Its
energy would be zero. Thus \(E>8\pi^2\) for the marked
Yamabe-zero problem.

The lowest eigenvalue of \(6\Delta_g+R_g\), with
\(\Delta_g=-\operatorname{div}_g\nabla\), has the sign of \(Y\).
For a positive unit-volume Yamabe factor \(u\), write a test
function as \(u f\) and integrate by parts using
\((6\Delta_g+R_g)u=Y u^3\). The numerator becomes
\(\int(6u^2|df|_g^2+Y u^4 f^2)\,dV_g\). This gives strict
positivity for \(Y>0\), and testing \(u\) treats \(Y\le0\).
The eigenvalue is continuous on a \(C^2\)-continuous path of
positive metrics. Locally in the parameter, compare the
quadratic forms after a fixed coercive shift and compare their
\(L^2\) denominators on the common \(H^1\) domain. The relative
errors tend to zero by uniform convergence of the inverse
metrics, volume densities and scalar curvatures. Infimizing the
shifted Rayleigh quotients gives continuity. It follows from
\eqref{var:wall} that the sign is constant on each connected
portion with \(E<8\pi^2\), where \(|Y|\ge\sqrt{48(8\pi^2-E)}\).
A positive sign known at one point
of that portion can be used for recognition at any time where
the strict \(E<4\pi^2\) criterion is met.

\begin{proposition}[Estimates for a supplied subcritical positive class]\label{var:gauges}
Suppose \(Y([g])>0\) and \(E\le8\pi^2-\delta_0\) for \(\delta_0>0\).
In the unit-volume Yamabe representative the following bounds hold,
with the Sobolev inequality valid for all \(v\in H^1(M,g_Y)\).
\begin{equation}\label{var:sobolev}
 \sqrt{48\delta_0}\le Y\le8\sqrt6\,\pi,\qquad
 \|v\|_4^2\le \frac6{\sqrt{48\delta_0}}\|dv\|_2^2+\|v\|_2^2 .
\end{equation}
There is also a unit-volume representative \(g_\sigma\) with
constant \(q=\sigma_2(A)=(8\pi^2-E)/2\), where
\(A=(\operatorname{Ric}-R g/6)/2\) and the symmetric functions
are those of \(g_\sigma^{-1}A\). In this representative the
least Ricci eigenvalue is at least \(2q/\sigma_1(A)>0\).
If in addition \(R_{g_\sigma}\le\Lambda\), then
\begin{equation}\label{var:diameter}
 \begin{gathered}
 \operatorname{Ric}_{g_\sigma}\ge(6\delta_0/\Lambda)g_\sigma,
 \qquad \operatorname{diam}(M,g_\sigma)\le
 D:=\pi\sqrt{\Lambda/(2\delta_0)},\\
 \operatorname{Vol}_{g_\sigma}B(x,\rho)\ge(\rho/D)^4
 \quad(0<\rho\le D).
 \end{gathered}
\end{equation}
\end{proposition}

\begin{proof}
Equation~\eqref{var:wall} gives the lower bound on \(Y\).
Aubin's upper bound is the round Yamabe value
\(12\operatorname{Vol}(S^4,g_{\rm rd})^{1/2}=8\sqrt6\,\pi\).
The variational inequality
\(Y\|v\|_4^2\le6\|dv\|_2^2+Y\|v\|_2^2\)
in \(g_Y\) proves \eqref{var:sobolev}.
In this gauge \eqref{var:wall} also gives \(\|T\|_2^2\le4E\).
The curvature decomposition then yields
\(\|\operatorname{Ric}\|_2^2\le4E+96\pi^2\) and
\(\|\operatorname{Rm}\|_{L^2,(0,4)}^2\le16E+64\pi^2\), using
\(|\operatorname{Rm}|_{(0,4)}^2=|W|_{(0,4)}^2+2|T|^2+R^2/6\)~\cite[\S\S2--3]{LeeParker1987Yamabe}.

Chang--Gursky--Yang use the Schouten normalization \(2A\).
The integral of \(\sigma_2(2A)\) is \(16\pi^2-2E>0\) by
Chern--Gauss--Bonnet. Their Theorem 1.2 applies with this positive
integral and the given positive Yamabe sign. It yields a smooth
conformal metric with \(\sigma_2(2A)\) positive constant.
After volume normalization this gives the stated \(q\), since
\(\sigma_2(2A)=4\sigma_2(A)\)~\cite[\S1, (1.1) and Theorem 1.2]{ChangGurskyYang2003}.

In \(g_\sigma\) write \(S=\sigma_1(A)=R/6\). The identity
\(2q=S^2-|A|^2>0\) gives a nonvanishing \(S\) of constant sign.
The constant test function in the Yamabe quotient rules out
the negative sign. If \(a\) is the least Schouten eigenvalue,
then the other three eigenvalues have pair sum at most
\((S-a)^2/3\). The least Ricci eigenvalue is \(m=S+2a\), and
\[
 q\le a(S-a)+(S-a)^2/3=\frac{m(3S-m)}6 .
\]
Since \(S-a>0\), the bound gives \(m>0\). In fact
\(m\ge(3S-\sqrt{9S^2-24q})/2\ge2q/S\).
Using \(q\ge\delta_0/2\) and \(S\le\Lambda/6\) gives the Ricci
bound in \eqref{var:diameter}. Bonnet--Myers gives \(D\), and
the Bishop--Gromov relative volume comparison with an outer
radius larger than \(D\) gives the ball bound by letting the
outer radius decrease to \(D\). This uses completeness and unit
volume of \(g_\sigma\).
\end{proof}

The scalar upper bound in \eqref{var:diameter} is additional
input in \(g_\sigma\). At fixed \(q>0\), the three equal larger
Schouten eigenvalues can saturate the quadratic bound on \(m\)
while \(S\) tends to infinity and \(m\) tends to zero. This is
pointwise algebra. In particular the constant scalar value \(Y\)
of \(g_Y\) provides no such control in \(g_\sigma\) without
estimates on the conformal change.

\section{Sharpness of the Yamabe-zero threshold}\label{sec:wall}

Fix a smoothly parametrized closed coordinate ball
\(K=\psi(\overline B_1)\), where \(\psi:B_{1+d}\to S^4\) is an
orientation-preserving chart and \(d>0\). The round marking is the
requirement \(\psi^*g=(2/(1+|x|^2))^2\delta\) on \(B_1\).
Write \(\mathfrak B(g)=\int |B_g|_g\,dV_g\) for the total Bach norm.

\begin{theorem}[Sharp marked Yamabe-zero infimum]\label{wall:main}
On the standard smooth \(S^4\), the infimum of \(E_+\) among
conformal classes with Yamabe constant zero admitting the round
marking on \(K\) is \(8\pi^2\). The infimum is unattained. There are
smooth marked metrics \(g_R\), \(R\to\infty\), of Yamabe constant
zero such that
\[
 E_+(g_R)=E_-(g_R)=8\pi^2+O(R^{-1}),\qquad
 \mathfrak B(g_R)=O(R^{-1}).
\]
Their smooth unit-volume scalar-flat representatives \(k_R\)
satisfy \(\int|\operatorname{Ric}_{k_R}|^2\,dV_{k_R}\to0\).
The unmarked infimum on the standard sphere is also \(8\pi^2\),
and its attainment is equivalent to the existence of a smooth
Ricci-flat metric there.
\end{theorem}

\subsection{Taub--NUT cores and their chiral energy}

On the one-centre Taub--NUT space identify the complement of the
nut with \((0,\infty)\times S^3\), using \(r=|z|^2/2\) in the
underlying smooth coordinates \(z\in\mathbb C^2\). Write \(\pi\) for
the Hopf map to the unit round sphere \((S^2,h_0)\). The connection
\(\eta\) is normalized by circle period \(2\pi\) and
\(d\eta=-\pi^*dA_0/2\) in the chosen base orientation. We use the
normalization
\begin{equation}\label{wall:tn}
 g_{\rm TN}=V(dr^2+r^2\pi^*h_0)+V^{-1}\eta^2,
 \qquad V=1+\frac{1}{2r}.
\end{equation}
For example, with \(\rho=|z|\) the Euclidean metric is
\(d\rho^2+\rho^2(\pi^*h_0/4+\eta^2)\), and
\(\alpha=\rho^2\eta\) is smooth at the nut. Indeed, up to sign it is
\(\operatorname{Im}(\bar z_1\,dz_1-\bar z_2\,dz_2)\).
Equation \eqref{wall:tn} becomes
\((1+\rho^2)\delta-(2+\rho^2)\alpha^2/(1+\rho^2)\), showing smooth
positive extension at \(z=0\). In the hyperkähler orientation
the metric is Ricci-flat and the three parallel Kähler forms
trivialize the self-dual connection, so \(W^+=0\)~\cite[\S2.1, (1)]{CherkisHurtubise2019Monads}.

In outward radial arc distance \(d\ell=\sqrt V\,dr\), a slice has
horizontal warping factor \(a=r\sqrt V\) and vertical factor
\(c=V^{-1/2}\). The radial curvature matrix is diagonal in
these distributions, with equal horizontal entries and vertical
entry \(K=-c''/c\). Here the prime denotes \(d/d\ell\), and the
normal Riccati equation gives
\[
 c'=(1+2r)^{-2},\qquad K=4(1+2r)^{-3}.
\]
Projection of the three radial unit bivectors onto \(\Lambda^-\)
gives orthogonal vectors of norm \(1/\sqrt2\). All curvature is in
this block, so its entries in the normalized basis are twice the
radial ones. Trace-freeness therefore gives the three eigenvalues
\((-K,-K,2K)\). The shell volume is \(8\pi^2 Vr^2\,dr\), by the
circle period and the area \(4\pi\) of the Hopf base. Thus
\begin{equation}\label{wall:core-energy}
 A_R:=E_-(\{r\le R\},g_{\rm TN})
 =96\pi^2\int_0^{2R}\frac{x}{(1+x)^5}\,dx
 =8\pi^2\left(1-\frac4{(1+2R)^3}+\frac3{(1+2R)^4}\right).
\end{equation}
This uses the Hilbert--Schmidt norm on the curvature operator.
It gives \(A_{1/2}=11\pi^2/2\) and agrees with the complete
Taub--NUT integral \(8\pi^2\) in this norm after orientation
translation~\cite[\S2 and \S5, (5.11), (5.18)]{AtiyahLeBrun2013Curvature}.

Two copies of \(r\le R\) are standard coordinate four-balls.
Join their boundaries through an \(S^3\)-cylinder by the identity
in the markings \(z/\sqrt{2R}\). This is the standard smooth
double, as is seen by using two latitude caps of the standard
sphere and their product complement. The two outward radial
directions point oppositely into the cylinder, so the global
orientation agrees with the hyperkähler orientation on one
ball and reverses it on the other. For all \(R\ge1\) the balls
can be identified with fixed \(R=1\) balls without moving the
inner balls \(r\le1/2\). For instance, send the reference \(r\) to
\(r[1+(R-1)\zeta(r)]\), where \(\zeta\) is a nondecreasing smooth
cutoff vanishing near the inner ball and equal to one near
the reference boundary. This also preserves the seam coordinate
\(u=r/R\). Before any small marking modification, the retained
inner balls give a lower bound \(11\pi^2/2\) in each chirality.
More generally, retaining even one such conformal ball forces
this lower bound on both global energies on a homotopy sphere,
by the signature identity. In the ASD orientation \(W^+\) vanishes
on the ball, so the exterior alone carries at least \(11\pi^2/2\) of
positive chiral energy. Any extension retaining this finite conformal
ball therefore exceeds the \(4\pi^2\) energy threshold of
Proposition~\ref{prop:recognition}.

Such a retained ball exists as a metric prescription on any smooth
\(M\): embed a slightly larger coordinate ball and interpolate the
Taub--NUT metric with any background metric by a smooth convex cutoff
equal to one on the smaller ball. If its exterior positive chiral energy
is supported in a region \(T\), then
\(\Vol_g(T)\sup_T |W^+|_{\HS}^2\ge11\pi^2/2\). Confining that energy
to regions of shrinking physical volume forces the supremum to diverge.

\subsection{A joining metric of each Yamabe sign}

Work first in the metric scaled by \(R^{-2}\). Set \(V_R(u)=1+(2Ru)^{-1}\) and introduce the signed distance
\(s=D_R(u)=\int_1^u\sqrt{V_R(v)}\,dv\). Then
\[
 R^{-2}g_{\rm TN}=ds^2+\gamma_s,\qquad
 \gamma_s=a^2\pi^*h_0+c^2\eta^2,\qquad
 a^2=u^2V_R,\quad c^2=R^{-2}V_R^{-1}.
\]
These formulas extend smoothly to a fixed collar on either side
of \(u=1\). In a slice the curvature of a horizontal plane is
\(a^{-2}-3v^2/4\), and that of either mixed plane is \(v^2/4\),
where \(|v|=c/(2a^2)\) is the length of the vertical bracket of
an orthonormal horizontal pair. In fact the fibres are geodesic
and the vertical bracket is \(-d\eta(E_1,E_2)c\) times the unit
vertical vector. The constant-fibre submersion formulas apply
also when the base metric is distorted below~\cite[\S3, proof of Lemma 3.2]{ShankarTappTuschmann2005Circle}.
In particular the intrinsic scalar curvature and outward mean
curvature of the round-base slices are
\begin{equation}\label{wall:slice}
 S(\gamma_s)=\frac{2}{a^2}-\frac{c^2}{8a^4},\qquad
 H_s=\partial_s\log(a^2c)
 =V_R^{-1/2}\left(\frac2u-\frac{1}{4Ru^2V_R}\right).
\end{equation}
Both are uniformly positive and bounded on \(1\le u\le2\), \(R\ge1\).

For the positive endpoint, choose on \([-1,1]\) a fixed smooth
concave function \(\sigma\) that is \(t+1\) near the left end,
\(1-t\) near the right end and constant for \(|t|\le1/4\).
For example, integrate minus a smooth odd nondecreasing function
that is zero on \([-1/4,1/4]\) and equals \(\pm1\) near the respective
ends, with \(\sigma(-1)=0\). Then \(0\le\sigma\le1\),
\(|\sigma'|\le1\) and \(u_R(\sigma)\in[1,2]\), where \(u_R=D_R^{-1}\).
Join by \(dt^2+\gamma_{\sigma(t)}\). The seam coordinates are
exactly \(s=t+1\) and \(s=1-t\), so the metric agrees with the
Taub--NUT continuation on open collars. Its longitudinal
coefficient is one, including on the plateau.

For a unit radial metric \(ds^2+\gamma_s\), the Gauss formula
for the tangential scalar trace and the Riccati trace in the
normal direction give \(S=S(\gamma_s)-|A_s|^2-H_s^2-2\partial_s H_s\),
where \(A_s\) is the second fundamental form. The same formula
on the join uses \(\sigma'A_{\sigma}\) and \(\sigma'H_\sigma\).
Since the original scalar curvature vanishes, substitution yields
\[
 S(dt^2+\gamma_{\sigma(t)})
   =(1-\sigma'^2)S(\gamma_{\sigma(t)})-2\sigma''H_{\sigma(t)} .
\]
This is nonnegative and strictly positive on the plateau~\cite[Lemma 2.4]{EbertFrenck2021Surgery}.
The closed joined metric has positive Yamabe constant. Indeed
its conformal Laplacian has positive lowest \(L^2\) eigenvalue,
as a zero Rayleigh minimizer would have to be constant and
vanish on the positive-scalar-curvature region. This eigenvalue
and the compact Sobolev inequality give coercivity against the
critical \(L^4\) norm. A fixed smooth axial test supported in
the plateau has Yamabe numerator \(O(R^{-1})\) and denominator
comparable to \(R^{-1/2}\), since the slice volume there is
\(8\pi^2 a^2c=8\pi^2u^2\sqrt{V_R}/R\). Thus this positive
Yamabe constant is at most \(C R^{-1/2}\).

For the negative endpoint, let \(f=2+x_3\) on the unit round
base, \(\Delta_0=\operatorname{div}_{h_0}\nabla\), and
\[
 v=k\Delta_0(f^2),\qquad h=e^{2v}h_0.
\]
In dimension two the Dirichlet integral is conformally invariant
and \(S_h\,dA_h=(2-2\Delta_0v)dA_0\). It follows that
\[
 I_h=\int_{S^2}(6|df|_h^2+S_h f^2)\,dA_h
 =\int_{S^2}(6|df|_0^2+2f^2)\,dA_0
       -2k\int_{S^2}(\Delta_0(f^2))^2\,dA_0 .
\]
The last integral is positive because \(f^2\) is nonconstant.
Fix \(k\) so \(I_h=-m<0\). Choose a nonzero smooth nonnegative
bump \(w(t)=w_0(t/L)\) supported in \((-L,L)\) with \(L\) a
large fixed number. With \(J=\int_{S^2}f^2\,dA_h\), arrange
\(6J\int|w'|^2\le(m/2)\int w^2\), since the ratio scales as \(L^{-2}\).
On \([-L,L]\times S^3\) use
\(dt^2+\pi^*h+R^{-2}\eta^2\). Its scalar curvature is
\(S_h-|F|_{\Lambda^2 h}^2/(2R^2)\), where \(d\eta=\pi^*F\), by the
horizontal and mixed formulas above. Its volume density is
\(R^{-1}dt\,|\pi^*dA_h\wedge\eta|\). The separated test
\(w(t)\pi^*f\), extended by zero, therefore has Yamabe numerator
\[
 \frac{2\pi}{R}\left[
  6J\int |w'|^2+
   \left(I_h-\frac1{2R^2}\int_{S^2}|F|_{\Lambda^2 h}^2 f^2\,dA_h\right)
       \int w^2\right]
 \le -\frac{\pi m}{R}\int w^2 .
\]
The denominator is a positive constant times \(R^{-1/2}\).
Hence the negative endpoint has \(Y\le-c_0R^{-1/2}<0\).

\subsection{Common collars and a uniformly controlled interpolation}

Take the common neck \([-T,T]\times S^3\), \(T=L+2\), with
seam coordinates \(u_-=1+t+T\) on the left and \(u_+=1+T-t\)
on the right. On small fixed end collars use exactly the continuation
\begin{equation}\label{wall:seams}
 S_{R,\pm}=V_R(u_\pm)(dt^2+u_\pm^2\pi^*h_0)
             +R^{-2}V_R(u_\pm)^{-1}\eta^2.
\end{equation}
On the negative neck extend from these collars to the fixed
distorted product by a convex combination with disjoint smooth
end cutoffs, each equal to one near its own seam and zero on
\([-L,L]\). This gives agreement of the entire seam germs.
Pull the positive endpoint back from \([-1,1]\) by an increasing
map \(b_R:[-T,T]\to[-1,1]\) that equals
\(-1+D_R(u_-)\) and \(1-D_R(u_+)\) on the respective collars.
For a precise extension, let \(d_\pm=\sqrt{V_R(u_\pm)}\) there and
use fixed disjoint end cutoffs \(\beta_\pm\) on sufficiently
short supports. Set
\[
 b_R'=\beta_-d_-+\beta_+d_++\gamma_R(1-\beta_--\beta_+),
 \qquad
 \gamma_R=\frac{2-\int(\beta_-d_-+\beta_+d_+)}
                 {\int(1-\beta_--\beta_+)}.
\]
The collars can be chosen so the subtracted integral is below one
for all \(R\ge1\). Integrate with \(b_R(-T)=-1\). This gives the
prescribed end germs, with upper and positive lower derivative
bounds and bounds through order five for the map. The pulled-back
positive metric therefore agrees with \eqref{wall:seams} on the same
end collars as the negative one. Both are on the same smooth double
via the identity in \(u\) on the cores and the sphere marking.

Convexly interpolate the full positive metric forms on the neck,
leaving both scaled cores fixed. In a finite family of fixed base
charts write \(\eta=d\phi+\alpha\) and locally unwrap the fibre using
\(z=\phi/R\). Then \(R^{-1}\eta=dz+R^{-1}\alpha\). The coefficient
matrices of both endpoints have uniform positive lower and upper
bounds and four bounded coordinate derivatives in these charts.
For the positive endpoint this follows from \(du/ds=V_R^{-1/2}\),
the fixed smooth slice profile, and the five derivative bounds
for \(b_R\). For the negative endpoint it follows from the fixed
base metrics and cutoff functions and the bounded shell
coefficients in \eqref{wall:seams}. Matrix convexity keeps
uniform ellipticity along the interpolation, and differentiation
of the inverse gives the corresponding inverse bounds. Hence
curvature and its first two covariant derivatives are bounded
on the scaled neck uniformly in \(R\) and in the interpolation
parameter. This uses only local isometric lifts. The determinant
is bounded on the lifts, and their quotient period \(2\pi/R\)
in \(z\) makes the scaled neck volume \(O(R^{-1})\).

The Bach tensor is controlled by the second derivatives of
Schouten and the Schouten--Weyl product, and is zero on each
Ricci-flat core~\cite[\S1, (1.1)]{BahuaudGuentherIsenbergMazzeo2025Bach}.
The neck thus contributes \(O(R^{-1})\) to \(\mathfrak B\) and to
each chiral energy. A constant homothety restoring \(g_{\rm TN}\)
on the cores leaves these integrals and the Yamabe constant
invariant. In dimension four a positive conformal factor \(u^2\)
gives weight \(-2\) to the covariant Bach tensor, hence weight
\(-4\) to its norm, which cancels the volume weight. The Weyl
operator has weight \(-2\), giving the same cancellation for
its squared norm. We conclude that for the restored path
\(g^0_{R,\theta}\), \(0\le\theta\le1\),
\begin{equation}\label{wall:path}
 E_\pm(g^0_{R,\theta})=A_R+O(R^{-1}),\qquad
 \mathfrak B(g^0_{R,\theta})=O(R^{-1})
\end{equation}
uniformly in \(\theta\). The chiral energies are equal by
signature for every smooth metric on \(S^4\).

\subsection{Marking, the sign crossing and nonattainment}

Place the compact closure of a slightly shrunken extended
marked chart inside a fixed interior region of one unscaled
core. To do this, compress the chart by the flow of a cut off
radial contraction, then move its center to the core by a fixed
orientation-preserving round rotation. Pull the metrics back
by this single placement map. In oriented normal coordinates
\(x\) at the marked center the entire path has the same
Taub--NUT germ \(\delta+q(x)\), with \(q=O(|x|^2)\) and
\(Dq=O(|x|)\). Replace it by
\(\delta+\chi(x/\varepsilon)q(x)\), where \(\chi=0\) on \(B_1\)
and \(\chi=1\) outside \(B_2\). The positive metric is Euclidean
on \(B_\varepsilon\), and the derivative bounds through order
four on \(B_{2\varepsilon}\) are \(O(\varepsilon^{2-j})\) at
order \(j\). Curvature is bounded there and its second covariant
derivative is \(O(\varepsilon^{-2})\). The changed volume is
\(O(\varepsilon^4)\), so the energy costs are \(O(\varepsilon^4)\)
and the total Bach cost is \(O(\varepsilon^2)\), with constants
from the one common germ.

For the positive endpoint \(h=g^0_{R,0}\) before flattening,
write \(Q_h(\varphi)=\int(6|d\varphi|_h^2+S_h\varphi^2)\,dV_h\)
and \(D_h(\varphi)=\|\varphi\|_{L^4(h)}^2\).
The replacement \(h_\varepsilon\) differs by \(O(\varepsilon^2)\)
in metric coefficients and by \(O(\varepsilon^2)\) in the
\(L^2\) norm of the scalar density relative to \(dV_h\).
Indeed the scalar curvature is bounded on the changed region
and unchanged elsewhere. Uniformly over \(\varphi\), comparison
of inverse metrics and volume forms and Hölder's inequality give
\[
 |Q_{h_\varepsilon}-Q_h|\le C_R\varepsilon^2(Q_h+D_h),\qquad
 D_{h_\varepsilon}\le(1+C_R\varepsilon^2)D_h .
\]
Here \(S_h\ge0\) makes \(Q_h\) dominate the Dirichlet term.
As \(Q_h\ge Y([h])D_h\), choose one sufficiently small
\(\varepsilon(R)\le \varepsilon_*/R\) so
\(C_R\varepsilon(R)^2(1+Y([h]))\le Y([h])/2\). This preserves
positive Yamabe sign. The negative test is supported off the
change, so the same radius preserves its negative sign.

The negative endpoint has the two-sided scale
\(-C R^{-1/2}\le Y([g_R^-])\le-cR^{-1/2}\). For the
lower bound the scalar curvature vanishes on the retained cores
and is uniformly bounded on the scaled neck of volume \(O(R^{-1})\).
H\"older's inequality bounds the Yamabe quotient below by
\(-\|(\Scal)_-\|_2\). The cap replacement contributes at most
\(O(\varepsilon(R)^2)\) to this norm, since its curvature is bounded
on a set of volume \(O(\varepsilon(R)^4)\). The separated test gives
the other inequality. Both bounds pass to the final marked class.

The same local construction applies to an arbitrary fixed smooth metric
germ. In normal coordinates a flat replacement matches its zero and first
jets. A cutoff on a ball of radius \(\varepsilon\) bounds the curvature
by \(O(1)\) and the Bach tensor by \(O(\varepsilon^{-2})\) there. The ball
volume is \(O(\varepsilon^4)\). Thus the changes in \(E_\pm\) and
\(\mathfrak B\) are \(O(\varepsilon^4)\) and \(O(\varepsilon^2)\),
respectively. The reparametrization and conformal marking below can be
supported in the extended chart, so the original exterior metric is preserved.

The shared flat germ gives the marking for the whole path by
one diffeomorphism and one positive conformal factor at each
\(R\). To see this, the oriented transition germ from the
marked coordinates to the flat coordinates can be extended by a
localized flow. Join its derivative to the identity in
\(\mathrm{GL}^+(4,\mathbb R)\), then join the derivative to the
germ \(G\) by the maps \(x\mapsto \lambda^{-1}G(\lambda x)\)
for \(0\le\lambda\le1\), extended smoothly at zero. Cut off the
velocity of this local isotopy on a neighbourhood in the extended
chart. By first compressing the marked ball to this neighbourhood,
one obtains a pullback equal to \(a^2\delta\) in the marked
coordinates for some \(a>0\). Multiplication there by
\((2/(a(1+|x|^2)))^2\), extended as a smooth positive factor
using a cutoff of its logarithm, supplies the exact round metric.
All marking operations are supported in the extended chart
away from the neck. They preserve the conformal invariants.
The flattening and the marking produce a smooth marked path
\(g_{R,\theta}\), with the bounds \eqref{wall:path} and opposite
Yamabe signs at its endpoints.

For fixed \(R\), the lowest eigenvalue of
\(6(-\operatorname{div}_g\nabla)+S_g\) varies continuously along
the path. On the common \(H^1\) domain, shift the Rayleigh
numerator at a reference parameter by a large fixed multiple
of the \(L^2\) denominator to obtain a coercive form.
Smooth metric variation makes the shifted forms and \(L^2\)
norms relatively close, uniformly over test functions.
Taking infima of the shifted quotients gives continuity.
The first eigenvalue has the same sign as \(Y([g])\).
With a positive smooth unit-volume Yamabe factor \(a\), the
equation \((6(-\operatorname{div}\nabla)+S_g)a=Y([g])a^3\)
and integration by parts give
\[
 Q_g(\varphi)=6\int a^2|d(\varphi/a)|_g^2\,dV_g
                  +Y([g])\int a^2\varphi^2\,dV_g.
\]
This is positive in \(L^2\) for positive \(Y\). At zero the
test \(a\) gives the null eigenvalue, and for negative \(Y\)
it gives a negative test. Thus a parameter between the
endpoints has \(Y=0\). Choose one for each \(R\) to obtain
the sequence of marked metrics~\cite[\S3, Theorem 3]{Schoen1984Yamabe}
\cite[\S1]{LeeParker1987Yamabe}.

Finally, on any smooth homotopy four-sphere, choose the unit-volume
scalar-flat representative \(k\) in a Yamabe-zero class.
Chern--Gauss--Bonnet and zero signature give
\begin{equation}\label{wall:cgb}
 16\pi^2=2E_+([k])-\tfrac12\int|\operatorname{Ric}_k|_k^2\,dV_k,
 \qquad
 \int|\operatorname{Ric}_k|_k^2\,dV_k=4(E_+([k])-8\pi^2).
\end{equation}
Here the Euler integrand uses one quarter of the full four-index
Weyl norm, and that full squared integral equals
\(4(E_++E_-)=8E_+\). This proves the lower bound, and
\eqref{wall:path} at the sign crossings proves sharpness and
the Ricci-integral limit. Equality is equivalent to a Ricci-flat
representative, whose Yamabe numerator is the Dirichlet form.
If the class is marked, this representative is Bach-flat and
conformally flat on the open cap. The four-dimensional Bach
equation is the ambient obstruction equation. Analytic
continuation of Weyl-flatness for smooth obstruction-flat metrics
then gives full Weyl vanishing on the connected manifold,
contradicting the energy at equality~\cite[(0.4), Remark 2]{ChangGurskyYang2003}
\cite{GrahamHirachi2004Ambient}
\cite[Theorem 1.4]{LassasLiimatainen2019ConformalHarmonic}.
The scalar-flat representative used in \eqref{wall:cgb} is
conformally flat on the cap. The exact unit round marking, with
scalar curvature \(12\) there, belongs to a different representative.

The lower bound holds also without the cap and the constructed
sequence belongs to both classes on the standard smooth sphere.
Unmarked attainment therefore reduces to the Ricci-flat existence
question. The strict inequality for each marked Yamabe-zero class
holds on an arbitrary smooth homotopy sphere, while the construction
of a sharp sequence here uses the standard smooth sphere.

For the zero-sign sequence the cap-surviving compactness
hypotheses of Theorem~\ref{var:compactness} fail in every system of
conformal gauges on \(S^4\). Otherwise its vanishing total Bach norm
and a surviving common open part of the cap would force \(E_+\to0\),
contrary to \(E_+\to8\pi^2\). This applies also to its scalar-flat
representatives, since their conformal factors can be absorbed into
the comparison gauges.

\section{Radial strain at the Yamabe threshold}\label{mscost:section}

Surjectivity of the full Weyl linearization allows a local radial
prescription at each regular reference. The initial metric variation
required by this particular source can grow without bound. We obtain
an estimate in the reference metric and apply it to the marked
Taub--NUT doubles from Section~\ref{sec:wall}.

\subsection{The Bach pairing with the radial source}

Let \(M\) be a smooth closed oriented homotopy four-sphere and let
\(g\) be a smooth Riemannian metric. Use \(L_g=dF_g(0)\) on the full
trace-free domain of the fixed-volume chart \eqref{lift:chart}.
The comparison is an isometry on the Weyl fibers and preserves the
volume form, so
\[
 dE_+(g)[h]
 =2\left\langle W_g^+,L_gh\right\rangle_{L^2(g)}
 =-\tfrac12\int_M\langle B_g,h\rangle_g\,dV_g ,
\]
where the second equality is \eqref{var:bach}. In particular
\(L_g^*W_g^+=-B_g/4\) in the chiral operator normalization.
For every solution of \(L_gh=-W_g^+\) this gives
\begin{equation}\label{mscost:strain}
  4E_+(g)=\int_M\langle B_g,h\rangle_g\,dV_g
     \le \mathfrak B(g)\,\|h\|_{L^\infty(g)} .
\end{equation}
The identity extends from smooth variations to the Sobolev domain
of \(L_g\) at the grades in Section~\ref{sec:lift} by density and
Sobolev embedding.
If \(F_g(h(t))=(1-t)W_g^+\) with \(h(0)=0\), its derivative
\(h'(0)\) solves this equation whenever the lift is differentiable
in the chart. At a regular reference the normal slice gives such
an exact local lift by Proposition~\ref{lift:radial}. The estimate
applies to every choice of initial strain, including those obtained
by adding a homogeneous solution to a particular preimage.

\subsection{Regular references with Yamabe constant zero}

\begin{proposition}\label{mscost:zero}
Let \(r\) be a smooth Riemannian metric with \(Y([r])=0\) on a smooth
closed oriented homotopy four-sphere. Suppose \(r\) is round on a
marked closed coordinate ball \(K\). There are regular smooth
references arbitrarily close to \(r\) in \(C^\infty\), with the same
pointwise volume form and the same metric on a neighborhood of
\(K\), whose Yamabe constant is exactly zero.
\end{proposition}

\begin{proof}
Set \(\nu=dV_r\), \(H_r=6\Delta_r+\Scal_r\), where
\(\Delta_r=-\operatorname{div}_r\nabla\), and write \(\lambda_1(r)\)
for its first \(L^2\) eigenvalue. The sign comparison from
Section~\ref{sec:wall} gives \(\lambda_1(r)=0\). Write \(u>0\)
for a zero eigenfunction with \(\int_M u^2\nu=1\).
The ground state identity
\[
 \langle uf,H_r(uf)\rangle_{L^2(r)}
     =6\int_M u^2|df|_r^2\,\nu
\]
shows that the kernel is simple by connectedness. More generally,
subtracting the eigenvalue term in this identity shows that a
positive eigenfunction realizes the lowest eigenvalue.

Along a smooth finite-parameter family \(r_\xi\) with the same
density, the normalized first eigenpair is smooth near \(r\).
Indeed the eigenpair and normalization rows are
\((H_{r_\xi}f-\lambda f,\int f^2\nu-1)\).
At \((r,u,0)\) their derivative in \((f,\lambda)\) is
\[
 (\dot f,\dot\lambda)\longmapsto
       (H_r\dot f-\dot\lambda u,\;2\langle u,\dot f\rangle_{L^2(r)}).
\]
The component along \(u\) in the first row determines
\(\dot\lambda\). On \(u^\perp\) the elliptic inverse of \(H_r\)
determines \(\dot f\), with its component along \(u\) determined by
the second row. This is an isomorphism from \(H^{s+2}\times\mathbb R\)
to \(H^s\times\mathbb R\), for example for any fixed integer
\(s\ge5\). The implicit function theorem gives the eigenpair, and
Sobolev embedding keeps its eigenfunction positive. Elliptic
regularity gives spatial smoothness~\cite[(1)--(2)]{Melrose2017Lecture22}.

Put \(k=u^2r\), which is scalar-flat by the conformal scalar
curvature law~\cite[(2.6)--(2.7)]{LeeParker1987Yamabe}.
Differentiating the eigenvalue equation and testing against \(u\)
shows that \(d\lambda_1(r)[h]\) is the derivative of
\(\langle u,H_{r_\xi}u\rangle_{L^2(\nu)}\) with \(u\) held fixed.
For that fixed \(u\) the scalar law identifies this quadratic form
with \(\int_M\Scal_{u^2r_\xi}\,dV_{u^2r_\xi}\). In a variation
\(q=\dot k=u^2h\) the derivative of scalar curvature is
\[
 \dot{\Scal}_k=-\langle\Ric_k,q\rangle_k
          +\nabla^i\nabla^j q_{ij}
          -\nabla^i\nabla_i\tr_k q .
\]
The divergence terms integrate to zero. The volume variation has
zero coefficient since \(\Scal_k=0\). Consequently
\begin{equation}\label{mscost:eigenvariation}
 d\lambda_1(r)[h]=-\int_M\langle\Ric_k,u^2h\rangle_k\,dV_k .
\end{equation}

The Ricci tensor of \(k\) is nonzero somewhere in \(M\setminus K\).
If it vanished throughout that open exterior, the Schouten formula
would give \(B_k=0\) there. On \(K^\circ\), the metric \(k\) is
conformally flat, so \(B_k\) vanishes there as well. Smoothness
across the boundary and conformal covariance would then give
\(B_r=0\) on \(M\). Proposition~\ref{var:stationarity} would make
\(r\) conformally round, in contradiction with \(Y([r])=0\).
Choose a nonnegative cutoff \(\chi\) supported near a point of
nonzero exterior Ricci curvature and positive at that point.
The tensor \(w=\chi u^{-2}\Ric_k\) is \(r\)-trace-free and
\(d\lambda_1(r)[w]=-\int\chi|\Ric_k|_k^2\,dV_k<0\).

Use the exterior probes of Section~\ref{msreg:section}. Denote the
compared linearization at \(r(e^{j^\sharp}\cdot,\cdot)\) by
\(\mathscr A(j)\), where the operator at the new reference is
linearized in its own symmetric chart and transported to the
\(r\)-bundles. If the old cokernel has dimension \(m>0\), these probes
give variations \(v_N\) and exact old kernel columns \(k_{j,N}\)
for which the matrix
\[
 \bigl(\langle Z_i,D\mathscr A(0)[v_N] k_{j,N}\rangle_{L^2(r)}
       \bigr)_{i,j=1}^m
\]
has an invertible limit. Here \(Z_i\) is an old adjoint basis.
The perturbations satisfy \(\|v_N\|_\infty=O(N^{-2})\). The columns
have bounded \(L^2\) norms because the Lie variations are bounded
in \(L^2\) and their projection corrections tend to zero there.
By \eqref{mscost:eigenvariation},
\[
 \beta_N=-\frac{d\lambda_1(r)[v_N]}{d\lambda_1(r)[w]}=O(N^{-2}).
\]
When the variation is changed to \(v_N+\beta_Nw\), each matrix entry
changes by at most
\[
 |\beta_N|\,
     \|(D\mathscr A(0)[w])^*Z_i\|_{L^2(r)}
     \|k_{j,N}\|_{L^2(r)}=O(N^{-2}).
\]
The adjoint in this estimate has the fixed smooth coefficient \(w\).
Thus one can fix a finite \(N\) with invertible corrected matrix.

On \(r_{a,b}=r(e^{(av_N+bw)^\sharp}\cdot,\cdot)\) the implicit
function theorem gives \(b=b(a)\) with \(b(0)=0\),
\(b'(0)=\beta_N\) and \(\lambda_1(r_{a,b(a)})=0\).
The compared mixed derivative of the Weyl linearization along the
curve is precisely the corrected one. The column argument of
Section~\ref{msreg:section} therefore gives surjectivity for all
sufficiently small \(a\ne0\). The metrics coincide with \(r\) near
\(K\) because the supports are compact in the exterior, and the
trace-free exponents give exactly the density \(\nu\). The sign
comparison gives \(Y([r_{a,b(a)}])=0\). If \(m=0\), \(r\) itself is
regular. The parameter neighborhoods in the argument are local
to the supplied reference.
\end{proof}

\subsection{Diverging strain along marked doubles}

\begin{theorem}\label{mscost:doubles}
On standard smooth \(S^4\) with the fixed round marking of
Section~\ref{sec:wall}, there are smooth regular references
\(g_R^{(0)}\) and \(g_R^{(+)}\) with
\[
 \begin{aligned}
 Y([g_R^{(0)}])&=0, &
 0<Y([g_R^{(+)}])&=O(R^{-1/2}),\\
 E_+(g_R^{(\bullet)})&=8\pi^2+O(R^{-1}),&
 \mathfrak B(g_R^{(\bullet)})&=O(R^{-1}) .
 \end{aligned}
\]
For either sequence, every \(h_R\) with
\(L_{g_R^{(\bullet)}}h_R=-W^+_{g_R^{(\bullet)}}\) satisfies
\(\|h_R\|_{L^\infty(g_R^{(\bullet)})}\ge cR\) for large \(R\),
for a constant \(c>0\) independent of \(R\).
\end{theorem}

\begin{proof}
Start with the marked zero-sign metrics of Theorem~\ref{wall:main}
and apply Proposition~\ref{mscost:zero} at each \(R\).
Continuity of \(E_+\) and \(\mathfrak B\) in \(C^4\) near each
smooth positive metric allows the change of both integrals to be at
most \(1/R\). The resulting regular references have Yamabe constant
zero and retain the marking.

For the positive sequence take the marked positive endpoints of
the same family. The fixed axial test on the neck plateau obeys the
quotient bound \(CR^{-1/2}\) before marking. The local replacement is
disjoint from its support, and the conformal factor and pullback
transfer the test and its quotient to the marked representative.
Thus those endpoints already satisfy the displayed estimates.
Positivity is open by the Rayleigh form comparison used in
Section~\ref{sec:wall}. Apply the unconstrained exterior perturbation
of Section~\ref{msreg:section}. For each \(R\), choose a reference
sufficiently close in \(C^4\) to retain positivity and to change the
two curvature integrals and the transferred test quotient by at most
\(1/R\) each. No uniform eigenvalue gap is used.

For either sequence, \(E_+\ge4\pi^2\) for large \(R\) and
\(\mathfrak B\le C_B/R\) for a fixed \(C_B>0\). The radial source
has a preimage by regularity, so \eqref{mscost:strain} also gives
\(\mathfrak B>0\) at these references. The same inequality yields
\(\|h_R\|_\infty\ge 16\pi^2 R/C_B\) for every radial preimage.
\end{proof}

The estimate measures the prescribed direction through its pairing
with \(W_g^+\) itself. An \(L^2\) co-norm bound tests all target
directions and measures preimages in \(L^2\). In the long flat collars of
Section~\ref{sec:collars}, the test modes for small singular values
have zero pairing with that source. Here the strain lower bound
holds for the radial forcing at smooth references with an onto
linearization. It concerns the initial derivative, with the
Yamabe and marking conditions imposed on the reference metrics.
Quantitative estimates for a finite step, motions without a
differentiable initial tangent, and classes with a fixed positive
energy margin below \(8\pi^2\) remain separate questions.

\section{Smooth rescalings below the scalar-flat threshold}
\label{msale:section}

A small total Bach norm passes to a smooth pointed rescaling.
Below \(8\pi^2\), the inherited Yamabe and curvature bounds restrict
the topology of a rescaling with asymptotically locally Euclidean
(ALE) ends. We first derive these bounds on a supplied smooth limit
and then compactify the only end allowed by them.

\subsection{Scalar curvature, the Bach equation, and curvature budgets}

Let \(g_i\) be smooth unit-volume Yamabe minimizers on closed
oriented smooth homotopy four-spheres \(M_i\). The conformal classes
may have round markings. The argument applies to the underlying
classes and allows the smooth type of \(M_i\) to vary. Put
\[
 Y_i=Y([g_i])=s_{g_i}>0,\qquad
 E_i=E_+(g_i)=E_-(g_i),\qquad
 D_i=\int_{M_i}|\Ric_{0,g_i}|^2\,dV_{g_i}.
\]
Assume
\begin{equation}\label{msale:source}
 4\pi^2\le E_i\le 8\pi^2-\delta,\qquad
 \frac{\int_{M_i}|B_{g_i}|\,dV_{g_i}}{E_i}\longrightarrow0
 \quad\text{for a fixed }\delta>0.
\end{equation}
The marking and the Bach assumption are conformally compatible
because the total Bach norm in \eqref{var:bach} is invariant.
By \eqref{eq:GB} and the sharp round Yamabe bound
\(Y_{\rd}=8\sqrt6\,\pi\),
\begin{equation}\label{msale:source-euler}
 Y_i^2=48(8\pi^2-E_i)+12D_i,\qquad
 \sqrt{48\delta}\le Y_i\le Y_{\rd},\qquad
 D_i\le 4E_i .
\end{equation}
The same estimates, apart from the lower bound on \(E_i\), use
only the upper bound in \eqref{msale:source}. In the four-index
tensor norm the curvature decomposition gives
\[
 \int_{M_i}|\Rm_{g_i}|_{\mathrm{tens}}^2\,dV_{g_i}
 =8E_i+2D_i+\frac{Y_i^2}{6}.
\]

Suppose \(\lambda_i\downarrow0\) and
\((M_i,\lambda_i^{-2}g_i,x_i)\) converge smoothly to a pointed
connected oriented smooth four-manifold \((X,h,x)\). Here smooth
convergence supplies orientation-preserving embeddings of
neighborhoods of each fixed compact subset of \(X\) into \(M_i\),
with the base points identified and with smooth convergence of
the pullbacks. Take a subsequence with \(E_i\to e\) and \(Y_i\to y\).
By \eqref{msale:source-euler},
\[
 D_i\longrightarrow D=\frac{y^2}{12}-4(8\pi^2-e).
\]
The rescaled scalar curvature is \(\lambda_i^2Y_i\), so \(s_h=0\).
The rescaled total Bach norms tend to zero by \eqref{msale:source}.
Smooth convergence passes the Bach tensor on every embedded
compact set, where its \(L^1\) integral in the limit must vanish.
Thus \(B_h=0\).

Let \(Y_c(X,h)\) be the infimum of the Yamabe quotient on nonzero
compactly supported smooth functions. A test on \(X\) can be
transported to \(M_i\) by a comparison embedding and extended by
zero. The rescaled Yamabe inequality for this fixed test has
scalar term \(\lambda_i^2Y_i\int f^2\). Passing to the limit gives
\begin{equation}\label{msale:sobolev}
 y\|f\|_{L^4(h)}^2\le 6\|df\|_{L^2(h)}^2,\qquad
 Y_c(X,h)\ge y\ge\sqrt{48\delta}.
\end{equation}
The Yamabe inequality applies to real tests by approximation of
their absolute values by smooth positive functions. Taking the
constant test in \eqref{msale:sobolev} would contradict compactness
of a closed \(X\)~\cite[\S\S1--3]{LeeParker1987Yamabe}.

The two chiral Weyl integrals and the trace-free Ricci squared
integral are invariant under constant rescaling in dimension
four. The chiral labels are fixed by the orientation of the
embeddings. Smooth convergence on a fixed compact domain and
then monotone exhaustion give
\begin{equation}\label{msale:budgets}
 e_\pm:=E_\pm(h)\le e,\qquad
 d:=\int_X|\Ric_h|^2\,dV_h\le D,\qquad
 \int_X|\Rm_h|_{\mathrm{tens}}^2\,dV_h
       =4(e_++e_-)+2d\le8e+2D .
\end{equation}
The bound for the full Ricci tensor uses \(s_h=0\).
With this order of limits, the source integrals outside the
embedded domains tend to \(e-e_\pm\) for each block, \(D-d\)
for the trace-free Ricci tensor, and \(y^2\) for the scalar
curvature squared. In particular all scalar-square mass lies
outside the limiting exhaustion although the rescaled scalar
curvatures converge uniformly to zero. These are residuals
relative to the supplied convergence. Their location in the
source sequence is undetermined.

\subsection{An embedded ALE core and its boundary terms}

Suppose \((X,h)\) is complete and has finitely many ends, each
with an oriented coordinate identification with
\((\mathbb R^4\setminus\overline{B})/\Gamma_j\).
Here \(\Gamma_j\subset SO(4)\) is finite and acts freely on
\(S^3\). We require differentiable ALE decay of positive order
\(\alpha_j\), so on each lifted end
\begin{equation}\label{msale:end-decay}
 \partial^\beta(h_{ab}-\delta_{ab})
   =O(r^{-\alpha_j-|\beta|}),\qquad
 \alpha_j>0,\quad 0\le|\beta|\le5.
\end{equation}
Two orders suffice for the following boundary arguments. Write
\(k\) for the number of ends, \(q_j=|\Gamma_j|\), and
\(L_j=S^3/\Gamma_j\) for the outward oriented links. A large
compact truncation \(K\) is connected, is a deformation retract
of \(X\), and has boundary \(L=\bigsqcup_j L_j\). The convergence
embeds the entire \(K\) into one source homotopy sphere at a
time. This requires comparison maps on the whole truncation
including all of its boundary components.

For one such embedding write \(V\) for the closure of the
complement. The collars of \(L\) give a Mayer--Vietoris
decomposition. Each \(L_j\) is a rational homology sphere with
\(H_2(L_j;\mathbb Z)=0\), by duality and the finiteness of its
first homology. Since the ambient space is an integral homology
four-sphere, the sequence gives
\begin{equation}\label{msale:mv}
 H_2(K;\mathbb Z)=H_2(V;\mathbb Z)=0,\qquad
 H_1(L;\mathbb Z)\cong H_1(K;\mathbb Z)\oplus H_1(V;\mathbb Z).
\end{equation}
The injection \(H_0(L)\to H_0(K)\oplus H_0(V)\) prevents two
boundary components from meeting the same component of \(V\).
Every component of \(V\) meets the boundary, by connectedness
of the ambient manifold. Thus \(V\) has \(k\) components.
The integral relative fundamental class of \(K\) maps to the
sum of its boundary classes. Since \(H_1(K;\mathbb Z)\) is
finite, \(H_3(K,L;\mathbb Z)\cong H^1(K;\mathbb Z)=0\), so
\(H_3(K;\mathbb Z)\cong\mathbb Z^{k-1}\). The intersection form
on the image of relative
middle-degree cohomology has zero signature by
\eqref{msale:mv}. In particular
\begin{equation}\label{msale:core-topology}
 \chi(X)=\chi(K)=2-k,\qquad \tau(K)=0.
\end{equation}

Chern--Gauss--Bonnet and the signature formula with boundary give
\begin{align}
 8\pi^2\left(\chi(X)-\sum_{j=1}^k q_j^{-1}\right)
   &= e_++e_--\tfrac12d, \label{msale:ae-euler}\\
 \tau(K)
   &=\frac{e_+-e_-}{12\pi^2}
           -\sum_{j=1}^k\eta_{\mathrm{sig}}(L_j).
         \label{msale:ae-signature}
\end{align}
For the first identity, the curvature is \(O(r^{-2-\alpha_j})\).
The Euler boundary integral is invariant under a constant change
of scale. In a rescaled annulus around the cut it tends to that
of the outward Euclidean unit sphere divided by \(q_j\), by
\(C^2\) convergence and the expressions for the curvature and
second fundamental form. Since the Euclidean unit ball has zero
bulk Euler density and Euler characteristic one, the boundary
term is \(+8\pi^2/q_j\) in \eqref{eq:GB}'s normalization.
The scalar term vanishes for \(h\)~\cite[\S7, (7.2)]{TianViaclovsky2005ALE}.

Here \(\eta_{\mathrm{sig}}\) is the unhalved eta value of the
odd signature operator on even forms for the round quotient
metric, with zero eigenvalues omitted. In this convention the
product-boundary signature formula is
\(\tau(K)=\int_K L_1-\eta_{\mathrm{sig}}(\partial K)\), with
\(L_1=(|W^+|_{\HS}^2-|W^-|_{\HS}^2)dV/(12\pi^2)\)
as in \eqref{eq:sign}. To apply it, change the metric
conformally by \(r^{-2}\) on each end. In \(t=\log r\) the
new metric differs from the round product by \(O(e^{-\alpha_j t})\)
through two derivatives on collars of fixed \(t\)-width.
Interpolate to the exact round product in the last such collar
of each compact truncation. The product is locally conformally
flat, so the new collar has \(O(e^{-2\alpha_j t})\) Weyl energy.
The bulk integrals converge to \(e_\pm\) by conformal invariance,
finite energy, and the vanishing collar error. The signatures
of the truncations agree. The product formula then gives
\eqref{msale:ae-signature} with precisely the round link metrics
and the stated orientations~\cite[Theorem (4.14)]{APS1975SpectralI}.

In particular the source and limit Euler identities read
\begin{equation}\label{msale:euler-residual}
 8\pi^2\left(2-\chi(X)+\sum_j q_j^{-1}\right)
   =(e-e_+)+(e-e_-)-\tfrac12(D-d)+\tfrac1{24}y^2 .
\end{equation}
The signature identity measures the difference of the chiral
residuals through the sum of the link eta values.

\subsection{One end and a trivial spherical link}

\begin{theorem}[Topology of a smooth ALE limit]\label{msale:ends}
Under \eqref{msale:source}, suppose the complete smooth pointed
limit has finitely many differentiably ALE ends of positive order
and its full compact truncations embed by the convergence.
Then it has precisely one end, the link is \(S^3\), and its two
chiral Weyl energies agree.
\end{theorem}


For an end of group order \(q_j\), the compactly supported Yamabe
constant satisfies
\begin{equation}\label{msale:end-yamabe}
 Y_c(X,h)\le \frac{Y_{\rd}}{\sqrt{q_j}} .
\end{equation}
Indeed, take the radial Euclidean sharp test
\(U(x)=(1+|x|^2)^{-1}\), whose quotient is \(Y_{\rd}\).
Cut it off on \(\epsilon<r<2L\), with transition annuli
\((\epsilon,2\epsilon)\) and \((L,2L)\). As
\(\epsilon\downarrow0\) and \(L\uparrow\infty\), the tests converge
to \(U\) in \(\dot H^1\cap L^4\). The squared Dirichlet errors
from the cutoff gradients are \(O(\epsilon^2)\) and \(O(L^{-2})\).
For a fixed cutoff profile, dilate its support out the given end.
The rescaled end metrics converge to the Euclidean metric on
the fixed annulus. The profile descends to the quotient, where
each integral is divided by \(q_j\), and the denominator of the
quotient is a square root. Let the dilation tend to infinity
first and then remove the cutoffs
~\cite{Talenti1976BestConstant}\cite[Theorem 3.3]{LeeParker1987Yamabe}.

Set \(q_{\max}=\max_j q_j\). Combining the Ricci budget in
\eqref{msale:budgets} with \eqref{msale:source-euler},
\eqref{msale:core-topology}, and \eqref{msale:ae-euler} gives
\begin{equation}\label{msale:end-count}
 (8\pi^2-e)+\frac{e_++e_-}{2}
 +4\pi^2\left(k-2+\sum_jq_j^{-1}\right)
 \le\frac{y^2}{48}\le\frac{8\pi^2}{q_{\max}} .
\end{equation}
The upper bound follows from \eqref{msale:sobolev} and
\eqref{msale:end-yamabe}. If \(k\ge2\), the left side minus the
rightmost term is at least
\(\delta+4\pi^2(k-2)(1+q_{\max}^{-1})\), using
\(q_j^{-1}\ge q_{\max}^{-1}\) and \(e_++e_-\ge0\).
Thus \(k=1\).

The link of the unique end smoothly embeds in an integral homology
four-sphere, by restriction of the full-core embedding. Hillman's
locally flat embedding theorem with this ambient hypothesis gives,
in the finite fundamental group case, \(1\), \(Q_8\), or the
binary icosahedral group for the Poincar\'e homology sphere
~\cite[Theorem 5.2(1) and Corollary 4.4.2]{Hillman2024LocallyFlat}.
For that last case the embedded core is an integral homology ball,
by \eqref{msale:mv} and the relative homology calculation for
\(k=1\). Its spin structure restricts to the unique spin structure
on the link. By uniqueness of smoothing in dimension three, the
link is diffeomorphic up to orientation to the boundary of the
negative \(E_8\) plumbing of disk bundles of Euler number
\(-2\). This plumbing is simply connected with even intersection
form, so it is spin and has signature \(-8\). Gluing to the core
with matching opposite boundary orientations would give a smooth
closed spin four-manifold of signature \(\pm8\), in contradiction
with Rokhlin divisibility by \(16\)
~\cite[Descriptions (1) and (4)]{KirbyScharlemann1979Poincare}%
\cite[\S1.2]{Rokhlin1971Submanifolds}.

For a free orthogonal action of \(Q_8\), the central involution
acts as \(-\mathrm{Id}\). The operators \(I,J\) for the quaternionic
generators are skew-orthogonal, satisfy \(I^2=J^2=-\mathrm{Id}\),
and anticommute. For any unit vector \(v\) the vectors
\((v,Iv,Jv,IJv)\) form an orthonormal basis identifying the action
with left quaternion multiplication. Thus the round metric on
the quotient is determined up to orientation. In the quaternion
orientation, a generator of each of the three order-four cyclic
subgroups has oriented rotation angles \((\pi/2,\pi/2)\). For an
order-\(m\) cyclic subgroup the signature eta value on the round
quotient is
\[
 -\frac1m\sum_{\ell=1}^{m-1}
     \cot(\ell\theta_1/2)\cot(\ell\theta_2/2).
\]
It is \(-1/2\) in each order-four case and \(0\) for the common
order-two antipodal group. The \(Q_8\) eta function is the average
of the equivariant eta functions, by projection onto invariants
on each eigenspace. The three cyclic subgroups cover \(Q_8\) and
their pairwise intersections are the common order-two subgroup.
Consequently the value at zero is
\(\tfrac18[3\cdot4(-1/2)-2\cdot2(0)]=-3/4\).
Orientation reversal negates it
~\cite[\S2, Proposition (2.12)]{APS1975SpectralII}.
Equations \eqref{msale:core-topology} and
\eqref{msale:ae-signature} would then give
\(|e_+-e_-|=9\pi^2\), impossible since both energies belong to
\([0,8\pi^2-\delta]\).

The link is therefore \(S^3\) and the end is asymptotically
Euclidean (AE). The even-form signature eta value on the round
sphere is zero by orientation-reversing isometry. By
\eqref{msale:core-topology}, \eqref{msale:ae-euler}, and
\eqref{msale:ae-signature},
\begin{equation}\label{msale:ae-budgets}
 e_+=e_-=e_X,\qquad \chi(X)=1,\qquad d=4e_X .
\end{equation}
Thus a Ricci-flat or half-conformally-flat limit here has zero
full curvature. Van Kampen across the \(S^3\) link makes the
fundamental group of the core a free factor of \(\pi_1(M_i)=0\).
The complete flat limit is consequently Euclidean.

\paragraph{Equality in the two-end test.}
If the upper margin is removed and \(e=8\pi^2\), two ends in
\eqref{msale:end-count} require equal group orders, zero Weyl energy,
and equality in the Ricci and Yamabe bounds. The scalar-flat
Schwarzschild metric \(h=(1+b/r^2)^2g_{\rm eucl}\), \(b>0\), on
\(\mathbb R^4\setminus\{0\}\) realizes the limit-side constants
\cite[remark after Theorem 1.3]{TianViaclovsky2005ALE}.
Its harmonic conformal factor makes \(W=0\) and \(\Scal=0\).
The variable \(z=r-b/r\) gives
\(h=dz^2+(z^2+4b)g_{S^3}\) and displays the two complete Euclidean
ends. The Euler formula with \(\chi=0\) and \(q_1=q_2=1\) gives
\(\int|\Ric|^2=32\pi^2\). Conformal covariance on compact tests
and the zero capacity of the puncture give \(Y_c=Y_{\rd}\).
Existence of a closed Yamabe sequence realizing this model at
equality is a separate question.

\subsection{Smooth conformal compactification}

\begin{proposition}\label{msale:compact}
Let \((X,h)\) be smooth, complete, connected, oriented, scalar-flat
and Bach-flat, with one AE end satisfying \eqref{msale:end-decay}.
There is a smooth one-point conformal compactification
\((\widehat X,\bar h)\), with \(\bar h=a^2h\) on \(X\) for a smooth
positive function \(a\). The extended metric is scalar-flat near
the added point and Bach-flat everywhere. Its Yamabe constant is
\(Y_c(X,h)>0\), and its chiral Weyl integrals equal those of \(h\).
If the full core embeds smoothly in a homotopy four-sphere, then
\(\widehat X\) is a smooth homotopy four-sphere.
\end{proposition}

\begin{proof}
Write \(\mathcal D_h=\operatorname{div}_h\nabla\) and fix
\(0<\epsilon<\min\{1,\alpha_1\}\). In the end coordinates
\(\mathcal D_h r^{-2}=O(r^{-4-\alpha_1})\) while
\[
 \mathcal D_h r^{-2-\epsilon}
   =\epsilon(2+\epsilon)r^{-4-\epsilon}
                  +O(r^{-4-\epsilon-\alpha_1}).
\]
On a sufficiently distant end,
\(\tfrac12(r^{-2}+r^{-2-\epsilon})\) and
\(2(r^{-2}-r^{-2-\epsilon})\) are positive sub- and supersolutions
below and above \(r^{-2}\). Solve
\(\mathcal D_h u=0\) on concentric bounded annuli with inner
radius fixed and outer radius tending to infinity, prescribing
\(r^{-2}\) on the two boundary spheres. The maximum principle
gives the same two barriers on every annulus. Interior estimates
give a positive harmonic limit bounded above and below by fixed
positive multiples of \(r^{-2}\). On dilated annuli, the metric
bounds in \eqref{msale:end-decay} are uniform and the rescaled
inverse metric and first-order coefficients in the harmonic
equation are bounded in \(C^{3,\beta}\) for \(0<\beta<1\).
Interior estimates give \(\partial^j u=O(r^{-2-j})\) through
order five.

To obtain the leading coefficient, cut \(u\) off smoothly across
the inner end to a function on \(\mathbb R^4\), still denoted by
\(u\). Then \(F=\mathcal D_{\mathrm{Eucl}}u=O(r^{-4-\alpha_1})\)
has finite absolute \(\epsilon\)-moment. The Euclidean Newtonian
potential \(\Phi*F\), with \(\Phi\) the fundamental solution of
\(\mathcal D_{\mathrm{Eucl}}\), has the expansion
\[
 (\Phi*F)(x)=\Phi(x)\int_{\mathbb R^4}F+O(r^{-2-\epsilon}).
\]
On \(|z|\le r/2\) the kernel difference is bounded by
\(C r^{-2-\epsilon}|z|^\epsilon\). On the remaining region,
within \(|x-z|<r/4\) use \(F=O(r^{-4-\alpha_1})\) and the
integrability of the kernel there, and outside that ball use
the bound \(C r^{-2}\) times the tail mass. The potential
therefore tends to zero. Its difference from the cut-off \(u\)
is entire harmonic and tends to zero, hence vanishes.
The lower barrier gives a positive coefficient \(c\) of \(r^{-2}\).
Apply the same rescaled interior estimates to
\(u-cr^{-2}\), which solves
\(\mathcal D_h(u-cr^{-2})=-c\mathcal D_hr^{-2}\) on the end.
Its zeroth-order error is \(O(r^{-2-\epsilon})\) and the rescaled
forcing is smaller by the choice of \(\epsilon\).
Divide by \(c\) and extend smoothly and positively across the
core. The resulting \(a\) satisfies
\begin{equation}\label{msale:harmonic}
 \mathcal D_h a=0\ \text{near infinity},\qquad
 \partial^j(a-r^{-2})=O(r^{-2-\epsilon-j}),\quad 0\le j\le5.
\end{equation}

With \(\Delta_h=-\mathcal D_h\), conformal covariance gives
\(s_{a^2h}=a^{-3}(6\Delta_h+s_h)a\). Thus \(a^2h\) is
scalar-flat near infinity and Bach-flat on \(X\).
In the inverted coordinate \(z=x/|x|^2\), with \(\rho=|z|\),
the derivative of inversion is \(\rho^{-2}\) times an orthogonal
reflection. The leading conformal and inversion factors cancel,
and differentiating \eqref{msale:end-decay} and
\eqref{msale:harmonic} gives
\[
 \partial_z^j(\bar h_{ab}-\delta_{ab})
                  =O(\rho^{\epsilon-j}),\qquad 0\le j\le5.
\]
The added point therefore satisfies the \(C^0\) orbifold
coordinate condition with trivial group of Ache and Viaclovsky,
whose removability theorem requires constant scalar curvature and
obstruction-flatness. In four dimensions obstruction-flatness is
Bach-flatness. Their finite-derivative form uses \(2(k+1)+1=5\)
derivatives here, for \(k=1\). The metric extends smoothly in the
resulting smooth chart, and the curvature equations extend by
continuity~\cite[Definition 1.6, Theorem 1.8 and Remark 7.2]
{AcheViaclovsky2011Obstruction}. The harmonic gauge provides the
constant scalar curvature needed for this application.

For the topology, use \eqref{msale:mv} for the embedded truncation
with \(S^3\) boundary. Van Kampen shows the core and its complement
are simply connected, and Mayer--Vietoris makes each an integral
homology ball. The one-point addition fills the end by the cone
on \(S^3\) topologically, with its smooth chart supplied by the
removability theorem. The compactification is simply connected
and has the integral homology of \(S^4\). Hurewicz and Whitehead
then give a homotopy equivalence with \(S^4\). This supplies the
admissibility of the capped smooth type independently of the
smooth type of the original ambient complement.

Conformal covariance of the Yamabe quotient identifies compact
tests on \(X\) for \(h\) with compact tests away from the added
point for \(\bar h\), by multiplication by \(a\) in the direction
from \(\bar h\) to \(h\). A cutoff changing between radii
\(\rho\) and \(2\rho\) near the added point has squared
Dirichlet norm \(O(\rho^2)\). Multiplying any fixed smooth
global test by this cutoff approximates it in \(H^1\) and
\(L^4\). Thus
\begin{equation}\label{msale:cap-yamabe}
 Y([\bar h])=Y_c(X,h).
\end{equation}
The Weyl densities are conformally invariant and the point has
measure zero.

For positivity even without a source sequence, put \(G=a^{-1}\).
It solves \((6\Delta_{\bar h}+s_{\bar h})G=0\) off the
added point \(p\). It is locally integrable since
\(G\,dV_{\bar h}=a^3 dV_h\). On the expanding coordinate spheres
with outward normal, the conformal changes of normal and measure
give \(\partial_{\nu_{\bar h}}G\,dS_{\bar h}
=-\partial_{\nu_h}a\,dS_h\), whose integral tends to \(4\pi^2\).
In Green's identity the other boundary term tends to zero, since
\(G=O(r^2)\), the boundary area of \(\bar h\) is \(O(r^{-3})\),
and a smooth test has bounded \(\bar h\)-gradient. Therefore
\((6\Delta_{\bar h}+s_{\bar h})G=24\pi^2\delta_p\) in
distributions. Pairing with a positive first eigenfunction gives
a strictly positive first eigenvalue. The compact Sobolev
inequality and this spectral gap give positivity of the
Yamabe constant~\cite[\S\S2, 6]{LeeParker1987Yamabe}.
\end{proof}

\subsection{Strict residual energy and the Green construction}

For the supplied limit in \eqref{msale:source}, set
\(\widehat Y=Y([\bar h])\). Equations \eqref{msale:sobolev} and
\eqref{msale:cap-yamabe} give \(\widehat Y\ge y\).
The Ricci budget \(4e_X=d\le D\) of \eqref{msale:ae-budgets}
and \eqref{msale:source-euler} yield
\begin{equation}\label{msale:residual}
 e-e_X\ge \frac{Y_{\rd}^2-y^2}{48}.
\end{equation}
When \(h\) is nonflat, \(e_X>0\) by \eqref{msale:ae-budgets}.
Then the compactification has nonzero Weyl energy, and Schoen's
strict Yamabe inequality gives
\begin{equation}\label{msale:strict-residual}
 e-e_X\ \ge\ \frac{Y_{\rd}^2-\widehat Y^2}{48}>0 .
\end{equation}
In fact equality \(e=e_X\) in the energy budget would force
\(y=Y_{\rd}=\widehat Y\), conformal roundness of the
compactification, and hence a Euclidean limit by
\eqref{msale:ae-budgets}. In the positive source band even a
flat limit has strict residual Weyl energy
~\cite[Theorem 2]{Schoen1984Yamabe}.

Conversely, on a supplied smooth closed connected oriented
positive-Yamabe Bach-flat four-manifold \((N,\gamma)\), choose \(p\in N\) and
a conformal normal representative at \(p\). Let \(G_p\) be the
positive conformal-Laplacian Green function with leading term
\(\rho^{-2}\) in these coordinates, where \(\rho\) is the distance
to \(p\). Then \(G_p^2\gamma\) is scalar-flat, Bach-flat,
and complete on \(N\setminus\{p\}\) with one AE end of order two.
Indeed \(\gamma_{ij}=\delta_{ij}+O(\rho^2)\) and
\(G_p=\rho^{-2}+A+O(\rho)\) in dimension four. Applying rescaled
elliptic estimates on annuli to the bounded error
\(G_p-\rho^{-2}\) gives the differentiated bounds needed in
\eqref{msale:end-decay} after inversion. The compact test Yamabe
constant and chiral integrals agree with their closed
counterparts. If \(N\) is a homotopy four-sphere, all full compact
truncations embed in \(N\). A nonround closed class therefore
supplies a nonflat AE metric with these properties
~\cite[Lemma 6.4(a) and Theorem 6.5]{LeeParker1987Yamabe}.
The converse is a solution of a puncture construction on a given
closed class. Smooth extraction of rescalings and their end
hypotheses from a small total Bach norm remain separate questions.

\section{Closed Bach-flat critical values and energy loss}
\label{msgap:section}

A conformal compactification as in Section~\ref{msale:section}
relates a noncompact concentration limit to a closed Bach-flat
metric. The closed problem imposes a lower energy threshold on
that component. The Yamabe constant of the compactification also
controls the energy left outside it.

\subsection{The closed threshold and the energy of one component}

\begin{theorem}[Closed Bach-flat gap]\label{msgap:closed}
Let \(g\) be a smooth Riemannian Bach-flat metric of positive class
Yamabe constant on a smooth closed oriented homotopy four-sphere.
Then either the class is conformally diffeomorphic to the round
sphere or
\[
 E_+(g)=E_-(g)>4\pi^2.
\]
\end{theorem}

\begin{proof}
By \eqref{eq:sign}, the four-index Weyl integral is \(8E_+(g)\).
Below \(4\pi^2\) in a single block, Proposition~\ref{prop:recognition}
identifies the smooth type with \(S^4\). At \(E_+=4\pi^2\) the
four-index integral equals \(16\pi^2\chi(M)\).
If the smooth type were neither \(S^4\) nor \(\mathbb{RP}^4\),
the equality theorem of Chang, Gursky and Yang would identify the
manifold after a conformal change with a complex projective plane
or with a manifold isometrically covered by
\(S^3\times S^1\)~\cite[Theorem B]{ChangGurskyYang2003}.
The former has nonzero second Betti number, and the latter
is impossible for a simply connected base. Simple connectivity
also excludes \(\mathbb{RP}^4\). Thus in either energy range
\(E_+\le4\pi^2\) the underlying smooth manifold is standard \(S^4\).
The closed Bach-flat rigidity theorem of Chang and Chen applies
there at the inclusive threshold \(32\pi^2\) in the full tensor
norm and gives \(W=0\)~\cite[Remark 2 and Theorem 2.1]{ChangChen2021RicciSphere}.
Proposition~\ref{prop:recognition} identifies the resulting
conformal class. In particular the nonzero endpoint
\(E_+=4\pi^2\) is excluded.
\end{proof}

For the limit statements, use smooth volume-one Yamabe minimizers
\(g_i\), with \(Y_i=Y([g_i])>0\), and
\[
 E_i=E_+(g_i)=E_-(g_i)\in[4\pi^2,8\pi^2-\delta],
 \qquad E_i\longrightarrow e,\qquad Y_i\longrightarrow y
\]
on a common subsequence, where \(\delta>0\) is fixed. Write
\(Y_{\rd}=8\sqrt6\,\pi\). For a supplied nonflat smooth ALE
rescaling satisfying the hypotheses of Section~\ref{msale:section},
write \(e_X=E_\pm(X,g_X)\) for the common chiral energy and
\(\widehat Y_X=Y([\widehat g_X])\) for its compactification's
Yamabe constant. The one-point compactification is smooth,
Bach-flat and of positive Yamabe constant, with the topology
of a homotopy four-sphere. Its two block energies equal those
of \(X\). In the scalar-flat open Gauss--Bonnet formula the
Euclidean end contributes one, and \(\chi(X)=1\), so
\[
 \int_X|\Ric_{g_X}|^2\,dV_{g_X}=4e_X .
\]
If \(e_X=0\), this identity gives flatness. Consequently every
nonflat such limit has \(e_X>4\pi^2\) by
Theorem~\ref{msgap:closed}, and its compactification has
\(\widehat Y_X<Y_{\rd}\) by sharp Yamabe
rigidity~\cite[Theorem 2]{Schoen1984Yamabe}. Using the
inequality \(y\le\widehat Y_X\) and \eqref{msale:residual} gives
\begin{equation}\label{msgap:strict}
 e-e_X\ \ge\ \frac{Y_{\rd}^2-y^2}{48}
 \ \ge\ \frac{Y_{\rd}^2-\widehat Y_X^2}{48}
 \ >\ 0 .
\end{equation}
This comparison uses the full Ricci integral of the scalar-flat
limit and the trace-free Ricci integral of the closed source.
In fact the latter controls the former on each embedded compact
truncation by smooth convergence and constant-scale invariance.
It follows that \(e-e_X<4\pi^2-\delta\).
Moreover two nonflat components cannot have simultaneously
disjointly accounted cores, since each contributes strictly more than
\(4\pi^2\) to either chiral energy. This use of disjointness concerns
oriented embeddings of fixed truncations into the same source.

\subsection{Compactness for exactly Bach-flat sources}

\begin{lemma}[An exact compactness alternative]\label{msgap:compactness}
Let \(g_i\) be volume-one Yamabe minimizers on smooth closed oriented
homotopy four-spheres \(M_i\). Suppose that the metrics are exactly
Bach-flat, \(Y_i>0\), and the common chiral energies are in the
fixed band \([4\pi^2,8\pi^2-\delta]\).
After passing to a subsequence, one of the following holds.
\begin{enumerate}[label=\textup{(\roman*)}]
\item The metrics converge smoothly by diffeomorphisms from a
      closed smooth homotopy four-sphere. The limit is an exact
      Bach-flat volume-one Yamabe minimizer with
      \(Y=\lim_iY_i>0\) and chiral energy \(\lim_iE_i\).
\item There is a nonflat complete smooth pointed rescaling with
      scalar curvature zero, vanishing Bach tensor and finitely
      many positive-order differentiable ALE ends. Smooth
      oriented comparison maps exist on neighborhoods of
      every fixed compact subset. The end restrictions of
      Section~\ref{msale:section} give one Euclidean end,
      and Proposition~\ref{msale:compact} gives the smooth
      conformal compactification.
\end{enumerate}
\end{lemma}

\begin{proof}
By \eqref{var:wall} and the sharp sphere bound,
\[
 Y_*:=\sqrt{48\delta}\le Y_i\le Y_{\rd},\qquad
 \|\Ric_{0,i}\|_2^2\le4E_i,\qquad
 \int_{M_i}|\Rm_{g_i}|_{\mathrm{tens}}^2\,dV_{g_i}
       \le16E_i+64\pi^2 .
\]
For the last bound the curvature decomposition gives
\(8E_i+2\|\Ric_{0,i}\|_2^2+Y_i^2/6
=64\pi^2+4\|\Ric_{0,i}\|_2^2\).
The Yamabe minimizing inequality at volume one is
\[
 Y_i\|f\|_4^2\le 6\|df\|_2^2+Y_i\|f\|_2^2 .
\]
It gives the normalized Sobolev estimate
\(\|f\|_4\le \sqrt{6/Y_*}\|df\|_2+
\Vol(M_i)^{-1/4}\|f\|_2\), with volume and norms taken
in the selected metric. It is valid also after any constant rescaling.
In particular the Sobolev constant, total curvature and first
Betti number meet the hypotheses of exact Bach-flat compactness.
These are also the inputs to the first-bubble construction in
Chang--Qing--Yang~\cite[\S\S2 and 4]{ChangQingYang2006Gap}.

We give the curvature-maximum argument to specify the comparison maps.
Let \(V_x(r)=\Vol B(x,r)\) in the volume-one gauge.
Use a test supported in \(B(x,r)\), equal to one on \(B(x,r/2)\)
and with Lipschitz bound \(2/r\). For \(r\le1\) the Yamabe
inequality gives
\[
 Y_*\,V_x(r/2)^{1/2}\le (24+Y_{\rd})r^{-2}V_x(r).
\]
Thus \(V_x(r)/r^4\ge c[V_x(r/2)/(r/2)^4]^{1/2}\) for
\(c=Y_*/(4(24+Y_{\rd}))\). Iterating and using the Euclidean
small-ball asymptotic at the center gives \(V_x(r)\ge c^2r^4\).
This bound survives a constant rescaling \(A_i g_i\), \(A_i\ge1\),
for \(r\le\sqrt{A_i}\).
Packing disjoint balls of fixed small radius along a minimizing
geodesic bounds the diameter uniformly in the volume-one scale.

Let \(K_i=\max_{M_i}|\Rm_{g_i}|\). If \(K_i\) is bounded, work
at the original scale. If \(K_i\to\infty\), rescale by \(K_i\) and
base the metrics at maximum points. In each case the pointwise
curvature is bounded in the selected scale. The small-ball volume
lower bound and the sectional curvature bound give a uniform
injectivity radius~\cite[Theorem 4.7]{CheegerGromovTaylor1982}.
Bishop volume comparison gives a small-ball upper bound \(Cr^4\)
and a positive diameter lower bound, using that the total volume
in the selected scale is at least one. Fix a small radius below
half this lower diameter bound. By reducing it further, the
curvature integrals on all balls of this radius are uniformly small.
The constant-scalar Bach-flat equation and the local regularity
estimates then bound every covariant curvature derivative
~\cite[\S2.1 and Theorem 4.1]{TianViaclovsky2004Moduli}.
Smooth pointed Cheeger--Gromov compactness gives a complete
smooth limit with comparison embeddings on compact sets
~\cite[Definition 1.1 and \S2]{AndersonCheegerGromovAuthorSurvey}.
Choose the orientation of the limit consistently with the
embeddings along the subsequence.

In the bounded case, the diameter bound makes the limit compact.
An embedding of the whole compact limit has open and closed image
in each connected source and is a diffeomorphism.
Smooth convergence preserves volume, scalar curvature, the Bach
equation and the chiral integrals. The source Yamabe inequality
passes to the limit for each smooth test. Its constant test
achieves equality because the limiting scalar is \(y=\lim_i Y_i\)
and the volume is one. This proves the first alternative.

In the unbounded case, the rescaled volumes \(K_i^2\) tend to
infinity and the pointed limit has curvature norm one at its
base point. The rescaled diameter tends to infinity by Bishop
volume comparison. The limit is noncompact, scalar-flat and
Bach-flat by smooth convergence and \(Y_i/K_i\to0\).
For its ends apply the pointed theorem of Tian and Viaclovsky
to the rescaled sequence. Its assumptions are the uniform
normalized Sobolev constant, uniform full \(L^2\) curvature,
positive lower volume bound and bounded source \(b_1\).
The remark following their compact theorem extends the pointed
conclusion to constant-scalar-curvature Bach-flat metrics with
order-zero ALE asymptotics. The pointed metric limit there
agrees with the smooth limit above, so in particular its
finitely many end charts have differentiable order zero
~\cite[Theorems 1.1--1.3 and \S7]{TianViaclovsky2004Moduli}.
In dimension four Bach-flatness is obstruction-flatness,
and scalar flatness permits the decay theorem of Ache and
Viaclovsky. This improves the ALE order to two in suitable
coordinates~\cite[Definition 1.1 and Theorem 1.2]{AcheViaclovsky2011Obstruction}.
The comparison maps embed entire fixed truncations, so the
arguments of Section~\ref{msale:section} apply. Their topology
uses one comparison embedding at a time and admits varying
ambient smooth homotopy spheres. This proves the second alternative.
\end{proof}

\begin{proposition}[Uniform Yamabe deficit and residual]\label{msgap:yamabe}
For each fixed \(\delta>0\) there is \(d_\delta>0\) such that every
closed exactly Bach-flat positive-Yamabe metric on a smooth
homotopy four-sphere with energy in the band of
Lemma~\ref{msgap:compactness} satisfies \(Y\le Y_{\rd}-d_\delta\).
Consequently every supplied nonflat smooth ALE limit in this band
under the convergence and end hypotheses of Section~\ref{msale:section}
satisfies
\begin{equation}\label{msgap:uniform-residual}
 e-e_X\ \ge\ \kappa_\delta
       :=\frac{Y_{\rd}^2-(Y_{\rd}-d_\delta)^2}{48}>0 .
\end{equation}
\end{proposition}

\begin{proof}
If there were a sequence with \(Y_i\to Y_{\rd}\), the compact
alternative would give a smooth closed limit of Yamabe constant
\(Y_{\rd}\). The sharp equality theorem makes this conformal
class round~\cite[Theorem 2]{Schoen1984Yamabe}. This contradicts
smooth convergence of its Weyl energy from \(E_i\ge4\pi^2\).
In the noncompact case the compactly supported Yamabe inequality
passes to the rescaling with constant \(Y_{\rd}\), since the
rescaled scalar term tends to zero. The one-ended smooth
compactification has Yamabe constant equal to that compact-support
infimum, by conformal covariance and zero capacity of the added
point. Thus its Yamabe constant must also equal \(Y_{\rd}\).
Closed sharp rigidity gives \(W_{g_X}=0\), and the scalar-flat open
Euler identity then gives \(\Ric_{g_X}=0\), contrary to the
nonzero curvature at the base point. This excludes both cases.
The assertion is vacuous if the exact class is empty, and we
may take \(0<d_\delta<Y_{\rd}\).

For the second assertion it is the compactification of the
limit that is exactly Bach-flat. Its energy satisfies
\(4\pi^2<e_X\le e\le8\pi^2-\delta\), so its Yamabe constant
satisfies the exact-class bound. Insert that bound into
\eqref{msgap:strict}. If the exact class is empty, there can be
no supplied nonflat limit of the stated type.
\end{proof}

\subsection{Attainment of the first subcritical critical value}

Let
\begin{equation}\label{msgap:critical-set}
 \mathcal A=\{\,E_+(g): 0<E_+(g)<8\pi^2,\quad B_g=0,\quad Y([g])>0\,\}.
\end{equation}
Here \(g\) ranges over smooth Riemannian metrics on all smooth closed
oriented homotopy four-spheres.

\begin{theorem}[The first positive subcritical critical value]
\label{msgap:minimum}
If \(\mathcal A\ne\varnothing\), its infimum \(a\) is attained and
\(4\pi^2<a<8\pi^2\). Every sequence of metrics realizing values
tending to \(a\), taken in volume-one Yamabe minimizing gauges,
has a smoothly convergent subsequence up to oriented diffeomorphisms,
with an attaining metric as limit.
\end{theorem}

\begin{proof}
If \(\mathcal A\ne\varnothing\), the closed gap gives
\(a\ge4\pi^2\), and \(a<8\pi^2\) follows from any member.
Choose a minimizing sequence of values and discard a finite
prefix so that \(E_i\le8\pi^2-\delta\) for one \(\delta>0\).
Bach-flatness and the energies are conformally invariant,
so unit-volume Yamabe minimizers satisfy
Lemma~\ref{msgap:compactness}. If the rescaled case occurred,
the compactification of its nonflat limit would provide a
positive exact closed value \(e_X< a\) by \eqref{msgap:strict}
with \(e=a\). Since \(e_X>0\) and the compactification is a
positive-Yamabe smooth homotopy sphere, this value lies in
\(\mathcal A\), contradicting the definition of \(a\).
Thus the compact case gives an attaining metric. The closed
gap applied to that metric excludes \(a=4\pi^2\).
\end{proof}

With \(a=+\infty\) when \(\mathcal A\) is empty, there is a
non-effective interval of positive width above \(4\pi^2\)
free of exact closed critical values. The same interval is
excluded for an approximately Bach-flat sequence in a supplied
compact gauge that yields a smooth positive-Yamabe Bach-flat
limit on a homotopy four-sphere and preserves its full Weyl
energy. Indeed its limiting energy would be a member of
\(\mathcal A\). In the supplied smooth ALE regime a nonflat
component instead requires \(e_X\in\mathcal A\) and, by
\eqref{msgap:uniform-residual}, \(e\ge a+\kappa_\delta\) when
\(a<+\infty\). If \(\mathcal A\) is empty, neither of these two
types of nonzero limit can occur below \(8\pi^2\) under the
stated hypotheses. The minimum in Theorem~\ref{msgap:minimum}
is among critical values.

\subsection{Accounting with a smooth compactified base}

Consider now a sequence of volume-one Yamabe minimizers on a
smooth oriented homotopy four-sphere in the fixed band, allowing
\(\int |B_{g_i}|\,dV_{g_i}\to0\). Suppose the following
convergence data are supplied along the common subsequence.
There is a base on a closed smooth oriented homotopy sphere
minus finitely many points, whose limiting metric is conformal on
the punctured domain to a smooth positive metric on the closed
one. On each fixed compact base truncation suppose there is
smooth comparison convergence under embeddings into the source,
allowing positive conformal changes of the pullback metrics.
Suppose there are also complete smooth scalar-flat ALE components,
compared to the source by oriented compact embeddings and constant
length rescalings tending to zero. Each nonflat component has one
Euclidean end of positive differentiable order and a smooth
conformal one-point compactification of homotopy-sphere topology.
Convergence on a component may omit finitely many already smooth
points of that component. Assume that every fixed finite family of
matched compact truncations, including the base, has comparison
images with disjoint interiors for all sufficiently large \(i\).

The topology premise for the base holds, for example, if a whole
compact base truncation embeds with boundary a union of smooth
\(S^3\)'s and its compactification fills those boundaries by
smooth balls. Mayer--Vietoris in degree zero gives one
complementary component for each boundary sphere. The degree-two
sequence gives \(H_2=0\) for the core, and van Kampen across the
simply connected boundary spheres makes its fundamental group a
free factor of the trivial ambient group. Capping preserves these
groups. Poincar\'e duality and the simply connected homology
criterion then give a smooth homotopy four-sphere. This calculation
specifies the capped topology and permits the complementary pieces
to have their given smooth structures.

The compactifications are Bach-flat and have positive Yamabe
constant. In fact, for a test of fixed compact support in a
comparison domain, take the Yamabe quotient in the conformally
changed pullback. Conformal covariance transfers it to a test on
the source. Its quotient is at least \(Y_i\), hence every limit
has compact-support Yamabe constant at least \(y>0\).
The cutoffs at the finite set of smooth points change the squared
Dirichlet norms by a vanishing amount, since the squared cutoff
gradient integral is \(O(\varepsilon^2)\) per point. Therefore the
punctured infimum is the closed Yamabe constant. For the Bach
equation, the conformally invariant \(L^1\) Bach density has
vanishing integral in the limit on each fixed compact set.
Smooth convergence and conformal covariance give Bach-flatness
on the convergence domains, and smoothness gives it across the
added points. Constant-scale comparison of the component metrics
also transfers the trace-free Ricci integral bound needed for
\eqref{msgap:strict}, even if finitely many smooth points were
omitted from the convergence.

Let \(e_0\) be the common chiral energy on the base compactification
and \(e_j\) the common energy on each component. Equality of the
chiralities follows from the closed signature formula on the
compactifications and conformal invariance. A flat component
contributes zero. Nonnegativity, disjoint interiors of the source
images and smooth convergence on each fixed family of truncations
give, in either chirality,
\begin{equation}\label{msgap:accounting}
 e_0+\sum_j e_j\le e .
\end{equation}
Embedded boundaries have zero source Weyl measure. First let
\(i\to\infty\) at fixed truncations, then exhaust the comparison
domains. Finite sets of already smooth deleted points have zero
limiting Weyl measure. For a countable family the sum is the
supremum over finite subfamilies. In a finite family the difference
in \eqref{msgap:accounting} is the iterated limit of the corresponding
source Weyl energy outside the union of the images, with the source
index tending to infinity before removing the truncation.
A full energy identity is the additional assertion of equality in
\eqref{msgap:accounting}.

If there is a nonflat component, it is unique among the components
in this accounting, since each nonflat component costs strictly more
than \(4\pi^2\). For that component \(X\) the residual satisfies
\(e-e_X<4\pi^2-\delta\), so the smooth positive-Yamabe Bach-flat base
has \(e_0\le e-e_X<4\pi^2\). Theorem~\ref{msgap:closed} makes
the base conformally round and \(e_0=0\). All other components
have zero energy. Thus the unaccounted energy equals \(e-e_X\)
and is strictly positive by \eqref{msgap:strict}, with the
uniform bound \eqref{msgap:uniform-residual}. This contradicts a
full energy identity. If there are no nonflat components, an energy
identity instead puts all of \(e\) on the smooth compactified base,
so \(e=e_0\in\mathcal A\). The argument concerns the supplied
embeddings, smooth extensions and accounting. Extracting comparison
maps from the almost-Bach hypothesis and locating the residual mass
would require further estimates.

\section{Compact conformal gauges and shrinking round caps}
\label{msgauge:section}

The compactness question has a different conclusion when the sequence
has a common metric gauge on the entire manifold. In that gauge the
metric difference is an admissible first-variation test. It can recover
strong convergence of the Weyl tensor from weak convergence of the
second metric derivatives. The round marking then needs to retain
positive volume to force smooth recognition.

\subsection{Strong Weyl convergence in a weak metric gauge}

\begin{theorem}[Strong Weyl convergence in a supplied compact gauge]
\label{msgauge:compact}
Let \(M\) be a closed, connected, oriented smooth homotopy four-sphere
and let \(\bar g_i\) be smooth Riemannian metrics on \(M\). Suppose
smooth conformal representatives \(g_i=u_i^2\phi_i^*\bar g_i\), with
\(u_i>0\) and \(\phi_i\) diffeomorphisms of \(M\), satisfy
\begin{equation}\label{msgauge:hypotheses}
 c r\le g_i\le C r,\qquad
 \|g_i\|_{W^{2,q}(r)}\le C_q,\qquad
 \int_M |B_{g_i}|_{g_i}\,dV_{g_i}\longrightarrow0
\end{equation}
for some smooth reference \(r\), constants \(0<c\le C\) and finite
\(q>2\). A subsequence converges weakly in \(W^{2,q}\) and strongly in
\(C^0\cap W^{1,4}\) to a positive metric \(g\). The full Weyl tensors
converge strongly in \(L^2(r)\). The limit is critical for the full
Weyl functional and has a smooth Bach-flat representative \(\gamma\)
on \(M\) after a \(W^{2,q}\) conformal rescaling and a \(W^{3,q}\)
diffeomorphism. The energies \(E_\pm(g_i)\) and Yamabe constants
\(Y([g_i])\) converge to the corresponding invariants of \(\gamma\).

If along the subsequence there are measurable sets \(K_i\) with
\(W_{g_i}=0\) almost everywhere on \(K_i\) and
\(\Vol_{g_i}(K_i)\ge v>0\), then \(\gamma\) is conformally
diffeomorphic to the round sphere and \(E_\pm(g_i)\to0\).
\end{theorem}

The last condition in \eqref{msgauge:hypotheses} is gauge invariant.
Indeed the covariant Bach tensor has weight \(-2\) under \(g\mapsto
u^2g\). Its pointwise norm has weight \(-4\), canceled by the volume.
Oriented diffeomorphisms preserve both chiral energies. Orientation
reversal exchanges them, which makes no difference on \(M\) by
\eqref{eq:sign}.

\begin{proof}
All reference tensor pairings and integrals use \(r\) and
\(d\mu=dV_r\). Reflexivity, Morrey's embedding and Rellich compactness
give, after extraction,
\begin{equation}\label{msgauge:weak-metrics}
 g_i\rightharpoonup g\text{ in }W^{2,q},\qquad
 g_i\longrightarrow g\text{ in }C^0\cap W^{1,2q}.
\end{equation}
In dimension four the strict inequality \(q>2\) gives the compact
embedding of first derivatives from \(W^{1,q}\) to \(L^{2q}\).
The continuous limit retains the positive metric bounds.

In a fixed chart, or with a fixed reference derivative \(D\), write
the full covariant Weyl tensor as
\begin{equation}\label{msgauge:curvature-map}
 W_g=A(g)D^2g+J(g,Dg).
\end{equation}
The principal coefficient is algebraic in \(g,g^{-1}\). The lower
terms are polynomial of degree at most two in \(Dg\), with smooth
coefficients in the metric and the reference data. This follows
from the connection difference formula for curvature and its
algebraic Weyl projection. Scalar curvature has an expression of
the same type. In \eqref{msgauge:weak-metrics} the principal terms
converge weakly in \(L^q\) and the lower terms strongly in \(L^q\).
In particular \(W_i=W_{g_i}\rightharpoonup W_g\) and
\(R_{g_i}\rightharpoonup R_g\) in \(L^q\).
The same products make \(g\mapsto W_g\) a \(C^1\) map from the open
cone of positive \(W^{2,q}\) metrics to \(L^q\).

Set \(h_i=g_i-g\). In the derivative \(DW_{g_i}[h_i]\), the
second-derivative contribution \(A(g_i)D^2h_i\) cancels the
corresponding part of \(W_i-W_g\). More precisely,
\[
 \begin{split}
 D_g(A(g_i)D^2g_i)[h_i]
     -\big(A(g_i)D^2g_i-A(g)D^2g\big)
  =A'(g_i)[h_i]D^2g_i-(A(g_i)-A(g))D^2g .
 \end{split}
\]
Both right-hand terms tend to zero in \(L^2\) by uniform convergence
of the metric coefficients and bounded second derivatives in \(L^2\).
The difference of the lower terms tends to zero in \(L^2\)
as well. Their derivative is bounded there by a constant times
\[
 \|h_i\|_\infty(1+\|Dg_i\|_4^2)
    +\|Dh_i\|_4(1+\|Dg_i\|_4),
\]
which vanishes. A finite atlas therefore gives
\begin{equation}\label{msgauge:metric-test}
 DW_{g_i}[h_i]-(W_i-W_g)\longrightarrow0\quad\text{in }L^2(r).
\end{equation}

Let \(H(g)\) denote the positive symmetric reference fiber map
incorporating the four inverse metrics and the density ratio
in \(\mathcal W(g)=\int_M |W_g|_{g,\mathrm{tens}}^2\,dV_g\).
The \(C^1\) derivative is
\begin{equation}\label{msgauge:first-variation}
 D\mathcal W_g[h]
  =2\int_M\langle H(g)W_g,DW_g[h]\rangle_r\,d\mu
    +\int_M\langle H'(g)[h]W_g,W_g\rangle_r\,d\mu .
\end{equation}
For smooth \(g_i\) the Bach identity extends from smooth \(h\) to
\(W^{2,q}\) directions by density, since the approximation converges
uniformly as well. Thus \(h_i\) is an admissible unrestricted test, and
\[
 |D\mathcal W_{g_i}[h_i]|
     \le C'\|h_i\|_\infty\int_M|B_{g_i}|_{g_i}\,dV_{g_i}\longrightarrow0 .
\]
The second integral of \eqref{msgauge:first-variation} is
\(O(\|h_i\|_\infty\|W_i\|_2^2)=o(1)\). Equations
\eqref{msgauge:metric-test} and \eqref{msgauge:first-variation} give
\(\int_M\langle H(g_i)W_i,W_i-W_g\rangle_r\,d\mu\to0\).
Also \(H(g_i)\to H(g)\) uniformly, so the weak convergence of
\(W_i-W_g\) gives
\(\int_M\langle H(g_i)W_g,W_i-W_g\rangle_r\,d\mu\to0\).
Subtracting and using uniform positive definiteness gives
\begin{equation}\label{msgauge:strong-weyl}
 W_i\longrightarrow W_g\quad\text{in }L^2(r).
\end{equation}
The chiral projections and contractions only involve zeroth-order
metric coefficients, so \eqref{msgauge:strong-weyl} also passes
each chiral energy to the limit.

For a fixed smooth symmetric test \(h\), \(DW_{g_i}[h]\)
converges weakly in \(L^2(r)\) to \(DW_g[h]\). The only term
containing \(D^2g_i\) has uniformly convergent coefficient
\(A'(g_i)[h]\), and the differentiated lower terms converge
strongly. Now \eqref{msgauge:strong-weyl} passes the pairing in
\eqref{msgauge:first-variation} to the limit. Its value is zero
by the Bach bound in \eqref{msgauge:hypotheses}. Smooth density
then gives criticality on the full cone of \(W^{2,q}\) metrics.

\paragraph{Yamabe constants.}
For the weak metric use the same quotient as in
\eqref{eq:energies}, with quadratic form
\(Q_g(f)=\int_M(6|df|_g^2+R_gf^2)\,dV_g\) and domain \(H^1(M,r)\).
Smooth tests are dense for this quotient by the Sobolev embedding and
\(R_g\in L^q\). Testing by fixed smooth functions and using weak scalar
convergence gives \(Y(g)\ge\limsup_i Y([g_i])\).
For the converse take approximate minimizers \(f_i\) on a subsequence
realizing the lower limit, normalized by \(\|f_i\|_{4,g_i}=1\).
Their kinetic energies are bounded because the quotient is bounded
above by the constant test and the scalar terms are bounded in
absolute value by Hölder. Thus \(f_i\) is bounded in \(H^1(r)\).
Writing \(q'=q/(q-1)\), Rellich gives a subsequence converging strongly
in \(L^{2q'}\), since \(2q'<4\). The squares \(f_i^2\) can therefore
test the weak convergence of the scalar densities in \(L^q(r)\).
The kinetic coefficients and volume ratios converge uniformly.
It follows that
\(Q_{g_i}(f_i)-Q_g(f_i)\to0\) and \(\|f_i\|_{4,g}^2\to1\).
The quotient at \(g\) gives \(\liminf_i Y([g_i])\ge Y(g)\).

\paragraph{A smooth representative.}
The global regularity theorem in the October 2026 preprint of
Avalos, Cogo and Royo Abrego applies to the critical \(W^{2,q}\) metric on this smooth compact
four-manifold. It gives \(u>0\) in \(W^{2,q}\) and a global \(C^1\)
diffeomorphism \(\Phi\), with \(\Phi,\Phi^{-1}\in W^{3,q}\), such that
\(\gamma=\Phi^*(u^2g)\) is smooth on \(M\)%
~\cite[Theorem 1.1, \S2.3 and \S5.1]{AvalosCogoRoyoAbrego2026Bach}.
In particular \(\Phi\) is bi-Lipschitz, and \(u\) has positive
minimum by continuity. We verify the covariance needed to use
this smooth representative.

Multiplication by \(u^2\) and pullback by \(\Phi\) give isomorphisms
on \(W^{2,q}\) metric directions%
~\cite[Lemmas 2.1 and 2.4]{AvalosCogoRoyoAbrego2026Bach}.
One can pass the smooth covariance rules by approximation. Approximate
\(g,u\) strongly in \(W^{2,q}\), with positivity retained by uniform
convergence. Approximate \(\Phi\) in an embedding of \(M\) and apply a
tubular neighborhood retraction. Strong \(W^{3,q}\) approximation
implies \(C^1\) closeness. The approximations are local diffeomorphisms
of degree \(\pm1\) and hence global diffeomorphisms. Their Jacobians
and inverse Jacobians are uniformly bounded. The pullbacks are
uniformly bounded on \(W^{2,q}\): their chain rules use the \(C^0\)
bound on the first derivatives and \(L^{2q}\), \(L^q\) bounds on the
second and third derivatives of the maps. This gives strong
convergence of the pullbacks first on smooth tensors and then
by density. The \(L^q\) pullbacks of curvature tensors use only the
Jacobian and first derivative bounds. By strong continuity in
\eqref{msgauge:curvature-map} this proves scalar and Weyl naturality
and the conformal rule \(W_{u^2g}=u^2W_g\) for the covariant tensor
at this regularity. Consequently the change to \(\gamma\) preserves
\(\mathcal W\). Since it is also an isomorphism of metric directions,
it preserves criticality. Thus \(B_\gamma=0\). The chiral energies
are unchanged up to exchange of orientation.

Likewise the conformal scalar rule gives
\(Q_{u^2g}(f)=Q_g(uf)\) and \(\|f\|_{4,u^2g}=\|uf\|_{4,g}\)%
~\cite[\S2, (2.6)--(2.7)]{LeeParker1987Yamabe}.
The identities follow by strong approximation for smooth tests,
then by \(H^1\) density. The multipliers \(u,u^{-1}\) act boundedly
on \(H^1\) since their gradients lie in \(L^4\) and the factors
are bounded. Composition by the \(C^1\) diffeomorphism is an
\(H^1\) isomorphism and preserves the quotient by change of variables.
Therefore \(Y([\gamma])=Y(g)\).

\paragraph{A Weyl-flat region of positive measure.}
If the sets in the last assertion have \(\Vol_{g_i}(K_i)\ge v\), then
\(\mu(K_i)\ge v_0>0\). For any \(a>0\),
\[
 \mu\big(K_i\cap\{|W_g|_r>a\}\big)
   \le a^{-2}\|W_i-W_g\|_{2,r}^2\longrightarrow0 .
\]
It follows that \(\mu(\{|W_g|_r\le a\})\ge v_0\). Letting \(a\)
decrease to zero shows that the zero set has positive measure.
The covariance just proved and the bi-Lipschitz change of variables
give a Weyl zero set of positive measure for \(\gamma\).
In one connected conformal harmonic chart that meets this set
in positive measure, the smooth Bach-flat metric has a real-analytic
determinant-normalized conformal representative%
~\cite[Theorems 2.2 and 3.7(c)]{LassasLiimatainen2019ConformalHarmonic}.
Each Weyl component is analytic. A nontrivial real-analytic function
on a connected chart has a null zero set: each zero lies in a
regular zero hypersurface of one of its derivatives, obtained
from a minimal nonvanishing derivative. Thus all Weyl components
vanish on the chart. Lemma~\ref{var:continuation} propagates the
vanishing of \(W^+\), and Proposition~\ref{prop:recognition} gives
the conformally round conclusion.
\end{proof}

For sequences in the band
\[
 Y([\bar g_i])>0,\qquad
 4\pi^2\le E_+(\bar g_i)\le 8\pi^2-\delta_0,\qquad \delta_0>0,
\]
the condition \(\mathfrak B(\bar g_i)/E_+(\bar g_i)\to0\), where
\(\mathfrak B(g)=\int_M|B_g|_g\,dV_g\), implies the Bach condition
in \eqref{msgauge:hypotheses}. If the common gauges exist, the
Weyl-flat sets of any such sequence must have gauged volume tending
to zero, because a subsequence of positive lower volume would
have energy tending to zero. In particular, a fixed marked round
cap before applying the gauges can still shrink in these gauges.
Whether the caps shrink or are absent, the smooth Bach-flat
representative has the limiting energy and
\(Y([\gamma])\ge\sqrt{48\delta_0}\) by \eqref{var:wall}.
The closed gap of Theorem~\ref{msgap:closed} excludes the lower
energy endpoint for this representative. The common global
gauge and its metric bounds are additional hypotheses.

\subsection{A shrinking round cap from a closed Bach-flat class}

\begin{proposition}[Conditional construction inside the band]
\label{msgauge:reverse}
Suppose \(M\) admits a smooth Bach-flat metric \(g\) with
\[
 Y([g])>0,\qquad 4\pi^2<E=E_+(g)<8\pi^2-\delta_0,\qquad \delta_0>0.
\]
At any prescribed \(p\in M\) there are smooth Riemannian metrics
\(r_\varepsilon\) with a unit-round cap of coordinate radius
comparable to \(\varepsilon\), and
\begin{equation}\label{msgauge:splice-rates}
 E_+(r_\varepsilon)=E+O(\varepsilon^4),\qquad
 \mathfrak B(r_\varepsilon)=O(\varepsilon^2),\qquad
 Y([r_\varepsilon])=Y([g])+O(\varepsilon^2).
\end{equation}
These metrics have uniformly bounded second derivatives in a
fixed smooth atlas and converge to \(g\) in \(C^1\).
They can be chosen to be regular centers for the full fixed-volume
prescription. The mark can also be installed on a prescribed
fixed ball at unit round scale by changing representatives.
In unit-volume Yamabe representatives the conformally round
cap has diameter of order \(\varepsilon\) and volume of order
\(\varepsilon^4\).
\end{proposition}

\begin{proof}
In oriented \(g\)-normal coordinates at \(p\) let
\(s=(1+|x|^2/4)^{-2}\delta_{\rm eucl}\), the unit sectional curvature
metric in rescaled stereographic coordinates. The zero and first
metric jets of \(s\) and \(g\) agree at \(0\). For a smooth radial
\(0\le\chi\le1\) equal to zero on \(B_1\) and one outside \(B_2\),
set \(r_\varepsilon=s+\chi(x/\varepsilon)(g-s)\) in the chart
and extend by \(g\). This convex combination is uniformly positive
definite. The product rule gives
\begin{equation}\label{msgauge:splice-jets}
 |D^j(r_\varepsilon-g)|\le C\varepsilon^{2-j},\qquad
        0\le j\le4
\end{equation}
on the splice. For zero or one derivative on \(g-s\), use the
matched first jets. Higher derivatives of \(g-s\) are bounded in
the fixed chart. The bound persists in fixed smooth coordinates by
the chain rule. A closed ball of radius \(\varepsilon/2\) has
a round neighborhood. Curvature is bounded on the splice, and its
first and second covariant derivatives are \(O(\varepsilon^{-1})\)
and \(O(\varepsilon^{-2})\). The Bach tensor has a term with two
Schouten derivatives and a curvature product, hence is
\(O(\varepsilon^{-2})\). Outside the splice and in the core it
vanishes%
~\cite[\S1, (1.1)]{BahuaudGuentherIsenbergMazzeo2025Bach}.
The affected volume is \(O(\varepsilon^4)\). Integration gives the
first two rates in \eqref{msgauge:splice-rates}.

For the Yamabe rate use the fixed \(g\) form. Since \(Y([g])>0\),
\(Q_g(f)\ge Y([g])\|f\|_{4,g}^2\) controls \(\|f\|_{2,g}^2\),
and the gradient term then controls \(\|f\|_{H^1(g)}^2\).
The kinetic coefficients in the perturbed form differ by
\(O(\varepsilon^2)\) in \(C^0\). The scalar densities differ
by \(O(\varepsilon^2)\) in \(L^2(g)\), since both scalar curvatures
are bounded and the difference is supported on volume
\(O(\varepsilon^4)\). Hölder and Sobolev give
\[
 |Q_{r_\varepsilon}(f)-Q_g(f)|
       \le C\varepsilon^2\|f\|_{H^1(g)}^2
       \le C'\varepsilon^2 Q_g(f).
\]
Compare the \(L^4\) and \(L^2\) denominators by uniform volume
convergence. Taking infima proves the Yamabe rate and a uniform
positive lower bound on the first conformal-Laplacian eigenvalue.
The strict energy margins keep these classes in the original band.

The fixed marking uses the coordinate germ and compression from
Proposition~\ref{var:marking}. Take the compressed region containing
the support of the conformal cutoff wholly inside the round core.
On that region in marked coordinates the pulled-back metric is
\(c_\varepsilon^2(1+c_\varepsilon^2|x|^2/4)^{-2}\delta_{\rm eucl}\),
where \(c_\varepsilon\) is a sufficiently small fixed multiple of
\(\varepsilon\). The local scale to the prescribed unit round
metric is \(2(1+c_\varepsilon^2|x|^2/4)/
(c_\varepsilon(1+|x|^2))\). Extending its logarithm with the
cutoff produces a bounded total volume: the factor is
\(O(c_\varepsilon^{-1})\) where the pulled-back volume density
is \(O(c_\varepsilon^4)\). On the complement only the pullback
acts. This representative preserves the three invariants in
\eqref{msgauge:splice-rates} and has first conformal-Laplacian
eigenvalue bounded below by \(Y/\Vol^{1/2}>0\), using its own
Yamabe constant and volume. The literal unit round metric on
the fixed ball has fixed positive physical volume. Pullback
alone can fix the label of the shrinking ball while retaining its
small physical scale.

To justify that scale in the Yamabe gauge, write
\(g_\varepsilon^Y=u_\varepsilon^2r_\varepsilon\) for unit-volume
minimizers. Their factors form a minimizing sequence for the
fixed quotient at \(g\) by the form comparison. Since \(E>0\),
the strict Yamabe inequality gives \(0<Y([g])<Y_S\), where \(Y_S\)
is the round value%
~\cite[Theorems 2--3]{Schoen1984Yamabe}.
Take a weak \(H^1(g)\) limit \(u\) along a subsequence and let
\(a=\|u\|_{4,g}^4\). The squared gradient energy splits into the
limit term and the term for \(u_\varepsilon-u\) up to \(o(1)\).
The \(L^4\) masses also split: mixed terms vanish by approximating
\(u\) by bounded functions and using strong convergence below
exponent four. Aubin's almost-sharp Sobolev inequality on the
fixed background, with the coefficient loss tending to zero
and the \(L^2\) remainder vanishing on the difference, gives
\[
 Y([g])\ge Y([g])\sqrt a+Y_S\sqrt{1-a}.
\]
Hence \(a=1\), and the convergence is strong in \(H^1\) and
\(L^4\)%
~\cite[\S2, Theorems 2.2--2.3 and \S3]{LeeParker1987Yamabe}.

For uniform pointwise bounds, test the minimizing equation
\(6\Delta_{r_\varepsilon}u_\varepsilon+
R_{r_\varepsilon}u_\varepsilon=Y([r_\varepsilon])u_\varepsilon^3\)
against \(\eta^2u_\varepsilon^3\), using a smooth local cutoff
\(\eta\). Integration by parts with the uniform Sobolev
inequality gives
\[
 \|\eta u_\varepsilon^2\|_4^2
 \le C\int_M (1+|d\eta|^2)u_\varepsilon^4\,dV_{r_\varepsilon}
     + C\|u_\varepsilon\|_{L^4(\operatorname{supp}\eta)}^2
             \|\eta u_\varepsilon^2\|_4^2 .
\]
The strong \(L^4\) convergence allows a fixed finite cover on
which the last term is absorbed. It gives bounded \(L^8\) norms.
The right side of the Laplacian equation is then bounded in
\(L^{8/3}\). Apply the global elliptic estimate in the fixed
metric \(g\); the operator difference from
\(\Delta_{r_\varepsilon}\) is \(o(1)\) from \(W^{2,8/3}\) to
\(L^{8/3}\) by \(C^1\) convergence, and is absorbed. The resulting
\(W^{2,8/3}\) bound gives a uniformly convergent subsequence.
The limit is a nonzero nonnegative minimizer, smooth by
elliptic bootstrapping and strictly positive by the minimum
principle. Applying this reasoning to any subsequence gives
uniform upper and positive lower bounds on \(u_\varepsilon\)%
~\cite[\S2, Theorems 2.2 and 2.5--2.6]{LeeParker1987Yamabe}.
Thus the raw and unit-volume Yamabe representatives are uniformly
comparable, proving the stated physical cap scales.

Finally apply Section~\ref{msreg:section} separately to each
\(r_\varepsilon\), retaining the round neighborhood of the
smaller closed ball. The exterior perturbations can preserve
density and energy exactly and be arbitrarily small in \(C^4\)
at each radius. Choose them of \(C^4\) size at most
\(\varepsilon^4\) and with Bach integral error at most
\(\varepsilon^2\), by continuity at each fixed metric. The scalar
form estimate and the preceding compactness of Yamabe factors
persist, using the same bounded second jets and the additional
uniform scalar error \(O(\varepsilon^4)\). Surjectivity passes to
smooth conformal rescalings and oriented pullbacks. In fact,
under \(r'=u^2r\), the identification of the trace-free increments
\(h'=u^2h\) leaves the logarithmic endomorphisms unchanged and
\(L_{r'}(u^2h)=u^{-2}L_rh\) in the common Weyl endomorphism bundle.
The source and target identifications are invertible at each
fixed reference.
\end{proof}

At the regular centers every differentiable radial tangent has
intrinsic supremum norm at least \(c\varepsilon^{-2}\) by
\eqref{mscost:strain} and \eqref{msgauge:splice-rates}. The
original positive-energy Bach-flat class is stationary, so it
cannot have a differentiable radial tangent: such a tangent
would give derivative \(-2E\). These are conditional
examples of small total Bach norm and large radial strain
inside the band. The construction assumes the closed class
exists. In the raw gauge the full curvature is uniformly bounded,
so any smooth pointed rescaling at scales tending to zero is flat
in that gauge.

\paragraph{Energy at the endpoints.}
The closed Bach-flat gap excludes a positive-Yamabe input with
\(E=4\pi^2\) by Theorem~\ref{msgap:closed}. At the upper endpoint
\(E=8\pi^2-\delta_0\), the splice at a point where \(W_g^+(p)\ne0\)
strictly increases energy for all sufficiently small radii.
To see this, set \(d_\varepsilon=r_\varepsilon-g\) and let
\(\mathcal F(d)\) be the positive symmetric transport of \(W^+_{g+d}\)
to the bundle at \(g\), multiplied by the square root of the
volume-density ratio. Then
\(E_+(g+d)=\|\mathcal F(d)\|_{2,g}^2\). Put
\(\mathcal D=D\mathcal F(0)\). By the affine second-jet dependence
of curvature and \eqref{msgauge:splice-jets},
\[
 \mathcal F(d_\varepsilon)=W_g^+ +\mathcal D d_\varepsilon
                       +O(\varepsilon^2)
\]
pointwise with error supported on the splice. The coefficient
variation multiplying \(D^2d_\varepsilon\) is \(O(\varepsilon^2)\)
because it depends only on the metric values. The lower-jet Taylor
errors have the same bound. The integral linear term vanishes
by Bach stationarity, so
\[
 E_+(r_\varepsilon)-E
     =\|\mathcal D d_\varepsilon\|_{2,g}^2+O(\varepsilon^6).
\]
In the round core \(\mathcal D d_\varepsilon=-W_g^+
+O(\varepsilon^2)\), which gives a lower bound
\(c\varepsilon^4\) for the squared integral at this point.
The equal-energy exterior perturbation preserves this overshoot.
Convergence of the invariants still gives a fixed smaller positive
margin, for example \(\delta_0/2\). At \(E=8\pi^2\), convergence
alone gives no positive margin.

\section{Outstanding challenges}\label{sec:outlook}

\paragraph{Continuing a radial contraction.}
The local prescription theorem gives a regular metric at every marked
reference energy. Each new center has an interval on which its full
curvature chord lifts. The interval depends on the center, and the
lift preserves \(W^+=0\) on the marked ball while its metric there can
deform. A global construction must control the successive step sizes,
the metric and its compatible connection, and the cost of restoring
a round mark whenever the next step uses one. For a chord with initial energy \(E_0>0\), the subcritical and
recognition energy inequalities respectively require
\(t>1-\sqrt{8\pi^2/E_0}\) and \(t>1-\sqrt{4\pi^2/E_0}\) in its supplied
interval \(0\le t\le1\). Strict inequality excludes an equality slice.
A contraction through a specified singular reference also requires information about the sign
of its adjoint fields as in Section~\ref{mscone:section}. At an ordinary
differentiable start the necessary condition is \(W_r^+\perp\ker L_r^*\).
The construction smooth in \(\sqrt t\) addresses a wider set of starts.

\paragraph{The Bach gradient at a fixed energy margin.}
Can a sequence of round-marked positive-Yamabe classes on a given smooth
homotopy sphere satisfy
\[
 4\pi^2\le E_+\le8\pi^2-\delta,\qquad
 \frac{\int |B|\,dV}{E_+}\longrightarrow0?
\]
A single positive-Yamabe closed Bach-flat class strictly inside this
band on that manifold would give such a sequence by the shrinking-mark
construction. Existence of a positive closed critical value below
\(8\pi^2\) remains open here; Theorem~\ref{msgap:minimum}
describes the least value if one exists. A positive lower bound on
the displayed ratio would remove the strain obstruction in
\eqref{mscost:strain}. A bound from above on the metric cost of
solving the radial equation would still require an inverse estimate
on that particular direction.

\paragraph{Energy between concentration scales.}
A supplied nonflat smooth ALE component with full embedded core leaves
a uniformly positive Weyl energy outside that component. The closed
critical gap prevents a smooth Bach-flat homotopy-sphere base from
carrying the remainder if it is the only other term of an energy identity.
For an approximately critical sequence the estimates leave open both
the extraction of a smooth component and the location of any residual
energy. A concrete test is a long annulus on which the total Bach norm
tends to zero. In a common conformal coordinate gauge, can one test
the Weyl first variation against the cut-off difference between the
metric and its Euclidean model and control the two boundary terms
independently of the annulus length? A geometric application also
requires estimates that furnish the common gauge across scales.
Orbifold intermediate regions require their own boundary and energy
analysis.

\paragraph{Recognizing the original manifold.}
The metric constructed at the end of a continuation, or recovered in
a compactness argument, must determine the smooth type of the starting
manifold. A smoothly capped blow-up core can itself be a homotopy sphere
of unknown smooth type. The compact-gauge theorem in
Section~\ref{msgauge:section} retains the original manifold, provided
its uniform ellipticity and \(W^{2,q}\) bounds are available. A surviving
marked volume then forces conformal roundness. Along a deformation the
alternative is a positive Yamabe checkpoint and continuation within
\(E_+<8\pi^2\) to \(E_+<4\pi^2\), where Proposition~\ref{prop:recognition}
applies. Producing the required starting geometry and either form of
control on an arbitrary smooth homotopy four-sphere remains the
existence problem behind both approaches.

\clearpage
\setlength{\bibitemsep}{.36\baselineskip}
\printbibliography[heading=bibintoc]
\end{document}
